\chapter{Energy systems integration}
\label{vol04:ch14}

\epigraph{If everyone does a little, we'll achieve only a little. We
must do a lot.}{David J.\ C.\ MacKay, \emph{Sustainable Energy without
the Hot Air} (UIT Cambridge, 2009), p.~3
\cite[p.~3]{gen:mackay2009}.}

\chapmeta{Half-life: medium. Archetypes: balance, transport, scaling, control. Exercise target: 34. Project track: simulation plus household-data analysis (hybrid). Hazard class: low (analytical and modelling work; hands-on rooftop solar work is hazard class medium-to-high and outside the chapter's scope). Reader path default: standard, with sections explicitly tagged. Named cases: Three Mile Island (1979) as a reactor outage exemplar (cited); large grid blackouts (Northeast 2003, Texas 2021) described by mechanism without invoking the named case (Texas 2021 has a registry entry, Northeast 2003 does not; battery-storage fire cases McMicken 2019 and Moss Landing 2021/2024 have no registry entry).}

\noindent
The energy system is the integration of generation, transmission,
distribution, storage, and use across electrical, thermal, and
chemical energy carriers. Volume IV's preceding chapters supply the
component-scale physics; this chapter takes the components and runs
them together as a system, with the operational constraints of grid
stability, demand-side management, intermittency of renewable
sources, and the engineering of the energy transition that is the
defining infrastructure question of the early twenty-first century.
The half-life of the contents is medium rather than durable: the
energy mix, the storage technologies, the grid topology, and the
regulatory framework will all look different by 2046, and the
chapter's specific numbers carry their year accordingly.

\section{The grid: generation, transmission, distribution}\pathtag{core}

The \engterm{electric power grid} is a synchronous AC network at
$50\,\si{\hertz}$ (most of the world) or $60\,\si{\hertz}$ (North
America, parts of Japan, parts of South America) that connects
generation to load over distances from metres to thousands of
kilometres. Five working layers are conventionally distinguished:

\begin{itemize}[leftmargin=*]
  \item \engterm{Generation}: the production of electrical energy from
    a primary source (chemical, nuclear, hydraulic, wind, solar,
    geothermal). Generation is conventionally measured in
    $\si{\mega\watt}$ or $\si{\giga\watt}$, with annual energy
    production in $\si{\tera\watt\hour}$.
  \item \engterm{Transmission}: long-distance high-voltage transfer of
    bulk electrical power, typically at $115$ to $765\,\si{\kilo\volt}$
    AC or up to $\pm 800\,\si{\kilo\volt}$ HVDC. The voltage levels
    minimise $I^{2}R$ losses; a $400\,\si{\kilo\volt}$ line carries
    $50\,\si{\times}$ the power of a $20\,\si{\kilo\volt}$ line at the
    same conductor loss.
  \item \engterm{Sub-transmission}: regional voltage levels of $34$ to
    $115\,\si{\kilo\volt}$ feeding metropolitan and industrial loads.
  \item \engterm{Distribution}: medium-voltage feeders at $11$ to
    $35\,\si{\kilo\volt}$ that fan out from substations to the
    secondary distribution transformers serving neighbourhoods.
  \item \engterm{Service}: low-voltage delivery to end consumers at
    $120/240\,\si{\volt}$ (North America), $230/400\,\si{\volt}$
    (Europe and most of the world).
\end{itemize}

A balanced three-phase generator-line-load chain delivers power
$P = \sqrt{3}V_{L}I_{L}\cos\phi$ as developed in Chapter 9. The
chain's losses are dominated by transmission-line $I^{2}R$ losses and
by transformer iron and copper losses; a modern bulk-power
transmission system has end-to-end losses of $\sim 5\,\%$ in the
United States and Europe (current as of 2025), distribution losses
add another $\sim 3\,\%$, and end-use efficiency varies widely with
the load.

\engterm{Frequency regulation} is the central operational discipline
of synchronous AC grids. All generators and motors on the grid are
mechanically synchronised through the magnetic coupling of their
rotors to the grid's rotating-field reference. A momentary imbalance
between generation and load (a generator trip, a large industrial
load switching on) causes a small change in grid frequency, which
acts as a global error signal to which all generators respond. Three
control time scales operate:
\begin{itemize}[leftmargin=*]
  \item \engterm{Primary frequency response}: governor-based, droop
    control, seconds. Each generator's mechanical input is reduced
    in proportion to its over-frequency or increased in proportion to
    its under-frequency, with a droop of about $5\,\%$.
  \item \engterm{Secondary frequency response} or
    \engterm{automatic generation control}: control-room or
    automated dispatch, minutes. Designated regulating generators
    adjust their setpoints to restore the frequency to nominal.
  \item \engterm{Tertiary response}: market-based dispatch, hours.
    Generators are ramped in or out of merit order according to
    forecast loads and generation availability.
\end{itemize}

A working frequency band of $\pm 0.05\,\si{\hertz}$ for steady-state
operation and $\pm 0.5\,\si{\hertz}$ for emergency operation is
standard in interconnected European and North American grids. Outside
the emergency band, generators trip on their own protective relays
to avoid catastrophic equipment failure; the consequence is a
cascading separation of the grid into electrically isolated
``islands.''

% Figure: grid frequency response to a sudden generator trip. Frequency
% (Hz) vs time (s) showing the nadir, primary response arrest, and
% secondary recovery to nominal. Nominal 50 Hz.
\begin{figure}[htbp]
\centering
\begin{tikzpicture}
\begin{axis}[
  width=13cm, height=8cm,
  xlabel={time after trip (s)},
  ylabel={grid frequency (Hz)},
  xmin=0, xmax=60, ymin=49.5, ymax=50.1,
  xtick={0,10,20,30,40,50,60},
  ytick={49.5,49.6,49.7,49.8,49.9,50.0,50.1},
  axis lines=left,
  tick label style={font=\scriptsize},
  label style={font=\small},
]
% nominal line
\draw[enggray, dashed] (axis cs:0,50) -- (axis cs:60,50);
\node[enggray, font=\scriptsize, anchor=south east] at (axis cs:60,50) {nominal 50 Hz};
% the trip happens at t=2; frequency falls to nadir, arrested, recovers
\addplot[engblue, very thick, smooth] coordinates {
  (0,50.0)(2,50.0)(3,49.85)(4,49.72)(5,49.66)(6,49.65)(8,49.68)(12,49.74)(20,49.82)(30,49.90)(45,49.97)(60,50.0)
};
% nadir marker
\filldraw[engred] (axis cs:6,49.65) circle [radius=1.6pt];
\node[engred, font=\scriptsize, anchor=north] at (axis cs:6.5,49.64) {nadir};
% phase annotations
\draw[engamber, |-|] (axis cs:2,49.55) -- (axis cs:8,49.55);
\node[engamber, font=\scriptsize, anchor=north] at (axis cs:5,49.55) {primary (inertia + governor)};
\draw[englime!60!black, |-|] (axis cs:8,49.55) -- (axis cs:60,49.55);
\node[englime!60!black, font=\scriptsize, anchor=north] at (axis cs:34,49.55) {secondary (AGC)};
% trip arrow
\draw[engred, -{Stealth}] (axis cs:2,49.95) -- (axis cs:2,50.0);
\node[engred, font=\scriptsize, anchor=north] at (axis cs:2,49.93) {trip};
\end{axis}
\end{tikzpicture}
\caption[Grid frequency response to a generator trip]{Grid frequency
after a sudden generator trip. The rotor inertia of the synchronous fleet
limits the initial rate of fall; governor (primary) response arrests the
decline at the nadir within seconds; automatic generation control
(secondary) restores the frequency to nominal over tens of seconds to
minutes. A shallower nadir means more inertia and faster primary
response. Illustrative trajectory.}
\label{fig:vol04ch14:frequency-response}
\end{figure}


Figure \ref{fig:vol04ch14:frequency-response} traces the frequency
trajectory after a sudden loss of generation. The rotor inertia of the
synchronous fleet sets the initial rate of fall, primary governor
response arrests the decline at the nadir within seconds, and secondary
control restores the nominal value over the following minute. The
nadir is the quantity that decides whether automatic under-frequency
load shedding fires. The single-bus swing-equation model in
\texttt{code/frequency\_response.py} integrates the linearised dynamics
and shows the nadir deepening as the inertia constant $H$ falls, the
regime that high-inverter-penetration grids are moving into; a measured
event trajectory is tabulated in \texttt{data/grid\_frequency\_event.csv}.

\section{Power generation: thermal, hydro, wind, solar, nuclear, geothermal}\pathtag{standard}

A modern grid mixes generation from several primary sources, each
with its own working physics and operational characteristics. Two
numbers characterise any generator before its physics: its
\engterm{capacity factor}, the ratio of the energy it actually produces
in a year to the energy it would produce running at rated power
continuously, and its \engterm{levelized cost of energy}, the constant
per-unit price that makes lifetime discounted revenue equal lifetime
discounted cost. Figure \ref{fig:vol04ch14:capacity-factor} compares
capacity factors across technologies and Figure
\ref{fig:vol04ch14:lcoe-bars} compares the levelized costs; the
\texttt{code/capacity\_factor.py} and \texttt{code/lcoe.py} routines
compute both from the inputs tabulated in
\texttt{data/generation\_mix.csv}.

% Figure: capacity-factor comparison across generation technologies.
% Horizontal bars, percent. Representative values current as of 2024.
\begin{figure}[htbp]
\centering
\begin{tikzpicture}
\begin{axis}[
  width=13cm, height=8cm,
  xlabel={capacity factor (\%, current as of 2024)},
  symbolic y coords={solar PV, wind onshore, wind offshore, hydro, coal, gas combined-cycle, geothermal, nuclear},
  ytick=data,
  xmin=0, xmax=100,
  xbar,
  bar width=11pt,
  nodes near coords,
  nodes near coords style={font=\scriptsize},
  tick label style={font=\scriptsize},
  label style={font=\small},
  enlarge y limits=0.12,
]
\addplot[draw=englime!60!black, fill=englime!50] coordinates {
  (18,solar PV)
  (35,wind onshore)
  (45,wind offshore)
  (40,hydro)
  (50,coal)
  (55,gas combined-cycle)
  (80,geothermal)
  (92,nuclear)
};
\end{axis}
\end{tikzpicture}
\caption[Capacity-factor comparison]{Representative capacity factors by
technology (current as of 2024), the ratio of actual annual energy
produced to the energy that would be produced running at rated power
continuously. Intermittent sources sit low; dispatchable thermal sits
middle; nuclear, run as baseload, sits highest. Capacity factor is the
bridge between a plant's rated power and the energy it actually delivers.
Representative figures only.}
\label{fig:vol04ch14:capacity-factor}
\end{figure}


% Figure: levelized cost of energy (LCOE) comparison across technologies.
% Bars show approximate unsubsidised LCOE midpoints (USD/MWh), current as
% of 2024 (Lazard-style ranges). Range whiskers drawn manually to avoid
% the error-bar shorthand clashing with symbolic coordinates.
\begin{figure}[htbp]
\centering
\begin{tikzpicture}
\begin{axis}[
  width=13cm, height=8cm,
  xlabel={levelized cost of energy (USD/MWh, current as of 2024)},
  symbolic y coords={gaspeaker, nuclearnew, coal, gascc, windon, solarpv},
  yticklabels={solar PV (utility), wind onshore, gas combined-cycle, coal, nuclear (new), gas peaker},
  ytick={solarpv, windon, gascc, coal, nuclearnew, gaspeaker},
  xmin=0, xmax=250,
  xbar,
  bar width=11pt,
  nodes near coords,
  nodes near coords style={font=\scriptsize},
  tick label style={font=\scriptsize},
  label style={font=\small},
  enlarge y limits=0.18,
]
\addplot[draw=engblue, fill=engblue!55] coordinates {
  (42,solarpv)
  (50,windon)
  (75,gascc)
  (110,coal)
  (165,nuclearnew)
  (180,gaspeaker)
};
% range whiskers
\draw[enggray, thick] (axis cs:25,solarpv) -- (axis cs:59,solarpv);
\draw[enggray, thick] (axis cs:27,windon) -- (axis cs:73,windon);
\draw[enggray, thick] (axis cs:50,gascc) -- (axis cs:100,gascc);
\draw[enggray, thick] (axis cs:68,coal) -- (axis cs:152,coal);
\draw[enggray, thick] (axis cs:141,nuclearnew) -- (axis cs:189,nuclearnew);
\draw[enggray, thick] (axis cs:115,gaspeaker) -- (axis cs:245,gaspeaker);
\end{axis}
\end{tikzpicture}
\caption[LCOE comparison across technologies]{Approximate unsubsidised
levelized cost of energy by technology (current as of 2024). Bars give
representative midpoints; grey whiskers give the spread driven by
resource quality, capacity factor, fuel price, and capital cost. Utility
solar and onshore wind are the lowest-cost new generation in most
markets, though their values do not include the system cost of
intermittency (storage, transmission, backup). Order-of-magnitude
figures only.}
\label{fig:vol04ch14:lcoe-bars}
\end{figure}


The capacity factor is a single average, and the average hides the shape
that actually decides what a grid builds. The shape lives in the
\engterm{load-duration curve}: take the $8760$ hourly demands of a year and
re-sort them from the highest hour to the lowest, discarding the time
order. The result is a monotone-decreasing curve whose left edge is the
annual peak, whose area is the annual energy, and whose ratio of mean to
peak is the system load factor (Figure
\ref{fig:vol04ch14:load-duration}). Reading it in horizontal bands assigns
generation to demand. The bottom band, the load present in nearly every
hour of the year, is served cheapest by high-capacity-factor base-load
plant whose capital is repaid over many running hours; nuclear, run-of-river
hydro, and (historically) the most efficient coal sit there. The middle
band, present for a large fraction of the year but not all of it, is served
by intermediate cycling plant, typically combined-cycle gas. The thin
sliver at the top left, the peak that appears for only a few hundred hours a
year, is served by low-capacity-factor peaking plant and by demand
response, because no plant whose capital must be repaid over many hours can
justify itself on a few hundred hours of running.

% Figure: load-duration curve. The chronological load is re-sorted from
% highest to lowest hour, producing a monotone curve whose area is the
% annual energy and whose shape sets the capacity factor of each layer of
% the dispatch stack. Horizontal bands show base, intermediate, and peak.
% Values illustrative for a temperate grid (current as of 2025).
\begin{figure}[htbp]
\centering
\begin{tikzpicture}
\begin{axis}[
  width=13cm, height=8cm,
  xlabel={hours of the year sorted by load (\% of $8760$ h)},
  ylabel={load (GW)},
  xmin=0, xmax=100, ymin=0, ymax=42,
  xtick={0,20,40,60,80,100},
  ytick={0,10,20,30,40},
  axis lines=left,
  tick label style={font=\scriptsize},
  label style={font=\small},
  legend style={font=\scriptsize, at={(0.97,0.97)}, anchor=north east, draw=none},
]
% monotone-decreasing load-duration curve
\addplot[engblue, very thick, smooth] coordinates {
  (0,40)(2,35)(5,31)(10,28)(20,25)(30,23)(40,21.5)(50,20)(60,18.5)(70,17)(80,15.5)(90,14)(97,12)(100,10)
};
\addlegendentry{load-duration curve}
% base-load band (served all year)
\draw[enggray, dashed] (axis cs:0,12) -- (axis cs:100,12);
\node[anchor=west, font=\scriptsize, color=enggray] at (axis cs:2,9.6) {base load (high CF)};
% intermediate band
\draw[enggray, dashed] (axis cs:0,25) -- (axis cs:100,25);
\node[anchor=west, font=\scriptsize, color=enggray] at (axis cs:2,18.6) {intermediate};
% peak band
\node[anchor=west, font=\scriptsize, color=engred] at (axis cs:5,33.5) {peak (low CF)};
% capacity-factor reading: average load marker
\addplot[engorange, very thick, dashed] coordinates {(0,21)(100,21)};
\node[anchor=west, font=\scriptsize, color=engorange] at (axis cs:55,23) {mean load $\to$ system CF};
\end{axis}
\end{tikzpicture}
\caption[Load-duration curve]{The load-duration curve re-sorts the
$8760$ hourly loads of a year from highest to lowest. Its area is the
annual energy, its height at the left is the peak demand, and the ratio of
its mean (the orange line) to its peak is the system load factor. Reading
the curve in horizontal bands assigns generation to it: the band served
for almost every hour is the base load that high-capacity-factor plant
covers, the middle band is intermediate cycling plant, and the thin sliver
at the top left is the peak that low-capacity-factor peaking plant and
demand response serve for a few hundred hours a year. The shape, not the
total, decides how much of each kind of capacity a grid needs. Values
illustrative of a temperate grid (current as of 2025).}
\label{fig:vol04ch14:load-duration}
\end{figure}


The load-duration curve also fixes the relationship between capacity factor
and position in the dispatch order, and the relationship runs the other
way from the intuition that a plant's capacity factor is a property of the
plant. A peaking plant has a low capacity factor not because it cannot run
continuously but because the load-duration curve gives it only the few
hundred hours at the top to run in; the same machine moved to the bottom of
a growing grid would run flat out. The capacity factor of a dispatchable
plant is set jointly by the plant (its availability and its place in the
merit order) and by the shape of the load-duration curve it serves. For a
variable renewable the logic inverts: its capacity factor is set by the
resource, not by the load, so its energy lands wherever the weather puts it
on the time axis, which is why the renewable capacity factor and the
capacity credit of the later mastery box are different numbers. The area
under the load-duration curve down to a given power level is the energy
that a band of capacity at that level would serve, and integrating the
curve is how a planner reads off how many gigawatts of each kind of plant
the demand actually calls for. The same construction, applied to net load
(demand minus variable renewables) rather than gross demand, is the tool
that the duck curve later in the chapter is one slice of: high renewable
penetration reshapes the net-load-duration curve, hollowing its middle and
steepening its top, and it is that reshaped curve, not the gross one, that
the dispatchable fleet must be built to serve.

\engterm{Thermal generation} from fossil fuels (coal, natural gas,
fuel oil) burns the fuel to raise steam (Rankine cycle, Chapter 3 of
this volume) or to drive a gas turbine (Brayton cycle). A modern
supercritical-coal plant achieves $\eta\sim 45\,\%$ thermal-to-
electric efficiency at $\sim 30\,\si{\mega\pascal}$ steam pressure
and $\sim 600\,\si{\celsius}$ steam temperature. A combined-cycle gas
turbine (Brayton-Rankine cascade, where the Brayton exhaust drives a
Rankine cycle) reaches $\eta\sim 64\,\%$ in the largest modern units,
the highest practical thermal efficiency of any single fossil-fuel
generator. The trade-off is fuel cost and carbon emissions; combined-
cycle gas is the marginal generator on most modern grids, setting
electricity prices in many wholesale markets (current as of 2025).

\engterm{Hydroelectric generation} extracts gravitational potential
energy from a falling water column through a turbine. The power is
$P = \eta\rho g Q H$, where $Q$ is the volumetric flow rate, $H$ is
the head (height of fall), and $\eta$ is the conversion efficiency
(typically $0.9$ to $0.95$ for modern Francis or Kaplan turbines, and
$0.85$ to $0.92$ for Pelton wheels at high heads). A $1\,\si{\meter
\cubed\per\second}$ flow at $100\,\si{\meter}$ head produces about
$1\,\si{\mega\watt}$. Hydroelectricity supplies $\sim 15\,\%$ of
global electricity generation (current as of 2024); pumped-storage
hydroelectricity is the dominant grid-scale storage technology by
energy capacity.

\engterm{Wind generation} extracts kinetic energy from a moving air
mass with a horizontal-axis turbine. The power is
\[
P = \tfrac{1}{2}C_{p}\rho A v^{3},
\]
where $C_{p}$ is the power coefficient (limited by the Betz
limit at $C_{p,\max} = 16/27 \approx 0.593$), $\rho$ is the air
density, $A$ is the rotor-swept area, and $v$ is the wind speed.
Modern utility-scale wind turbines (Vestas V150-4.2~MW or similar,
current as of 2024) have rotor diameters of $\sim 150\,\si{\meter}$,
rated power $\sim 4\,\si{\mega\watt}$, and capacity factors of $\sim
30\,\%$ onshore, $\sim 45\,\%$ offshore. The cube dependence on wind
speed means a doubling of average wind speed multiplies the power
yield by $8$; siting matters enormously.

% Figure: the Betz power coefficient as a function of the axial induction
% factor a. C_p(a) = 4a(1-a)^2, maximised at a=1/3 giving 16/27.
\begin{figure}[htbp]
\centering
\begin{tikzpicture}
\begin{axis}[
  width=12cm, height=7.5cm,
  xlabel={axial induction factor $a$},
  ylabel={power coefficient $C_p = 4a(1-a)^2$},
  xmin=0, xmax=1, ymin=0, ymax=0.65,
  xtick={0,0.2,0.333,0.5,0.8,1.0},
  xticklabels={0,0.2,$\tfrac{1}{3}$,0.5,0.8,1.0},
  ytick={0,0.2,0.4,0.593},
  yticklabels={0,0.2,0.4,$\tfrac{16}{27}$},
  axis lines=left,
  tick label style={font=\scriptsize},
  label style={font=\small},
  samples=100,
  domain=0:1,
]
\addplot[engblue, very thick] {4*x*(1-x)^2};
% Betz maximum
\draw[engred, dashed] (axis cs:0.3333,0) -- (axis cs:0.3333,0.5926);
\draw[engred, dashed] (axis cs:0,0.5926) -- (axis cs:0.3333,0.5926);
\filldraw[engred] (axis cs:0.3333,0.5926) circle [radius=1.8pt];
\node[engred, font=\scriptsize, anchor=south west] at (axis cs:0.34,0.55)
  {Betz limit $16/27\approx0.593$};
\end{axis}
\end{tikzpicture}
\caption[The Betz power coefficient]{The power coefficient of an ideal
actuator-disc wind turbine, $C_p = 4a(1-a)^2$, as a function of the
axial induction factor $a$ (the fractional slowing of the wind at the
disc). The maximum sits at $a = 1/3$, where $C_p = 16/27 \approx 0.593$:
no horizontal-axis turbine can extract more than this fraction of the
kinetic energy in the oncoming wind. Slowing the wind too much (large
$a$) reduces the mass flux through the disc and lowers the extracted
power.}
\label{fig:vol04ch14:betz-curve}
\end{figure}


The power coefficient $C_p$ is bounded above by the Betz limit, derived
in the exercises by treating the rotor as an actuator disc that slows
the wind from upstream to downstream. Figure
\ref{fig:vol04ch14:betz-curve} plots $C_p = 4a(1-a)^2$ against the
axial induction factor $a$; the maximum at $a = 1/3$ gives the famous
$16/27 \approx 0.593$. Real turbines reach $C_p \approx 0.45$ to
$0.50$ at their best operating point, below the ideal because of tip
losses, drag, and wake rotation.

\engterm{Solar photovoltaic generation} converts sunlight directly to
electricity through pn-junction cells (Chapter 12). At standard
$1\,\text{sun}$ AM1.5G illumination, modern silicon modules produce
$\sim 200\,\si{\watt\per\meter\squared}$ peak power at $\sim 20\,\%$
efficiency. Capacity factor varies from $\sim 12\,\%$ in cloudy
northern Europe to $\sim 25\,\%$ in the southwestern United States
and the Middle East. The technology has experienced dramatic cost
reduction: utility-scale solar PV electricity cost is now $\sim
\$25$ to $\$60\,\si{\per\mega\watt\hour}$ in 2024, below the cost of
new coal or nuclear generation in many markets (Lazard LCOE
analysis, current as of 2024).

\engterm{Concentrated solar thermal} uses mirrors to focus sunlight
onto a working fluid (molten salt, oil) that drives a Rankine cycle.
The advantage over PV is built-in thermal storage; the disadvantage
is requiring direct (cloud-free) sunlight and higher capital cost.
Concentrated solar has not scaled at PV's rate; commercial deployment
peaked $\sim 2015$ and has been roughly flat since.

The physics divides concentrated solar sharply from photovoltaics on two
counts, and the division decides where each belongs. The first is that a
concentrator works only on the \engterm{direct} beam, the light arriving
straight from the sun's disc, because only a parallel beam can be focused to a
point or a line; the \engterm{diffuse} light scattered by cloud and haze
cannot be concentrated and is lost to the receiver. A photovoltaic panel, by
contrast, collects both the direct and the diffuse components, so it produces
usefully under an overcast sky while a concentrator produces almost nothing.
This is why concentrated solar is confined to the high-direct-beam deserts and
why its output collapses on a hazy day that a PV array rides through. The
second is that a concentrator raises a working fluid to a high temperature,
$\sim 400$ to $565\,\si{\celsius}$ for molten salt, and a high-temperature
store buys the one thing PV lacks without a battery: cheap thermal storage.
The Carnot ceiling rewards the high temperature, so the Rankine cycle behind a
solar tower reaches $\sim 40\,\%$ heat-to-electric efficiency, and the hot
salt can be held in an insulated tank for hours at a low energy cost, then run
through the turbine after sunset. A tower plant with a large salt store can
therefore dispatch into the evening peak the way a battery-backed PV plant
does, but the store is a tank of hot salt rather than a bank of cells, with
the low energy cost of a thermal reservoir. The reason concentrated solar
nonetheless lost the race to PV is economic rather than physical: PV's
manufacturing learning curve drove its cost down an order of magnitude while
the concentrator's mirrors, tracking, receiver, and thermal plant stayed
site-built and capital-heavy, so a PV array plus a lithium-ion battery came to
undercut a tower plant at the same dispatchability. Concentrated solar retains
a niche where very-high-temperature process heat or long thermal storage is
the product rather than electricity, and it is the clearest case in the
chapter of a technology beaten not on its physics but on the cost curve of a
rival.

\engterm{Nuclear generation} uses controlled fission of uranium (and
in some designs plutonium or thorium) to raise steam for a Rankine
cycle. Modern light-water reactors (PWR, BWR) operate at thermal
efficiencies of $\sim 33\,\%$ and electrical outputs of $\sim 1\,\si{
\giga\watt}$ per unit. The fuel cycle is engineered for $\sim 18$-
month operating cycles with planned refuelling. Capacity factors of
modern units exceed $90\,\%$, the highest of any non-storage
generation technology. The economic and political constraints are
the high capital cost (\$5 to \$15 per watt of installed capacity)
and waste-management requirements; the operational record outside
of two named events (Chernobyl 1986, Fukushima 2011) is one of low
fatality rate per kilowatt-hour generated.

\engterm{Geothermal generation} extracts thermal energy from
high-temperature aquifers or hot dry rock and drives a Rankine cycle.
Conventional hydrothermal geothermal is limited to a few favourable
regions (Iceland, parts of California and Nevada, the Philippines,
Indonesia); enhanced geothermal systems (engineered fracture networks
in hot dry rock) are in early commercial deployment and are
geographically more flexible. Geothermal supplies $\sim 0.4\,\%$ of
global electricity generation (current as of 2024); the engineering
challenge is the high upfront drilling cost and the working-life
declines from reservoir depletion.

The physics that sets geothermal apart from the other Rankine sources is the
modest temperature of the heat. A geothermal reservoir delivers fluid at
$\sim 150$ to $300\,\si{\celsius}$, far below the $\sim 600\,\si{\celsius}$
of a supercritical-coal boiler, and the Carnot ceiling scales with the
temperature difference. A reservoir at $200\,\si{\celsius}$ ($473\,\si{
\kelvin}$) rejecting to a $30\,\si{\celsius}$ ($303\,\si{\kelvin}$)
condenser has a Carnot efficiency of $1 - 303/473 \approx 0.36$, and a real
binary-cycle plant reaches perhaps a third to a half of that, so a thermal-
to-electric efficiency of $\sim 10$ to $15\,\%$, well below a fossil or
nuclear Rankine plant. The low efficiency is not a defect to engineer away
but a consequence of the resource temperature, and it means a geothermal
plant must move a large mass of fluid per unit of electricity: at
$\eta \approx 0.12$ a $50\,\si{\mega\watt}$ plant must extract $\sim
420\,\si{\mega\watt}$ of heat from the reservoir, and at a fluid enthalpy
drop of $\sim 0.4\,\si{\mega\joule\per\kilogram}$ across the plant that is a
production flow of $\sim 1000\,\si{\kilogram\per\second}$, several large wells
running continuously. The wells are the cost: drilling scales roughly
exponentially with depth because the rock is harder, hotter, and slower to
penetrate, so a well to $3\,\si{\kilo\meter}$ might cost a few million dollars
and one to $5\,\si{\kilo\meter}$ several times that, which is why the
enhanced-geothermal ambition of reaching hot dry rock anywhere is a drilling-
cost problem before it is a reservoir problem. The reservoir also cools as it
is worked: pumping heat out faster than the surrounding rock conducts it back
draws the production temperature down over years, the \engterm{thermal
drawdown} that sets a geothermal field's working life, and the operator
manages it by reinjecting the spent cooled fluid to sweep fresh heat from a
wider volume of rock and by spacing the production and injection wells so the
cooled reinjection front reaches the producers slowly. Geothermal is
therefore the chapter's clearest case of a resource whose engineering is
governed by a slow underground transport process rather than by the power
plant at the surface.

\engterm{Biomass generation} burns plant matter (wood, agricultural residues,
dedicated energy crops, or the biogas from anaerobic digestion) to raise steam
for a Rankine cycle, so its power plant is a conventional thermal unit and its
distinctive engineering lives in the fuel rather than the boiler. Two
properties set biomass apart from the fossil fuels it can substitute for. The
first is energy density: dry wood carries $\sim 18\,\si{\mega\joule\per
\kilogram}$ against coal's $\sim 24$ to $30$, and green wood or wet residue
carries far less because the water in it must be boiled off before the fuel
burns, so a biomass plant handles a large mass of low-density fuel and its
economics are dominated by the cost and distance of collecting it. The
collection radius is the binding constraint: transporting a diffuse fuel more
than $\sim 50$ to $100\,\si{\kilo\meter}$ costs more energy and money than the
fuel returns, which caps a biomass plant at a modest size, typically tens of
megawatts, far below a coal or nuclear unit, because the fuel cannot be
gathered from a wide enough area to feed a large boiler. The second is the
carbon accounting: biomass is counted as low-carbon only if the plant matter
regrows and reabsorbs the carbon its combustion released, a balance that holds
for genuine residues and fast-rotation crops and fails for slow-growing forest
burned faster than it regrows, so the carbon merit of biomass is a
land-and-time question rather than a combustion one. Biomass earns its place in
the mix as a dispatchable, firm, storable-fuel renewable, the one variable-free
low-carbon thermal source, which is why it is valued for the firm capacity the
capacity-credit argument prices even where its energy is dearer than wind or
solar.

\engterm{Marine generation} extracts energy from the movement of the sea, and
its two working forms sit at opposite ends of the predictability the rest of the
chapter has cared about. \engterm{Tidal} generation, whether a barrage across
an estuary or a free-stream turbine in a tidal current, draws on the
gravitational pull of the moon and sun, so its resource is almost perfectly
predictable years ahead, a rare virtue among renewables, though it delivers in
the twice-daily pulses of the tide rather than on demand, and it is confined to
the few sites with a large tidal range or a fast tidal stream. The power in a
tidal stream follows the same cube law as wind, $P = \tfrac12 C_p \rho A v^3$,
but with seawater's density $\sim 800$ times that of air, so a slow
$2.5\,\si{\meter\per\second}$ current carries a useful power density and a
tidal turbine is smaller than a wind turbine of the same rating. \engterm{Wave}
generation extracts energy from the surface motion of wind-driven waves, a
denser and more persistent resource than the wind that made it but a mechanically
harsh one, because the device must survive the storm waves that carry many times
the energy of the average sea while capturing the ordinary ones, and the
survivability-versus-capture trade has kept wave power from scaling. Marine
generation supplies a negligible share of global electricity today and is the
chapter's clearest case of a resource that is physically real and partly
predictable but held back by the engineering cost of building machines that
survive the sea.

\section{Energy storage: batteries, pumped hydro, thermal, hydrogen}\pathtag{standard}

The growth of intermittent renewable generation (wind, solar PV)
makes energy storage essential; without storage, the grid must
either curtail generation when supply exceeds demand or maintain
high-capital-cost dispatchable backup generators.

% Figure: renewable output variability over a week. Wind and solar
% normalized output (fraction of rated) showing the uncorrelated,
% intermittent character that storage and dispatch must absorb.
\begin{figure}[htbp]
\centering
\begin{tikzpicture}
\begin{axis}[
  width=13cm, height=8cm,
  xlabel={day of week},
  ylabel={output (fraction of rated capacity)},
  xmin=0, xmax=7, ymin=0, ymax=1,
  xtick={0,1,2,3,4,5,6,7},
  xticklabels={Mon,Tue,Wed,Thu,Fri,Sat,Sun,},
  ytick={0,0.25,0.5,0.75,1.0},
  axis lines=left,
  legend pos=north east,
  legend cell align=left,
  tick label style={font=\scriptsize},
  label style={font=\small},
  legend style={font=\scriptsize},
]
% solar: daily bells, zero at night, reduced on cloudy days (Wed, Thu)
\addplot[engamber, thick, smooth] coordinates {
  (0,0)(0.25,0.55)(0.5,0)(1,0)(1.25,0.6)(1.5,0)(2,0)(2.25,0.2)(2.5,0)
  (3,0)(3.25,0.15)(3.5,0)(4,0)(4.25,0.58)(4.5,0)(5,0)(5.25,0.6)(5.5,0)
  (6,0)(6.25,0.5)(6.5,0)(7,0)
};
\addlegendentry{solar}
% wind: slow swings, high on the cloudy/stormy days
\addplot[englime!60!black, thick, smooth] coordinates {
  (0,0.3)(0.5,0.45)(1,0.35)(1.5,0.2)(2,0.4)(2.5,0.7)(3,0.85)(3.5,0.8)
  (4,0.55)(4.5,0.3)(5,0.25)(5.5,0.35)(6,0.5)(6.5,0.4)(7,0.3)
};
\addlegendentry{wind}
\end{axis}
\end{tikzpicture}
\caption[Renewable output variability over a week]{Normalised wind and
solar output across a week. Solar follows a clean daily cycle modulated
by cloud (the low Wednesday and Thursday). Wind varies on the multi-day
scale of weather systems and often rises when solar falls, as on the
cloudy stormy days here. The partial anti-correlation is one reason
mixed wind-and-solar portfolios are smoother than either alone, though
neither removes the need for storage and dispatch. Illustrative.}
\label{fig:vol04ch14:renewable-variability}
\end{figure}


Figure \ref{fig:vol04ch14:renewable-variability} shows the source of
the problem: wind and solar output vary on time scales from minutes to
days, and the variation is only partly predictable. Solar follows a
clean daily cycle modulated by cloud; wind swings on the multi-day
scale of weather systems. The two are partly anti-correlated (wind
often rises when a front blocks the sun), so a mixed portfolio is
smoother than either alone, but neither removes the need to move energy
in time. The storage technologies in \texttt{data/storage\_technologies.csv}
span the range from minute-scale to seasonal.

% Figure: battery charge/discharge profile over a day, with state of
% charge tracking the integral. Left axis = power (kW, charge positive),
% right axis = state of charge (kWh).
\begin{figure}[htbp]
\centering
\begin{tikzpicture}
\begin{axis}[
  width=13cm, height=5.5cm,
  name=power,
  xlabel={hour of day},
  ylabel={power (kW)},
  xmin=0, xmax=24, ymin=-6, ymax=6,
  xtick={0,4,8,12,16,20,24},
  ytick={-5,0,5},
  axis lines=left,
  tick label style={font=\scriptsize},
  label style={font=\small},
]
% charge (positive) midday from surplus PV; discharge (negative) evening
\addplot[engamber, very thick, const plot] coordinates {
  (0,0)(8,0)(9,2)(10,4)(11,5)(12,5)(13,4)(14,3)(15,1)(16,0)(17,0)(18,-4)(19,-5)(20,-4)(21,-2)(22,0)(24,0)
};
\draw[enggray, dashed] (axis cs:0,0) -- (axis cs:24,0);
\node[engamber, font=\scriptsize, anchor=south] at (axis cs:12,5) {charge};
\node[engamber, font=\scriptsize, anchor=north] at (axis cs:19,-5) {discharge};
\end{axis}
\begin{axis}[
  width=13cm, height=5.5cm,
  at={(power.below south)}, anchor=north, yshift=-0.6cm,
  xlabel={hour of day},
  ylabel={state of charge (kWh)},
  xmin=0, xmax=24, ymin=0, ymax=12,
  xtick={0,4,8,12,16,20,24},
  ytick={0,5,10},
  axis lines=left,
  tick label style={font=\scriptsize},
  label style={font=\small},
]
% SoC rising during charge, falling during discharge
\addplot[engblue, very thick, smooth] coordinates {
  (0,2)(8,2)(10,4.5)(12,8.5)(14,10)(16,10.2)(18,9.8)(19,5)(20,2.5)(21,1.5)(22,1.5)(24,1.5)
};
\draw[engred, dashed] (axis cs:0,10.5) -- (axis cs:24,10.5);
\node[engred, font=\scriptsize, anchor=south east] at (axis cs:24,10.5) {usable capacity};
\end{axis}
\end{tikzpicture}
\caption[Battery charge/discharge profile]{A daily battery cycle. The
upper panel shows charging power (positive, midday from surplus PV) and
discharging power (negative, early evening to serve the net-load peak).
The lower panel shows the state of charge as the running integral of
power, bounded above by the usable capacity. Round-trip losses mean the
discharged energy is below the charged energy. Values are illustrative.}
\label{fig:vol04ch14:storage-profile}
\end{figure}


Figure \ref{fig:vol04ch14:storage-profile} shows the daily cycle a
grid-scale or household battery performs to flatten the net load:
charge from the midday surplus, discharge into the early-evening peak.
The state of charge is the running integral of power, bounded above by
the usable capacity, and the discharged energy falls below the charged
energy by the round-trip loss.

\engterm{Pumped-storage hydroelectricity} is the dominant grid-scale
storage technology by energy capacity. Water is pumped from a lower
reservoir to an upper reservoir during periods of cheap electricity
(typically overnight or on a high-renewable-output day) and released
through a turbine to generate during high-demand periods. Round-trip
efficiency is $70$ to $80\,\%$; capacity factors of dedicated
pumped-storage plants are $\sim 15\,\%$. The standard installation
size is $1$ to $3\,\si{\giga\watt}$ over $4$ to $12\,\si{\hour}$
discharge. The geographic constraints are severe: two large
reservoirs at substantially different elevation in close proximity.

\engterm{Lithium-ion batteries} are the dominant grid-scale storage
technology by power and by recent deployment growth. A modern
utility-scale battery (Tesla Megapack, BYD CATL utility products) has
energy density of $\sim 250\,\si{\watt\hour\per\kilogram}$ at the
cell level, $\sim 150\,\si{\watt\hour\per\kilogram}$ at the system
level. Round-trip efficiency is $85$ to $90\,\%$ at moderate
charge-discharge rates. Cycle life at $\sim 80\,\%$ depth of discharge
is $5000$ to $8000$ cycles, equivalent to $15$ to $20$ years at one
cycle per day. Costs have fallen from \$1100 per kWh
in 2010 to $\sim$\$140 per kWh in 2024 (BloombergNEF, current as
of 2024). The chemistry of lithium-ion cells matters: LFP (lithium
iron phosphate) is the modern stationary-storage standard for its
safety and cycle life; NMC (nickel-manganese-cobalt) is the modern
EV-traction standard for its higher energy density.

\engterm{Thermal storage} is well-suited for concentrated solar
plants and for industrial heat. Molten salt at $\sim 565\,\si{\celsius}$
stores $\sim 150\,\si{\watt\hour\per\kilogram}$ thermally and is
discharged through a steam-Rankine cycle, recovering electricity at
$\sim 40\,\%$ efficiency (round-trip electrical) or $\sim 95\,\%$
efficiency (heat-to-heat, where the heat is the end product, such as
in district heating or industrial process heat). The Ivanpah and
Cerro Dominador plants are working installations.

\engterm{Hydrogen} is a chemical energy carrier produced by water
electrolysis from electrical input, stored as compressed gas
($\sim 70\,\si{\mega\pascal}$), liquid ($-253\,\si{\celsius}$), or
in ammonia or metal-hydride form, and recovered by combustion in a
fuel cell or in a turbine. Round-trip electrical efficiency is $\sim
40\,\%$ for state-of-the-art (alkaline or PEM electrolyser + PEM
fuel cell, current as of 2024); the wide use of hydrogen as a
storage medium is constrained by this low efficiency. Hydrogen is
better suited to long-duration storage (seasonal scale, $\sim$weeks
to months) where the per-kilowatt-hour energy cost dominates over
the round-trip-efficiency cost, and to industrial-feedstock uses
(ammonia synthesis, steelmaking) where hydrogen replaces fossil
methane or coke.

\begin{mastery}
\textbf{Power and energy are priced separately.} A storage system has
two capacities that scale and cost independently: the power rating
(how fast it can charge or discharge, $\si{\mega\watt}$) and the energy
rating (how much it can hold, $\si{\mega\watt\hour}$). Their ratio is
the discharge duration. Lithium-ion costs are dominated by the energy
(the cells), so a $4\,\si{\hour}$ system is roughly twice the cost of a
$2\,\si{\hour}$ system of the same power; pumped hydro and hydrogen
have a high fixed power cost (turbines, electrolysers) and a low
marginal energy cost (a bigger reservoir, a bigger tank), so their
cost-per-hour falls with duration. The crossover is why lithium-ion
wins the minutes-to-hours market and pumped hydro or hydrogen win the
hours-to-seasons market. A round-trip chain compounds its losses
multiplicatively: the green-hydrogen path of electrolysis,
compression, transport, and fuel cell multiplies to $\sim 0.65 \times
0.90 \times 0.95 \times 0.55 \approx 0.31$, so roughly two thirds of
the input electricity is lost. The \texttt{code/storage\_sizing.py}
routine computes both the chain efficiency and the energy a multi-day
shortfall demands.
\end{mastery}

\section{Demand-side management and grid control}\pathtag{standard}

The classical grid was generation-following-load: generation was
adjusted in real time to meet whatever demand presented itself. The
modern grid increasingly is also load-following-generation: demand is
shifted in time to align with available generation (especially
intermittent renewables). The mechanisms are collectively known as
\engterm{demand-side management} (DSM) and \engterm{demand response}
(DR).

% Figure: daily load curve and the net-load "duck curve" produced by
% subtracting midday solar PV from gross demand. Hours on the x-axis,
% power (GW) on the y-axis. Illustrative shapes, not a specific grid.
\begin{figure}[htbp]
\centering
\begin{tikzpicture}
\begin{axis}[
  width=13cm, height=8cm,
  xlabel={hour of day},
  ylabel={power (GW)},
  xmin=0, xmax=24, ymin=0, ymax=45,
  xtick={0,3,6,9,12,15,18,21,24},
  ytick={0,10,20,30,40},
  axis lines=left,
  legend pos=north west,
  legend cell align=left,
  tick label style={font=\scriptsize},
  label style={font=\small},
  legend style={font=\scriptsize},
]
% gross demand: morning ramp, midday plateau, evening peak
\addplot[engblue, very thick, smooth] coordinates {
  (0,22) (3,19) (6,21) (8,30) (10,33) (12,34) (14,33) (16,34) (18,40) (20,38) (22,30) (24,23)
};
\addlegendentry{gross demand}
% solar PV output: bell from sunrise to sunset
\addplot[engamber, thick, dashed, smooth] coordinates {
  (0,0) (5,0) (6,1) (8,7) (10,16) (12,22) (14,19) (16,11) (18,3) (19,0) (24,0)
};
\addlegendentry{solar PV output}
% net load = gross - solar (the duck): belly at midday, neck at evening
\addplot[engred, very thick, smooth] coordinates {
  (0,22) (3,19) (6,20) (8,23) (10,17) (12,12) (14,14) (16,23) (18,37) (20,38) (22,30) (24,23)
};
\addlegendentry{net load (duck)}
\node[engred, font=\scriptsize, anchor=north] at (axis cs:12,11) {belly};
\node[engred, font=\scriptsize, anchor=south] at (axis cs:18.5,38) {neck (steep ramp)};
\end{axis}
\end{tikzpicture}
\caption[Daily load curve and the duck curve]{The daily gross-demand
curve (blue), the midday solar-PV output (amber), and the resulting net
load (red) the dispatchable fleet must serve. Subtracting midday solar
hollows out the belly and steepens the early-evening ramp as the sun
sets while demand peaks; the shape is the ``duck curve'' that high solar
penetration imposes on the rest of the fleet. Shapes are illustrative.}
\label{fig:vol04ch14:duck-curve}
\end{figure}


The operational case for demand-side management is written in the shape
of the net load. Figure \ref{fig:vol04ch14:duck-curve} shows the
``duck curve'' that high solar penetration imposes: subtracting midday
solar from gross demand hollows out the middle of the day and steepens
the early-evening ramp, when the sun sets while demand peaks. The
dispatchable fleet must follow that ramp, and the steeper it grows, the
more flexible (and expensive) the generation and storage that serve it.

% Figure: generation dispatch stack over a day. Merit-order layers
% (nuclear baseload, hydro, wind, solar, gas) stacked to meet demand.
% Stacked area plot via cumulative coordinates.
\begin{figure}[htbp]
\centering
\begin{tikzpicture}
\begin{axis}[
  width=13cm, height=8cm,
  xlabel={hour of day},
  ylabel={cumulative generation (GW)},
  xmin=0, xmax=24, ymin=0, ymax=45,
  xtick={0,4,8,12,16,20,24},
  ytick={0,10,20,30,40},
  axis lines=left,
  stack plots=y,
  area style,
  legend pos=north west,
  legend cell align=left,
  tick label style={font=\scriptsize},
  label style={font=\small},
  legend style={font=\scriptsize, fill opacity=0.85, draw=none},
]
% nuclear baseload, flat
\addplot[draw=engblue, fill=engblue!55] coordinates {
  (0,10)(4,10)(8,10)(12,10)(16,10)(20,10)(24,10)} \closedcycle;
\addlegendentry{nuclear (baseload)}
% hydro, mild
\addplot[draw=enggray, fill=enggray!40] coordinates {
  (0,4)(4,4)(8,5)(12,5)(16,5)(20,6)(24,4)} \closedcycle;
\addlegendentry{hydro}
% wind, variable
\addplot[draw=englime!60!black, fill=englime!50] coordinates {
  (0,6)(4,7)(8,5)(12,4)(16,5)(20,7)(24,7)} \closedcycle;
\addlegendentry{wind}
% solar, midday bell
\addplot[draw=engamber, fill=engamber!50] coordinates {
  (0,0)(4,0)(8,3)(12,9)(16,5)(20,0)(24,0)} \closedcycle;
\addlegendentry{solar}
% gas, the flexible top of the stack filling the gap to demand
\addplot[draw=engred, fill=engred!45] coordinates {
  (0,2)(4,0)(8,7)(12,6)(16,9)(20,15)(24,5)} \closedcycle;
\addlegendentry{gas (flexible)}
\end{axis}
\end{tikzpicture}
\caption[Generation dispatch stack over a day]{A generation dispatch
stack over a day. Low-marginal-cost sources (nuclear, hydro, wind,
solar) fill the bottom of the stack; flexible gas fills the gap between
the must-run and variable layers and the demand curve. The dispatchable
gas layer ramps hardest in the early evening, when solar falls and
demand rises. Quantities are illustrative.}
\label{fig:vol04ch14:dispatch-stack}
\end{figure}


Figure \ref{fig:vol04ch14:dispatch-stack} shows the complementary
picture from the generation side: the dispatch stack, in which
low-marginal-cost must-run and variable sources fill the bottom and
flexible gas fills the gap to the demand curve. The gas layer ramps
hardest exactly where the duck's neck is steepest. Demand-side
management attacks the same problem from the load side, flattening the
net-load curve so the stack needs less flexible generation at the top.

The simplest DSM mechanism is \engterm{time-of-use pricing}: the
retail electricity tariff varies with the wholesale supply-demand
balance, with higher prices during peak hours and lower prices
overnight. Consumers respond by scheduling discretionary loads
(dishwasher, dryer, EV charging, water heating) to off-peak periods.
Even modest price differentials of $5$:$1$ between peak and off-peak
have measurable effects on residential load shapes; commercial and
industrial consumers are more responsive.

\engterm{Direct load control} (DLC) is the utility-side mechanism by
which the utility, with the consumer's prior consent, briefly
interrupts power to specific large loads (water heaters, air
conditioner compressors) during grid emergencies. The aggregated
capacity of a residential DLC programme can reach hundreds of
megawatts in a large utility. The mechanism is reliable but coarse-
grained; it is one of the working tools of last-resort grid
reliability.

\engterm{Smart meters and home energy management} provide the
information backbone for fine-grained DSM. A smart meter reports
interval consumption (typically $15$- or $30$-minute intervals) and
supports time-of-use billing; a home energy management system
coordinates appliances, EV charging, and rooftop PV in response to
prices and operational signals.

\engterm{Aggregators and virtual power plants} (VPPs) are commercial
entities that contract with thousands of distributed energy resources
(batteries, EVs, controllable loads, rooftop PV) and bid the
aggregated capacity into wholesale markets as if it were a single
generator. The VPP business model has matured rapidly: the global
VPP capacity, current as of 2024, exceeds $50\,\si{\giga\watt}$.

\begin{archetype}[Balance, transport, scaling, control]
The energy system is the four archetypes of this volume operating at
the largest scale. A \engterm{balance} relates instantaneous
generation to instantaneous demand at every moment. A
\engterm{transport} chain moves energy from where it is generated
to where it is used over distances of metres to thousands of
kilometres. \engterm{Scaling} laws govern the trade-off between
distributed and centralised generation, the economics of grid
expansion, and the time-and-space averaging of renewable variability.
\engterm{Control} mechanisms operate at every time scale from
millisecond to year. The discipline of energy-systems engineering is
the integration of all four at the level of a national grid.
\end{archetype}

\section{Off-grid and microgrid systems}\pathtag{standard}

Not every electrical consumer is connected to a synchronous AC grid.
The chapter's working categories are the off-grid system, which stands alone,
and the microgrid, which can stand alone or connect. What separates them from
the bulk grid is the loss of the vast synchronous machine that absorbs
disturbances. The bulk grid holds its frequency and voltage steady because
thousands of generators share every disturbance, and the inertia of their
combined spinning mass swamps any single event; a small isolated system has
neither the shared reserve nor the inertia, so a single load switching on is a
large fraction of its capacity, and the control that the bulk grid gets for
free from its scale must be engineered explicitly. An isolated system's
frequency swings harder after a disturbance because its inertia constant is
small, its voltage is more sensitive to reactive load because it has fewer
sources to hold it, and its reserve margin against the loss of a single unit is
a large percentage because the units are coarse relative to the load, the same
effect the small-system loss-of-load example showed. The engineering of an
isolated system is therefore the engineering of stability at small scale:
enough fast storage or a grid-forming inverter to supply synthetic inertia and
hold the frequency, enough reactive capability to hold the voltage, and a
control architecture that can black-start the system from nothing, because
there is no neighbouring grid to energise it after an outage.

An \engterm{off-grid system} provides all the energy needs of an
isolated load from local generation and storage. The canonical small
off-grid system is a rural household in a developing country with
$\sim 100\,\si{\watt}$ rooftop PV, $\sim 1\,\si{\kilo\watt\hour}$
battery, an inverter, and $\sim 100$ to $200\,\si{\watt}$ of
controllable LED lighting, mobile-phone charging, and a small TV.
The cost (current as of 2024) is $\sim$\$300 to \$500 installed.
Larger off-grid systems serve mines, remote installations,
expedition stations, and ships.

A \engterm{microgrid} is a small electrical network with its own
generation and storage, capable of operating either connected to the
larger grid or in isolated (islanded) mode. Microgrids are deployed
on military bases (for resilience to grid outage), in universities
and hospitals (for power-quality and outage independence), and in
remote communities (where the alternative is diesel generation only).
The control architecture must handle a wide dynamic range: from
grid-connected operation at near-zero local response to
islanded-operation full-control of frequency and voltage by the
microgrid's own resources. The IEEE 1547 series of interconnection
standards specifies the working interface between distributed
energy resources and the larger grid.

\engterm{Grid-forming inverters} are a key enabling technology for
high-renewable microgrids and for the large-scale grid. Unlike
\engterm{grid-following inverters} (the standard until recently),
which require a synchronous reference from the grid to operate,
grid-forming inverters generate the voltage and frequency reference
themselves, enabling 100$\,\%$ inverter-based operation of an
electrically isolated system. The Hawaii inverter-based-resource
demonstration (2022) was the first utility-scale demonstration of
this capability at a state-grid scale; rollout is ongoing in
high-renewable jurisdictions.

\engterm{Black start}, the restoration of a grid from a complete shutdown, is
the capability an isolated system most needs and the bulk grid most takes for
granted. After a total blackout the grid is dead: there is no voltage to
synchronise to, no frequency reference, and most generators cannot start
themselves because their auxiliaries (pumps, fans, control power) need
electricity the dead grid cannot supply. Restoration proceeds outward from a
few \engterm{black-start units} that can start from nothing, historically a
hydroelectric plant (a gate opens and water spins the turbine with no external
power) or a diesel or gas unit with its own starting battery and compressed
air. A black-start unit energises a transmission path to a larger plant, whose
auxiliaries then start on that path's power, and the restored capacity fans out
to bring more plant and load online in a controlled sequence, matching
generation to load at every step so the fragile early island does not trip on
the imbalance a large block of load would cause. The sequence is slow and
delicate, hours for a large grid, because each newly energised segment is a
small, low-inertia island vulnerable to exactly the frequency and voltage
swings the off-grid discussion described, so load is picked up in small
increments and the operators watch the frequency after each. The transition
complicates black start the way it complicates inertia and reactive support:
the retiring synchronous units were the black-start providers, and a
high-inverter grid must certify that its grid-forming inverters and storage can
start dead and build an island from nothing, a capability that grid-following
inverters, which need an existing reference, cannot supply. A microgrid earns
its resilience precisely by being able to black-start itself, which is why the
control architecture of the off-grid and microgrid systems above is written
around a grid-forming source that can raise the first voltage with no
neighbour to lean on.

\begin{mastery}
\textbf{Where the inertia goes.} A synchronous grid stores kinetic
energy in the spinning mass of every generator and motor, and that
stored energy is what slows the rate of frequency change after a trip.
The inertia constant $H$ measures it: the seconds of rated output the
stored kinetic energy could supply, typically $H \approx 2$ to $9\,\si{
\second}$ for a thermal fleet. The initial rate of change of frequency
after losing a block $\Delta P$ of generation is $df/dt = -f_0\,\Delta
P/(2H S)$ with $S$ the system rating, so halving the inertia doubles the
initial slope and deepens the nadir (Figure
\ref{fig:vol04ch14:frequency-response}). Replacing synchronous
generators with inverter-coupled wind and solar removes their physical
inertia, because the inverter decouples the source from the grid
frequency. A high-renewable grid must therefore manufacture the
stabilising response it used to get for free: synthetic (virtual)
inertia and fast frequency response programmed into grid-forming
inverters, faster primary control, and a cap on the size of the largest
single in-feed. The integration question for the energy transition is
not only how to generate the energy but how to keep the whole machine
stable once the spinning mass that used to stabilise it is gone.
\end{mastery}

\section{The energy transition}\pathtag{standard}

The defining infrastructure question of the early twenty-first
century is whether and how the world can replace fossil-fuel
generation with low-carbon generation in the time scale dictated by
the climate constraints of the Paris Agreement and successor
processes. The chapter does not advocate; it states the operational
content.

Global electricity generation, current as of 2024, is approximately
$30\,000\,\text{TWh}/\text{yr}$. The mix is approximately
$60\,\%$ fossil fuels (coal, gas, oil), $10\,\%$ nuclear, $15\,\%$
hydro, $7\,\%$ wind, $5\,\%$ solar, $3\,\%$ other (biomass,
geothermal, marine), according to Ember and IEA reporting current as
of 2024. Electricity is approximately $20\,\%$ of total final energy
consumption; the remainder is direct fuel use (transport, heating,
industrial heat).

% Figure: simplified Sankey-style diagram of national energy flow from
% primary sources through carriers to end uses, with rejected (waste)
% energy as a separate sink. Built from plain TikZ boxes and bands so it
% compiles without the sankey package. Widths are schematic.
\begin{figure}[htbp]
\centering
\begin{tikzpicture}[
  font=\scriptsize,
  src/.style={draw=engblue, fill=engblue!12, rounded corners=2pt,
    minimum width=2.5cm, minimum height=0.7cm, align=center},
  car/.style={draw=engamber, fill=engamber!15, rounded corners=2pt,
    minimum width=2.5cm, minimum height=0.7cm, align=center},
  use/.style={draw=englime!60!black, fill=englime!18, rounded corners=2pt,
    minimum width=2.5cm, minimum height=0.7cm, align=center},
  loss/.style={draw=engred, fill=engred!12, rounded corners=2pt,
    minimum width=2.5cm, minimum height=0.7cm, align=center},
  flow/.style={-{Stealth}, enggray, thick},
]
% primary sources (left column)
\node[src] (fossil) at (0,4)   {fossil fuels};
\node[src] (nuclear) at (0,2.6){nuclear};
\node[src] (renew) at (0,1.2)  {renewables\\ + hydro};

% carriers (middle column)
\node[car] (elec) at (5,3.3)  {electricity\\ ($\sim$20\% of final)};
\node[car] (direct) at (5,1.2){direct fuel\\ use};

% end uses (right column)
\node[use] (uses) at (10,3.3) {end uses:\\ transport, heat,\\ industry, light};
% rejected energy
\node[loss] (rej) at (10,1.0) {rejected energy\\ (conversion +\\ transmission losses)};

% flows: sources -> carriers
\draw[flow] (fossil.east) -- (elec.west);
\draw[flow] (nuclear.east) -- (elec.west);
\draw[flow] (renew.east) -- (elec.west);
\draw[flow] (fossil.east) -- (direct.west);
\draw[flow] (renew.east) -- (direct.west);
% carriers -> uses and losses
\draw[flow] (elec.east) -- (uses.west);
\draw[flow] (direct.east) -- (uses.west);
\draw[flow, engred] (elec.east) -- (rej.west);
\draw[flow, engred] (direct.east) -- (rej.west);

% column labels
\node[engblue, anchor=south] at (0,4.7) {primary sources};
\node[engamber, anchor=south] at (5,4.7) {energy carriers};
\node[englime!60!black, anchor=south] at (10,4.7) {end uses / losses};
\end{tikzpicture}
\caption[Simplified national energy-flow diagram]{A schematic
energy-flow (Sankey-style) diagram for a national system. Primary
sources feed two carriers, electricity and direct fuel; the carriers
feed end uses, with a large fraction lost as rejected energy in
conversion and transmission. Electricity is roughly one fifth of final
energy consumption (current as of 2024); electrification of transport
and heat shifts flow from the direct-fuel band to the electricity band.
Band positions are schematic, not to scale.}
\label{fig:vol04ch14:energy-sankey}
\end{figure}


Figure \ref{fig:vol04ch14:energy-sankey} sets out the whole-system
balance the transition must rework. Primary sources feed two carriers,
electricity and direct fuel; the carriers feed end uses, with a large
fraction lost as rejected energy in conversion and transmission. The
transition shifts flow from the direct-fuel band into the electricity
band (electrification of transport and heat) while decarbonising the
electricity band itself, and it shrinks the rejected-energy sink
because electric drives and heat pumps reject far less than the
combustion processes they replace.

The engineering content of the transition is fivefold:
\begin{itemize}[leftmargin=*]
  \item Electrification: more end uses are electrified (passenger
    transport via EVs, heating via heat pumps, light industrial
    process heat via electric furnaces). Electrification multiplies
    electricity demand even as it reduces total final energy demand
    (electric drives are more efficient than internal combustion;
    heat pumps deliver more thermal energy than they consume in
    electrical energy).
  \item Decarbonisation of electricity: the existing $20\,\%$ of
    final energy that is electricity is shifted from fossil to low-
    carbon sources. The pace of solar-PV cost decline ($\sim 90\,\%$
    over 2010-2024) makes this rapidly affordable in many regions.
  \item Decarbonisation of heat: gas-fired space heating gives way to
    electric heat pumps; district heating networks integrate waste
    heat from industry, data centres, and combined heat-and-power
    plants. Industrial heat above $\sim 400\,\si{\celsius}$ is
    harder; hydrogen, electric arc, and concentrated solar
    technologies are in early commercial deployment.
  \item Decarbonisation of transport: passenger cars and light-duty
    trucks transition to battery-electric drivetrains. Heavy
    transport (long-haul trucking, aviation, shipping) requires
    higher energy density than current batteries; hydrogen, ammonia,
    sustainable aviation fuels, and battery-electric for short-haul
    are the candidate solutions.
  \item Decarbonisation of industry: chemical-industry feedstocks
    (ammonia and steel in particular) are produced from natural gas and
    coke today; low-carbon hydrogen and direct-reduced-iron processes
    are the engineered alternatives.
\end{itemize}

\begin{mastery}
\textbf{Energy return on investment: an energy source must repay its own
energy first.} Before a generator delivers net energy to society it must
return the energy that built it, and the ratio of the two is the
\engterm{energy return on energy invested} (EROI): the lifetime energy a
source delivers divided by the energy spent to build, fuel, operate, and
retire it. A source with an EROI of $10{:}1$ returns ten units of energy for
every one it consumed, so a tenth of its output is, in effect, reinvested to
build the next generator, and the other nine tenths are the net energy the
rest of society runs on. The quantity sets a floor no cost accounting sees: a
source with an EROI near $1{:}1$ is an energy sink whatever its price, because
it barely returns the energy it took, and a modern industrial society is
generally reckoned to need a system-wide EROI well above $\sim 5{:}1$ to
$10{:}1$ to run its non-energy sectors. Historic conventional oil ran at a
high EROI, tens to one at the wellhead, which is part of why cheap fossil
energy powered the twentieth century; the figure has fallen as the easy
reserves were used and extraction moved to deeper, tighter, and more remote
resources. Modern wind sits at a healthy EROI, on the order of $\sim 20{:}1$
over its life, and utility solar somewhat lower but rising as manufacturing
energy falls and panel lifetimes lengthen. The subtlety the transition forces
is that storage and overbuild lower the delivered EROI of a variable source,
because the energy to build the battery and the curtailed energy of the
overbuild both count against the net delivered, so the honest EROI of a
high-renewable system is the EROI of the generation reduced by the storage and
curtailment it needs to be firm. The reading pairs with the capacity-credit
and value-factor arguments: a source must clear three separate bars to carry a
grid, an economic one (its value-adjusted cost), a reliability one (its
capacity credit), and a thermodynamic one (its EROI net of the firming it
needs), and a source can pass one and fail another.
\end{mastery}

\begin{mastery}
\textbf{Capacity credit: firm and variable capacity are not the same
currency.} A grid planner must meet peak demand with a margin, and the
question is how much of a generator's nameplate counts toward that
requirement. The measure is the \engterm{capacity credit} (or capacity
value): the firm capacity a resource can be relied on to deliver at the
moments the system is most stressed, expressed as a fraction of its
nameplate. A thermal plant available on demand has a capacity credit near
its forced-outage-adjusted nameplate, $\sim 90\,\%$. A solar array has a
capacity credit near zero on a winter-evening peak (the sun is down when a
cold northern grid peaks) and high on a summer-afternoon air-conditioning
peak (the sun is up); the same array is worth firm capacity in one system
and almost none in another. Wind sits in between and depends on the
correlation between wind and the peak. The capacity credit also falls as
penetration rises: the first gigawatt of solar on a summer-peaking grid
earns a high credit, but once midday is saturated the marginal solar adds
energy without adding firm capacity, and its capacity credit collapses
toward zero. This is why an energy LCOE comparison understates what a
dispatchable plant provides and overstates what an unpaired renewable
provides: they are priced in different currencies, energy and firm
capacity, and the exchange rate between them is the capacity credit, which
is system-specific and penetration-dependent. The planning consequence is
that a high-renewable grid must procure firm capacity separately, through
storage with enough duration to cover the peak, demand response, firm
low-carbon generation, or interconnection to a system whose peak does not
coincide. Sector coupling is the other side of the same coin: the cheap
surplus of an overbuilt renewable system (the curtailed energy of the
overbuild calculation) is worth little as electricity but valuable as a
feedstock for power-to-X conversion, electrolytic hydrogen, synthetic
fuels, thermal storage, and flexible EV charging, which absorb the surplus
and turn the curtailment problem into a supply of low-cost energy for the
hard-to-electrify sectors. The integration question of the transition is
to value energy, firm capacity, and flexibility as the three distinct
products they are.
\end{mastery}

\begin{mastery}
\textbf{Transmission is the cheap decarbonisation lever, and the slow one.}
The variability that storage buffers in time, transmission buffers in
space. Wind and solar are partly uncorrelated across a continent: a
weather system that becalms one region energises another, and the noon of
a western time zone is the evening peak of an eastern one. Averaging
generation over a wide synchronous area therefore smooths the aggregate
output and raises the capacity credit of the renewable fleet without
storing a single kilowatt-hour. The smoothing follows the statistics of
summing weakly correlated sources: the relative variability of the sum
falls roughly as $1/\sqrt{N_{\text{eff}}}$ in the number of effectively
independent regions, so the first few interconnections buy the most
smoothing and later ones less. A megawatt-kilometre of high-voltage
transmission is among the cheapest tools the transition has, often cheaper
per unit of firmed energy than the equivalent storage, because a line moves
energy that already exists rather than building a reservoir to hold it.
The constraint is not cost or physics but time and consent: a long
transmission corridor crosses many jurisdictions and many landowners, and
the permitting and right-of-way acquisition routinely take longer than the
generation it would connect. The build rate of transmission, not the cost
of panels or turbines, is the binding constraint on the pace of integration
in several large grids (current as of 2025). High-voltage direct current
is the technology of choice for the longest corridors, because it carries
bulk power over thousands of kilometres with lower loss than AC and ties
asynchronous grids together without synchronising them, the role the
missing interconnection played in the Texas 2021 event
\cite{acc:ferc-nerc-texas-2021}.
\end{mastery}

\begin{mastery}
\textbf{The value factor: a renewable eats its own price.} The levelized
cost prices a generator's energy at a constant per-unit figure, but the
market pays the hourly clearing price, and a variable renewable does not
collect the average price, because it produces most when others like it are
also producing and the price is depressed. The measure is the
\engterm{value factor}, the ratio of the average price a generator actually
captures to the time-averaged wholesale price. A dispatchable plant that can
choose its hours captures a value factor above one; a solar array, whose
output clusters in the sunny midday hours that the merit-order effect drives
toward the cheap left of the stack, captures a value factor below one, and
the factor falls as solar penetration rises, because each added array
deepens the midday price trough it must sell into. At low penetration solar
captures close to the average price; by the time midday solar routinely sets
a near-zero clearing price the marginal array captures a small fraction of
the average, so its market value collapses even as its levelized cost keeps
falling. The consequence sharpens the capacity-credit argument: a renewable
has two declining currencies at high penetration, its capacity credit (the
firm power it adds toward the peak) and its value factor (the price its
energy captures), and both fall for the same reason, that the resource
correlates with itself across the fleet. The levelized cost, which knows
nothing of when the energy lands, is a fair screen only at low penetration;
past that, the comparison that decides what to build is the value-adjusted
cost, the levelized cost divided by the value factor, and it is what makes
storage, curtailment-absorbing flexible loads, and a diversified wind-and-
solar-and-firm portfolio worth more than the cheapest single technology.
Storage raises a renewable's value factor by moving its energy from the
depressed midday into the valuable evening, which is one more reading of why
storage and variable renewables are complements rather than competitors:
the storage buys back the value the renewable's own abundance destroyed.
\end{mastery}

\begin{mastery}
\textbf{The abatement-cost curve orders the measures, and the order is not the
price order.} A region cutting its emissions has many measures available, and
the discipline that decides which to do first is not the cheapest capital or
the largest single cut but the \engterm{marginal abatement cost}, the net cost
per tonne of carbon dioxide a measure removes, worked for one measure in the
abatement-cost example. Ranking every measure by that number, from the most
negative to the most positive, and plotting each as a bar whose width is the
annual tonnes it abates and whose height is its cost per tonne, builds the
\engterm{marginal-abatement-cost curve}, a staircase that reads left to right
from the cheapest carbon to the dearest. The curve's first surprise is that its
left end sits below zero: a run of measures, chiefly efficiency, insulation,
heat pumps, and managed charging, cut carbon while saving money, because the
energy they avoid is worth more than their extra capital, so they should be
done for their own sake before any carbon price is considered. The second is
that the curve rises steeply at its right end, where the last tonnes are
removed by measures whose cost per tonne is high, the deep industrial
decarbonisation, the synthetic fuels, the direct air capture, so the marginal
cost of the last emissions cut is many times the cost of the first. A carbon
price works by drawing a horizontal line across the curve: every measure whose
bar sits below the price is worth doing at that price, and raising the price
walks the region rightward along the staircase into the more expensive
measures. The curve reframes the transition as a sequencing problem, do the
negative-cost measures now, then climb the staircase as the carbon price rises
or the technologies cheapen, and it explains why the cheap early progress of
the decarbonisation pathway gives way to a rising cost as the last emissions
are approached: the pathway is walking up its own abatement-cost curve. The
limits are that the curve is a snapshot, its bars move as technologies learn
down their cost curves (the learning-curve box), its widths interact because
one measure changes another's baseline, and it prices only the direct cost, not
the co-benefits or the system-integration costs the rest of this chapter
insists on, so it is a screen for ordering measures, not the whole plan.
\end{mastery}

The pace of the transition is constrained by the economics of new
generation, the speed of grid build-out, the supply chains for
batteries and rare-earth magnets, the workforce for installation and
maintenance, and the policy framework. The book takes no normative
position on whether the transition is occurring fast enough; it states
that the engineering of each component is well understood and the
integration questions are the working open problems.

\section{Worked examples: sizing the system}\pathtag{core}

The chapter's quantities combine into a handful of recurring
calculations. We work three that recur in practice: the capacity factor
to annual energy bridge, the storage energy to ride through a
shortfall, and the levelized cost comparison that decides what to
build.

\paragraph{From capacity factor to annual energy.} A $1\,\si{\giga
\watt}$ nameplate wind farm with a capacity factor of $35\,\%$ produces
\[
E = P_{\text{rated}}\,\text{CF}\,(8760\,\si{\hour})
  = (1000\,\si{\mega\watt})(0.35)(8760\,\si{\hour})
  \approx 3.07\,\text{TWh}/\text{yr}.
\]
A $1\,\si{\giga\watt}$ nuclear plant at a $92\,\%$ capacity factor
produces $\approx 8.06\,\text{TWh}/\text{yr}$, so matching the nuclear
plant's energy with wind needs $8.06/3.07 \approx 2.6\,\si{\giga\watt}$
of wind nameplate. Nameplate capacity and delivered energy are not the
same currency, and the capacity factor is the exchange rate; the
\texttt{code/capacity\_factor.py} routine carries the conversion both
ways.

\paragraph{Storage to ride through a shortfall.} Suppose a grid with a
$45\,\si{\giga\watt}$ peak demand and an average load of $\sim 70\,\%$
of peak must ride through a ten-day still-and-cloudy event entirely on
storage. The energy is
\[
E = (0.70)(45\,\si{\giga\watt})(24\,\si{\hour})(10) \approx 7.6\,
\text{TWh},
\]
which at a lithium-ion energy cost of $\$140\,\si{\per\kilo\watt\hour}$
(current as of 2024) costs $\sim \$1.06$ trillion. The number is the
quickest argument that lithium-ion is the wrong tool for multi-day
storage at national scale; long-duration chemistries with low
energy-cost (pumped hydro, hydrogen) or firm low-carbon generation
carry that role. The arithmetic is in \texttt{code/storage\_sizing.py}.

\paragraph{Levelized cost comparison.} The levelized cost of energy
discounts a plant's lifetime cost and lifetime energy to a single
per-unit price. For utility solar with a capital cost of $\sim \$900\,
\si{\per\kilo\watt}$, a fixed running cost of $\sim \$15\,\si{\per\kilo
\watt}$ per year, no fuel, an $18\,\%$ capacity factor, a 30-year life,
and a $7\,\%$ discount rate, the calculation in \texttt{code/lcoe.py}
returns $\sim \$56\,\si{\per\mega\watt\hour}$, in the range of Figure
\ref{fig:vol04ch14:lcoe-bars}. The levelized cost is a first-order
screen, not the whole decision: it omits the system cost of
intermittency, which is where the storage and dispatch of the rest of
this chapter re-enter.

\paragraph{Reliability: loss-of-load expectation versus autonomy.} The
levelized-cost screen prices energy; a stand-alone or storage-backed
system must also be priced for reliability, and the currency there is the
\engterm{loss-of-load expectation} (LOLE), the expected hours per year in
which load exceeds the available supply. Take a stand-alone
solar-plus-storage system carrying an average load $L$, with a battery of
usable energy $E_b = n L \cdot 24\,\si{\hour}$ for $n$ days of autonomy.
Reliability does not follow $E_b$ linearly. Suppose the daily net energy
(generation minus load) is a random draw with a long negative tail set by
multi-day low-sun weather; the system fails only when a run of poor days
drains the battery below empty, so the probability of a deficit-day-run
exceeding the stored energy falls steeply at first and then flattens into a
tail governed by the rare long events. Modelling the daily deficit run as
approximately geometric, the unserved hours scale as $\text{LOLE}(n)
\approx H_0\, p^{\,n}$ for a base rate $H_0$ and a per-day survival
factor $p < 1$, so each added day of autonomy multiplies the unserved
hours by $p$ rather than removing a fixed amount. With $H_0 \approx
1500\,\si{\hour}/\text{yr}$ and $p \approx 0.35$ at a PV overbuild of
$1.3\times$ (the values behind Figure
\ref{fig:vol04ch14:lolp-autonomy}), reaching a $99.95\,\%$ reliability
target (LOLE $\approx 4.4\,\si{\hour}/\text{yr}$, since $0.0005\times8760
\approx 4.4$) needs $n = \ln(4.4/1500)/\ln(0.35) \approx 5.6$, so about
$2.8\,$ days of autonomy. Pushing the same target on an array with no
overbuild ($p \approx 0.62$) needs $n = \ln(4.4/1800)/\ln(0.62) \approx
12.6$, more than four days of storage. The reliability of a renewable
system is bought jointly by overbuilding generation and adding storage,
and the two trade against each other along the curve of Figure
\ref{fig:vol04ch14:lolp-autonomy}; the last fraction of a percent is the
most expensive, which is the recurring argument for a small dispatchable
backup rather than a pure-renewable build to $100\,\%$.

% Figure: loss-of-load expectation (hours per year of unserved energy)
% against battery autonomy for a stand-alone solar-plus-storage system,
% shown for two PV overbuild ratios. The curve falls steeply then flattens
% into a long reliability tail, the structure that makes the last fraction
% of a percent of reliability expensive. Values illustrative.
\begin{figure}[htbp]
\centering
\begin{tikzpicture}
\begin{axis}[
  width=12.5cm, height=7.5cm,
  xlabel={battery autonomy (days of average load)},
  ylabel={loss-of-load expectation (h/yr)},
  xmin=0, xmax=5, ymin=0.5, ymax=2000,
  ymode=log,
  log basis y=10,
  xtick={0,1,2,3,4,5},
  ytick={1,10,100,1000},
  yticklabels={1,10,100,1000},
  grid=both,
  grid style={enggray!25},
  tick label style={font=\scriptsize},
  label style={font=\small},
  legend style={font=\scriptsize, at={(0.97,0.97)}, anchor=north east, draw=none},
]
% PV oversized 1.3x: faster decay, lower tail
\addplot[engblue, very thick, mark=*, mark size=1.3pt] coordinates {
  (0.25,1500)(0.5,700)(1,180)(1.5,55)(2,18)(2.5,7)(3,3.2)(3.5,1.8)(4,1.2)(5,0.8)
};
\addlegendentry{PV overbuild $1.3\times$}
% PV oversized 1.0x: higher tail, slower decay
\addplot[engamber, very thick, dashed, mark=square*, mark size=1.3pt] coordinates {
  (0.25,1800)(0.5,1100)(1,400)(1.5,160)(2,75)(2.5,42)(3,28)(3.5,21)(4,17)(5,13)
};
\addlegendentry{PV overbuild $1.0\times$}
% reliability target line (~4.4 h/yr ~ 99.95%)
\draw[engred, thick, dash dot] (axis cs:0,4.4) -- (axis cs:5,4.4);
\node[engred, font=\scriptsize, anchor=south west] at (axis cs:0.1,4.4) {$99.95\,\%$ target ($\approx 4.4$ h/yr)};
\end{axis}
\end{tikzpicture}
\caption[Loss-of-load expectation versus autonomy]{Loss-of-load expectation
against battery autonomy for a stand-alone solar-plus-storage system, at two
PV overbuild ratios (log vertical scale). Each added day of storage buys a
smaller reduction in unserved energy: the curve drops steeply then flattens
into a long tail set by rare multi-day low-sun events. Reaching a fixed
reliability target (red line) costs far more storage on the under-built
array than on the oversized one, which is why overbuilding generation and
sizing storage are coupled decisions rather than independent ones. Values
illustrative of a mid-latitude resource.}
\label{fig:vol04ch14:lolp-autonomy}
\end{figure}


\paragraph{Curtailment of overbuilt renewables.} Overbuilding generation
to cover the low-resource periods spills energy in the high-resource
periods, and that spilled energy is paid for and not delivered. Define
the \engterm{overbuild ratio} $r$ as annual renewable generation divided
by annual load. At $r = 1$ a system with perfect storage and no losses
would exactly balance, but with finite storage the surplus at times of
high output cannot all be stored, and the \engterm{curtailed fraction}
rises with $r$. A useful closed form treats the hourly net generation as
a distribution: the curtailed energy is the part of the surplus that
exceeds what the storage can absorb. For a system with $4\,\si{\hour}$ of
storage, raising the overbuild from $r = 1.5$ to $r = 2.0$ lifts annual
generation by a third but raises the curtailed fraction from $\sim
15\,\%$ to $\sim 30\,\%$ (Figure
\ref{fig:vol04ch14:overbuild-curtailment}), so the delivered energy grows
much more slowly than the nameplate. Adding storage (the $12\,\si{\hour}$
curve) absorbs more of the surplus and shifts the curtailment down at
every $r$. The engineering reading is that overbuild and storage are
substitutes for covering the low periods, but overbuild pays a rising
curtailment penalty while storage pays a rising capital penalty, and the
least-cost system sits where the two marginal costs meet. The curtailed
energy is not waste in a useless sense: at high penetration it is the
cheap surplus that powers flexible loads (electrolysis, thermal storage,
EV charging) which is the sector-coupling argument the next mastery box
develops.

% Figure: curtailed energy fraction against renewable overbuild ratio for a
% renewables-plus-storage grid. As generation is overbuilt to cover low
% periods, an increasing share of the annual energy is spilled at times of
% surplus. Two storage levels shift the curve. Values illustrative.
\begin{figure}[htbp]
\centering
\begin{tikzpicture}
\begin{axis}[
  width=12.5cm, height=7.5cm,
  xlabel={renewable overbuild ratio (generation / average load)},
  ylabel={curtailed fraction of generation (\%)},
  xmin=1, xmax=3, ymin=0, ymax=55,
  xtick={1,1.5,2,2.5,3},
  ytick={0,10,20,30,40,50},
  grid=both,
  grid style={enggray!25},
  tick label style={font=\scriptsize},
  label style={font=\small},
  legend style={font=\scriptsize, at={(0.03,0.97)}, anchor=north west, draw=none},
]
% low storage (4 h): curtailment rises fast
\addplot[engorange, very thick, mark=*, mark size=1.3pt] coordinates {
  (1.0,1)(1.25,7)(1.5,15)(1.75,23)(2.0,30)(2.25,36)(2.5,41)(2.75,45)(3.0,48)
};
\addlegendentry{storage $4$ h}
% higher storage (12 h): curtailment rises slower
\addplot[engblue, very thick, dashed, mark=square*, mark size=1.3pt] coordinates {
  (1.0,0.3)(1.25,3)(1.5,7)(1.75,12)(2.0,17)(2.25,22)(2.5,27)(2.75,31)(3.0,35)
};
\addlegendentry{storage $12$ h}
\end{axis}
\end{tikzpicture}
\caption[Curtailment versus renewable overbuild]{Curtailed share of annual
renewable generation against the overbuild ratio, for two storage durations.
Covering the low-resource periods by overbuilding generation spills a growing
fraction of the energy at times of surplus; more storage absorbs some of that
surplus and shifts the curve down. The curtailed energy is real cost: it is
generation paid for and not delivered, and it sets a practical ceiling on how
far overbuilding alone can substitute for storage or firm capacity. Values
illustrative of a wind-and-solar portfolio.}
\label{fig:vol04ch14:overbuild-curtailment}
\end{figure}


\paragraph{Pumped hydro versus batteries: where each wins.} The mastery
box on power-energy pricing asserted a crossover between lithium-ion and
long-duration storage; the arithmetic fixes where it sits. Model each
technology's installed cost as a fixed power term plus a per-hour energy
term, $C(d) = c_P + c_E\, d$ in dollars per kilowatt of power for a
discharge duration $d$ in hours. Lithium-ion (current as of 2024) carries
roughly $c_P \approx \$120\,\si{\per\kilo\watt}$ for the inverter and
balance-of-system and $c_E \approx \$140\,\si{\per\kilo\watt\hour}$ for
the cells, so a $4\,\si{\hour}$ system costs $120 + 140\times4 = \$680\,
\si{\per\kilo\watt}$ and an $8\,\si{\hour}$ system $\$1240\,\si{\per\kilo
\watt}$. Pumped hydro carries a high fixed term, $c_P \approx \$1200\,
\si{\per\kilo\watt}$ for the turbines, civil works, and powerhouse, and a
low energy term, $c_E \approx \$15\,\si{\per\kilo\watt\hour}$ for the
reservoir volume, so a $4\,\si{\hour}$ plant costs $1200 + 15\times4 =
\$1260\,\si{\per\kilo\watt}$ and an $8\,\si{\hour}$ plant $\$1320\,\si{
\per\kilo\watt}$. The crossover duration is where the two lines meet,
$c_{P,\text{Li}} + c_{E,\text{Li}}\,d = c_{P,\text{PH}} +
c_{E,\text{PH}}\,d$, giving $d^\ast = (c_{P,\text{PH}} -
c_{P,\text{Li}})/(c_{E,\text{Li}} - c_{E,\text{PH}}) = (1200 - 120)/(140
- 15) \approx 8.6\,\si{\hour}$. Below about nine hours of discharge,
lithium-ion is cheaper per kilowatt of firm power; above it, pumped hydro
wins, and the gap widens with every added hour because the energy term
dominates at long duration. The crossover is approximate and moves with
cell prices, which have been falling, but the structure is durable: the
storage market splits by duration, with lithium-ion holding the
minutes-to-hours band and reservoir-or-fuel storage holding the
hours-to-seasons band. Geography is the override: pumped hydro needs two
reservoirs at height, and where the terrain does not provide them the
crossover is moot.

\paragraph{Merit-order dispatch and the marginal price.} A wholesale
electricity market dispatches generators in order of rising marginal cost
until supply meets demand, and the marginal cost of the last unit
dispatched sets the price every generator is paid. Take an illustrative
grid with five generator blocks: $12\,\si{\giga\watt}$ of wind and solar at
a marginal cost of $\$0\,\si{\per\mega\watt\hour}$ (no fuel), $8\,\si{\giga
\watt}$ of nuclear at $\$8$, $10\,\si{\giga\watt}$ of efficient
combined-cycle gas at $\$35$, $10\,\si{\giga\watt}$ of older coal at $\$70$,
and $8\,\si{\giga\watt}$ of gas peakers at $\$140$. Stack them left to
right by cost, the supply curve of Figure
\ref{fig:vol04ch14:merit-order}. At a demand of $34\,\si{\giga\watt}$ the
dispatch fills the wind and solar ($12$), the nuclear (to $20$), the
combined-cycle gas (to $30$), and $4\,\si{\giga\watt}$ of the coal block,
so the marginal unit is coal at $\$70\,\si{\per\mega\watt\hour}$ and the
clearing price is $\$70$. Every dispatched generator earns that price,
including the zero-cost wind and solar, whose difference between the
clearing price and their own marginal cost is the inframarginal rent that
pays back their capital. The price is not the average cost of generation;
it is the cost of the marginal megawatt-hour. Two consequences follow.
First, adding zero-marginal-cost renewables pushes the whole stack to the
right and lowers the clearing price for a given demand, the
\engterm{merit-order effect} that depresses wholesale prices in
high-renewable hours. Second, when demand climbs into the steep right of
the stack the price jumps from $\$70$ to $\$140$ for a small increase in
load, the scarcity-pricing regime that rewards flexibility and makes the
case for demand response. The same dispatch logic, run hour by hour with a
storage state variable, is the national-scale extension of the household
loop in \texttt{code/household\_dispatch.py}.

% Figure: merit-order supply curve (dispatch stack as a step function of
% cumulative capacity versus marginal cost) crossed by a vertical demand
% line. The intersection sets the system marginal price. Step plot built
% with const plot. Values illustrative, current as of 2025.
\begin{figure}[htbp]
\centering
\begin{tikzpicture}
\begin{axis}[
  width=13cm, height=8cm,
  xlabel={cumulative available capacity (GW)},
  ylabel={marginal cost (USD/MWh)},
  xmin=0, xmax=50, ymin=0, ymax=200,
  xtick={0,10,20,30,40,50},
  ytick={0,40,80,120,160,200},
  axis lines=left,
  tick label style={font=\scriptsize},
  label style={font=\small},
  legend style={font=\scriptsize, at={(0.03,0.97)}, anchor=north west, draw=none},
]
% merit-order step curve: cheap renewables and nuclear first, gas last
\addplot[engblue, very thick, const plot] coordinates {
  (0,0)(12,0)(12,8)(20,8)(20,35)(30,35)(30,70)(40,70)(40,140)(48,140)(48,200)(50,200)
};
\addlegendentry{merit-order supply}
% demand line at 34 GW
\addplot[engred, very thick, dashed] coordinates {
  (34,0)(34,200)
};
\addlegendentry{demand $= 34$ GW}
% marker at the clearing point
\addplot[engorange, only marks, mark=*, mark size=2pt] coordinates {(34,70)};
\node[anchor=west, font=\scriptsize, color=engorange]
  at (axis cs:34.5,82) {clearing price $\approx 70$};
\end{axis}
\end{tikzpicture}
\caption[Merit-order supply curve and marginal price]{The merit-order
supply curve for an illustrative grid: generators stacked left to right in
order of rising marginal cost, from near-zero-cost wind and solar through
nuclear and coal to flexible gas (current as of 2025). The vertical demand
line cuts the stack at the marginal generator, and its marginal cost sets
the system clearing price that every dispatched generator is paid. As
demand rises into the steep right of the stack, the price climbs sharply,
which is the mechanism behind scarcity pricing and the case for demand
response. Quantities illustrative.}
\label{fig:vol04ch14:merit-order}
\end{figure}


\paragraph{Demand response: peak shaving and the avoided peaker.} The
steep right of the merit-order stack is also the expensive end of the
capacity build, because the peaker capacity that serves the last few hours
of the year runs at a tiny capacity factor and earns its keep only in those
hours. Demand response attacks the peak directly. Suppose a utility faces a
summer peak of $40\,\si{\giga\watt}$ that exceeds its firm dispatchable
capacity of $38\,\si{\giga\watt}$ for the $\sim 80$ highest-demand hours of
the year, and it can either build $2\,\si{\giga\watt}$ of new peaker
capacity or enrol demand response. A $2\,\si{\giga\watt}$ peaker at an
installed cost of $\sim \$700\,\si{\per\kilo\watt}$ costs $\$1.4$ billion,
which at the $9.4\,\%$ capital-recovery factor of the earlier examples is
$\sim \$132\text{M}/\text{yr}$ to serve $\sim 2\,\si{\giga\watt}\times 80\,
\si{\hour} = 160\,\si{\giga\watt\hour}/\text{yr}$, an annualised capital
cost of $\sim \$0.82\,\si{\per\kilo\watt\hour}$ on the energy it actually
delivers, before fuel. Demand response instead pays consumers to shed
$2\,\si{\giga\watt}$ during those $80\,\si{\hour}$: at a typical
incentive of $\$1.0$ to $\$1.5\,\si{\per\kilo\watt\hour}$ shed, the cost is
$\sim \$160$ to $\$240\text{k}$ per event hour, or $\sim \$13$ to $\$19
\text{M}/\text{yr}$ for the full $80\,\si{\hour}$, roughly a fifth to a
seventh of the peaker's annualised capital and with no fuel, no emissions,
and no asset to maintain for the other $8680\,\si{\hour}$ of the year. The
arithmetic is why peak demand, not energy, is the quantity demand response
is paid to move: the marginal cost of the peak hours is set by capacity
that exists almost entirely to serve them. The limit is reliability of the
response itself; enrolled load that fails to shed when called is worse than
useless, so demand-response programmes are sized with a derating for
non-performance and paired with a dispatchable reserve.

\paragraph{Capacity factor from a Weibull wind distribution.} The capacity
factor of a wind farm is not a single number read off a site; it is an
integral of the turbine's power curve against the site's wind-speed
distribution, which is well described by a Weibull law
$f(v) = (k/\lambda)(v/\lambda)^{k-1}\exp[-(v/\lambda)^k]$ with scale
$\lambda$ (close to the mean speed) and shape $k$ (typically $\sim 2$,
the Rayleigh case). Take a site with a mean wind speed of $8\,\si{\meter
\per\second}$, so $\lambda \approx 9.0\,\si{\meter\per\second}$ at $k = 2$,
and a turbine that cuts in at $3\,\si{\meter\per\second}$, reaches rated
power at $12\,\si{\meter\per\second}$, holds rated power to a cut-out of
$25\,\si{\meter\per\second}$, and produces nothing outside that band. The
capacity factor is
\[
\text{CF} = \int_{v_{\text{in}}}^{v_{\text{rated}}}
  \Bigl(\tfrac{v}{v_{\text{rated}}}\Bigr)^{\!3} f(v)\,dv
  + \int_{v_{\text{rated}}}^{v_{\text{out}}} f(v)\,dv,
\]
where the first term integrates the cube-law ramp below rated power and the
second counts the fraction of time at rated power. The Weibull cumulative
$F(v) = 1 - \exp[-(v/\lambda)^k]$ gives the second term in closed form:
$F(25) - F(12) = \exp[-(12/9)^2] - \exp[-(25/9)^2] \approx e^{-1.78} -
e^{-7.72} \approx 0.169$. Evaluating the cube-weighted first term
numerically (the routine that extends \texttt{code/capacity\_factor.py}
to a wind-speed distribution) returns $\sim 0.17$, so the capacity factor
is $\sim 0.34$, in the range quoted for good onshore sites. The structure
of the answer carries the engineering: about half the energy comes from the
hours near rated power and about half from the cube-law ramp, the cut-in
and cut-out speeds matter little to the total (they bound regions of low
probability or low power), and the shape parameter $k$ matters as much as
the mean, because a gusty site (low $k$) spends more time both becalmed and
above rated than a steady site of the same mean speed. Raising the mean
speed from $8$ to $9\,\si{\meter\per\second}$ lifts the capacity factor to
$\sim 0.42$, the cube law again rewarding the windier site out of
proportion to the speed difference.

\paragraph{HVDC versus HVAC over distance.} The choice of alternating or
direct current for a long line turns on where the loss-and-cost lines
cross. Compare a $\pm 500\,\si{\kilo\volt}$ HVDC bipole and a $500\,\si{
\kilo\volt}$ HVAC line, each carrying $2\,\si{\giga\watt}$ over $1000\,\si{
\kilo\meter}$ with conductors of $0.012\,\si{\ohm\per\kilo\meter}$. The
HVDC line carries direct current with unity effective power factor and no
reactive component, so its current is $I = P/(2V) = 2\times 10^9/(2\times
5\times 10^5) = 2000\,\si{\ampere}$ per pole, and the resistive loss is
$P_{\text{loss}} = 2 I^2 R = 2(2000)^2(0.012\times 1000)\approx 192\,\si{
\mega\watt}$, about $9.6\%$. The HVAC line at power factor $0.95$ carries
$I = P/(\sqrt3 V\cos\phi) = 2\times 10^9/(\sqrt3\cdot 5\times 10^5\cdot
0.95)\approx 2431\,\si{\ampere}$, and its three-phase resistive loss is $3
I^2 R = 3(2431)^2(12)\approx 213\,\si{\mega\watt}$, before the added
reactive-charging current that a long AC cable draws and that limits its
reach without shunt compensation. The resistive losses are comparable here;
the decisive difference is the converter stations. HVDC pays a large fixed
cost for the rectifier and inverter terminals (of order $\$300\text{M}$ for
a $2\,\si{\giga\watt}$ bipole) but a lower cost per kilometre of line and no
reactive-compensation penalty, so the total-cost lines cross at a break-even
distance of $\sim 600$ to $800\,\si{\kilo\meter}$ overhead (and far shorter,
$\sim 50\,\si{\kilo\meter}$, for cable, because AC cable charging current
grows with length). Below the break-even AC wins on the cheaper terminals;
above it DC wins on the cheaper line and the absence of charging current,
which is why the longest corridors and all long submarine links are DC.

\paragraph{Distribution-transformer loss across the daily load.} A
transformer's loss has a fixed part and a load-dependent part, and the daily
load shape decides its efficiency. A $1\,\si{\mega\watt}$-rated distribution
transformer has a no-load (iron) loss of $1.5\,\si{\kilo\watt}$, present
whenever it is energised, and a full-load (copper) loss of $9\,\si{\kilo
\watt}$ that scales with the square of the load fraction $x$, so the loss at
load fraction $x$ is $P_{\text{loss}}(x) = 1.5 + 9x^2\,\si{\kilo\watt}$.
Over a day that sits at $x = 0.3$ for $16\,\si{\hour}$ and $x = 0.8$ for
$8\,\si{\hour}$, the energy loss is $[(1.5 + 9\cdot 0.09)\cdot 16 + (1.5 +
9\cdot 0.64)\cdot 8] = [2.31\cdot 16 + 7.26\cdot 8] = 37.0 + 58.1 = 95.1\,
\si{\kilo\watt\hour}/\text{day}$, against a delivered energy of $(0.3\cdot
16 + 0.8\cdot 8)\cdot 1000 = 11{,}200\,\si{\kilo\watt\hour}/\text{day}$, an
efficiency of $\sim 99.2\%$. The instructive part is where the loss sits:
the iron loss runs all day and dominates at the light night load, while the
copper loss dominates at the heavy evening load because of the $x^2$. The
peak efficiency falls at the load where iron and copper losses are equal,
$1.5 = 9x^2$ so $x \approx 0.41$, which is why utilities size distribution
transformers so their typical load sits near $40\%$ of rating rather than
near $100\%$.

\paragraph{Heat-pump seasonal energy from degree-days.} Electrifying a
building's heat replaces a furnace burning fuel at $\sim 90\%$ efficiency
with a heat pump delivering several units of heat per unit of electricity,
and the seasonal saving follows from the heating demand and the seasonal
coefficient of performance. A house with a design heat loss of $8\,\si{
\kilo\watt}$ at a $30\,\si{\kelvin}$ indoor-outdoor difference has a loss
coefficient $U_{\text{tot}} = 8000/30 \approx 267\,\si{\watt\per\kelvin}$.
Over a heating season of $2500$ heating degree-days (the integral of the
indoor-outdoor temperature difference over time, here in
$\si{\kelvin}\cdot\text{day}$), the heat demand is
\[
Q = U_{\text{tot}}\times \text{HDD} = 267\,\si{\watt\per\kelvin}\times 2500
\,\si{\kelvin}\cdot\text{day}\times 24\,\si{\hour}/\text{day}
\approx 16{,}000\,\si{\kilo\watt\hour}/\text{yr}.
\]
A gas furnace at $90\%$ efficiency burns $16{,}000/0.9 \approx 17{,}800\,\si{
\kilo\watt\hour}$ of fuel for that heat. A heat pump at a seasonal COP of
$3.0$ (the season-averaged value, below the design-point COP because the
coldest hours are the least efficient) draws $16{,}000/3.0 \approx 5300\,\si{
\kilo\watt\hour}$ of electricity, a factor of $\sim 3.4$ less delivered
energy. Whether that saves carbon depends on the grid: at a grid intensity
below $\sim 3\times$ the fuel's intensity per delivered kilowatt-hour the
heat pump wins on emissions, which holds on almost every modern grid and
improves as the grid decarbonises, and the same factor-of-three is why
electrifying heat raises electricity demand far less than the naive
fuel-energy swap would suggest.

\paragraph{Self-consumption saturation with battery size.} The household
project asks for the self-consumption fraction against battery size and
expects it to saturate; the shape follows from the daily surplus the
battery can catch. Take a household whose PV generates a daily surplus
(generation above the simultaneous load) of $S \approx 12\,\si{\kilo\watt
\hour}$ on a typical day, available midday, and an evening-and-night deficit
of comparable size. A battery of usable energy $E_b$ can shift at most
$\min(E_b, S)$ of that surplus from day to night. The self-consumption
fraction is the share of generation used on site, which rises from the
no-battery baseline (direct midday use only) toward a ceiling as $E_b$
grows, $f_{\text{sc}}(E_b) \approx f_0 + (f_{\max}-f_0)\,\min(E_b/S, 1)$ in
the simplest one-day model, so it climbs roughly linearly until $E_b$
reaches the daily surplus and then flattens hard. With $f_0 \approx 0.35$,
$f_{\max}\approx 0.80$, and $S \approx 12\,\si{\kilo\watt\hour}$, a $6\,\si{
\kilo\watt\hour}$ battery reaches $\sim 0.58$ and a $12\,\si{\kilo\watt
\hour}$ battery $\sim 0.80$; doubling again to $24\,\si{\kilo\watt\hour}$
adds almost nothing, because there is no further daily surplus left to
store. The saturation point is the daily surplus, not a fixed kilowatt-hour
figure, so a larger array (more surplus) pushes the knee to the right, and
the marginal kilowatt-hour of battery is worth buying only up to the daily
surplus it can actually catch. The remaining winter deficit is a seasonal
problem that no daily-cycling battery removes, which is the project's
expected finding that the household stays grid-dependent in winter unless
the PV is heavily oversized.

\paragraph{The DC-to-AC ratio and how much clipping is worth accepting.}
The microgrid sizing above chose a DC-to-AC ratio of $\sim 1.25$ without
justifying the number; the trade behind it is worth working. An array
oversized relative to its inverter overshoots the inverter's alternating-
current limit around midday, and the overshoot is clipped and spilled, but
the same larger array fattens the morning and evening shoulders where the
inverter runs below its limit. Model a clear-sky day as a half-sinusoid of
peak $R$ times the inverter rating over a twelve-hour window,
$p(t) = R\sin(\pi t/12)$ for $t$ in hours from sunrise, so the delivered AC
power is $\min(1, p(t))$ in units of the inverter rating. At $R = 1.25$ the
generation exceeds the cap between the two solutions of $1.25\sin(\pi t/12) =
1$, that is $\sin(\pi t/12) = 0.8$, giving $\pi t/12 = 0.927$ and $2.214$
radians, so $t$ from $3.54$ to $8.46\,\si{\hour}$, a clipping window of
$\sim 4.9\,\si{\hour}$ centred on noon. The clipped energy is the area of the
DC bell above the cap over that window; integrating $1.25\sin(\pi t/12) - 1$
from $3.54$ to $8.46\,\si{\hour}$ gives $\sim 0.94$ inverter-hours spilled,
against a delivered AC energy of $\sim 6.6$ inverter-hours over the day, so
the clipping loss is $\sim 12\%$ of what the array made at midday but only
$\sim 3\%$ of the whole day's generation, because the shoulders the oversize
captured are not clipped. Figure \ref{fig:vol04ch14:inverter-clipping} is the
picture. The oversize pays because the inverter is a large share of the
system cost and the array is cheap: paying for a smaller inverter and
accepting a few percent of clipped energy lowers the levelized cost, and the
optimum ratio rises toward $1.4$ for cloudy sites and winter-weighted loads,
where the clear-sky peaks that clip are rarer. The clipping loss is not
waste in the way a curtailed grid surplus is; it is a deliberate purchase of
a cheaper inverter and a flatter delivery profile.

% Figure: daily DC generation of an oversized array (DC-to-AC ratio > 1)
% against the inverter's AC clipping limit. The clear-sky DC bell exceeds
% the inverter cap around midday, so the flat-topped AC output clips the
% peak (shaded overshoot is spilled) while the fattened shoulders capture
% more morning and evening energy. Half-sinusoid DC model over the daylight
% window. Values illustrative.
\begin{figure}[htbp]
\centering
\begin{tikzpicture}
\begin{axis}[
  width=13cm, height=8cm,
  xlabel={hour of day},
  ylabel={power (fraction of AC rating)},
  xmin=5, xmax=19, ymin=0, ymax=1.4,
  xtick={5,7,9,11,13,15,17,19},
  ytick={0,0.2,0.4,0.6,0.8,1.0,1.2,1.4},
  axis lines=left,
  tick label style={font=\scriptsize},
  label style={font=\small},
  legend style={font=\scriptsize, at={(0.03,0.97)}, anchor=north west, draw=none},
]
% clear-sky DC bell, peak 1.25 (DC-to-AC ratio), zero outside 6..18
\addplot[engblue, thick, dashed, domain=6:18, samples=80]
  {1.25*sin(deg((x-6)/12*pi))};
\addlegendentry{DC generation (oversized array)}
% AC output = min(DC, 1.0): clipped at the inverter cap
\addplot[engorange, very thick, domain=6:18, samples=120]
  {min(1.0, 1.25*sin(deg((x-6)/12*pi)))};
\addlegendentry{AC output (clipped at inverter)}
% inverter cap reference line
\addplot[enggray, thin, dotted] coordinates {(5,1.0)(19,1.0)};
\node[anchor=south, font=\scriptsize, color=engred] at (axis cs:12,1.27)
  {clipped (spilled)};
\end{axis}
\end{tikzpicture}
\caption[Inverter clipping of an oversized array]{Daily output of a solar
array oversized relative to its inverter, a DC-to-AC ratio above one. The
clear-sky direct-current generation (dashed) rises above the inverter's
alternating-current limit around midday, so the delivered output (solid) is
flat-topped: the midday peak is clipped and the overshoot is spilled. The
oversize is bought deliberately, because the same larger array fattens the
morning and evening shoulders, capturing more energy in the many hours the
inverter is below its limit than it loses to the few hours of clipping. The
optimum DC-to-AC ratio trades the spilled midday peak against the fattened
shoulders and the cheaper inverter, and it rises for sites and seasons with
lower average irradiance. Quantities illustrative.}
\label{fig:vol04ch14:inverter-clipping}
\end{figure}


\paragraph{Seasonal heat-pump energy by the temperature-bin method.} The
degree-day estimate earlier collapsed the whole heating season into one
average coefficient of performance; a sharper estimate bins the hours by
outdoor temperature and applies the COP that holds in each bin, because the
COP falls exactly where the heat demand rises. Take the same house of loss
coefficient $U_{\text{tot}} = 267\,\si{\watt\per\kelvin}$ and split a heating
season into three temperature bins: a mild bin near $8\,{}^{\circ}\mathrm{C}$
for $1400\,\si{\hour}$, a cold bin near $0\,{}^{\circ}\mathrm{C}$ for
$900\,\si{\hour}$, and a severe bin near $-8\,{}^{\circ}\mathrm{C}$ for
$300\,\si{\hour}$. The heat demand in each bin is $U_{\text{tot}}$ times the
indoor-outdoor gap (indoor $20\,{}^{\circ}\mathrm{C}$) times the hours:
$267\times 12\times 1400 \approx 4.5$, $267\times 20\times 900 \approx 4.8$,
and $267\times 28\times 300 \approx 2.2\,\si{\mega\watt\hour}$, a season total
of $\sim 11.5\,\si{\mega\watt\hour}$. The COP tracks the Carnot fraction of
Figure \ref{fig:vol04ch14:hp-cop-temperature}: at $0.45$ of Carnot with an
indoor sink of $293\,\si{\kelvin}$ the unit reaches $\text{COP} \approx 4.7$
in the mild bin, $\approx 3.6$ in the cold bin, and $\approx 2.4$ in the
severe bin. The electricity in each bin is the heat divided by its COP:
$4.5/4.7 + 4.8/3.6 + 2.2/2.4 \approx 0.96 + 1.33 + 0.92 \approx 3.2\,\si{
\mega\watt\hour}$, so the seasonal COP is $11.5/3.2 \approx 3.6$, above the
flat $3.0$ the degree-day estimate assumed, because most of the hours sit in
the mild bin where the unit runs efficiently. The bin method also exposes the
sizing trap: the severe bin holds only $\sim 19\%$ of the season's heat but
sets the design capacity and the resistance-backup demand, so a unit sized to
carry the severe bin on the compressor alone is oversized and cycles
inefficiently in the mild bin, which is why cold-climate installations size
the compressor near the cold bin and let resistance backup cover the severe
tail.

% Figure: heat-pump coefficient of performance against outdoor temperature.
% The Carnot ceiling (dashed) rises steeply as the outdoor temperature
% approaches the indoor setpoint; a real air-source unit tracks a fixed
% fraction of Carnot (solid) and drops toward the resistance-backup floor
% (COP = 1) in the cold. The crossover with the gas-furnace primary-energy
% equivalent marks where the heat pump stops beating the furnace on primary
% energy for a given grid. Values illustrative.
\begin{figure}[htbp]
\centering
\begin{tikzpicture}
\begin{axis}[
  width=13cm, height=8cm,
  xlabel={outdoor temperature (${}^{\circ}\mathrm{C}$)},
  ylabel={coefficient of performance},
  xmin=-20, xmax=15, ymin=0, ymax=12,
  xtick={-20,-15,-10,-5,0,5,10,15},
  ytick={0,2,4,6,8,10,12},
  axis lines=left,
  tick label style={font=\scriptsize},
  label style={font=\small},
  legend style={font=\scriptsize, at={(0.03,0.97)}, anchor=north west, draw=none},
]
% Carnot ceiling for indoor setpoint 20 C = 293.15 K: COP_c = 293.15/(20 - T)
\addplot[engblue, thick, dashed, domain=-20:15, samples=80] {293.15/(20 - x)};
\addlegendentry{Carnot ceiling ($T_h=20\,{}^{\circ}\mathrm{C}$)}
% Real air-source unit: 45% of Carnot, with a resistance-backup floor at COP = 1
\addplot[engorange, very thick, domain=-20:15, samples=80] {max(1, 0.45*293.15/(20 - x))};
\addlegendentry{real unit ($0.45\times$Carnot, floor $1$)}
% Gas-furnace primary-energy break-even for a grid at 2.5x fuel intensity
\addplot[engred, thick, dotted] coordinates {(-20,2.5)(15,2.5)};
\addlegendentry{furnace break-even (grid $2.5\times$)}
% marker at design point -10 C
\addplot[engblue, only marks, mark=*, mark size=2pt] coordinates {(-10,2.9)};
\node[anchor=west, font=\scriptsize, color=engblue]
  at (axis cs:-9.5,3.6) {design point};
\end{axis}
\end{tikzpicture}
\caption[Heat-pump COP against outdoor temperature]{The coefficient of
performance of an air-source heat pump falls as the outdoor air cools,
because the Carnot ceiling (dashed) collapses toward one as the source and
sink temperatures converge on the heating side. A real unit tracks a fixed
fraction of Carnot (solid, here $0.45$) and is held above the electric-
resistance floor of $\text{COP}=1$ by the backup element in the coldest
hours. The dotted line marks the primary-energy break-even against a gas
furnace on a grid whose electricity carries $2.5\times$ the fuel's carbon
intensity per delivered unit: above the line the heat pump wins, and the
window narrows in the cold, which is why cold-climate sizing must check the
seasonal average COP rather than the mild-day value. Quantities illustrative.}
\label{fig:vol04ch14:hp-cop-temperature}
\end{figure}


\paragraph{The transmission break-even distance from a two-term cost model.}
The HVDC-versus-HVAC comparison earlier quoted a break-even of $\sim 600$ to
$800\,\si{\kilo\meter}$; the number falls out of the same fixed-plus-slope
cost structure the storage crossover used. Model each option's total cost as
a fixed terminal cost plus a per-kilometre line cost, $C(d) = C_0 + c\,d$.
Alternating current carries a low terminal cost, $C_{0,\text{AC}} \approx
\$0.20\text{B}$ for the substations, and a steep slope, $c_{\text{AC}}
\approx \$1.30\text{M}/\si{\kilo\meter}$ once the line, the reactive-
compensation stations spaced along it, and the higher loss are counted.
Direct current carries a high terminal cost, $C_{0,\text{DC}} \approx
\$0.80\text{B}$ for the rectifier and inverter stations, and a shallow slope,
$c_{\text{DC}} \approx \$0.65\text{M}/\si{\kilo\meter}$, with no charging
current to compensate. The break-even distance is where the two lines meet,
$C_{0,\text{AC}} + c_{\text{AC}}d = C_{0,\text{DC}} + c_{\text{DC}}d$, so
\[
d^\ast = \frac{C_{0,\text{DC}} - C_{0,\text{AC}}}{c_{\text{AC}} -
c_{\text{DC}}} = \frac{0.80 - 0.20}{(1.30 - 0.65)\times 10^{-3}}
\approx 920\,\si{\kilo\meter},
\]
the crossing in Figure \ref{fig:vol04ch14:transmission-cost-distance}. Below
it the cheaper terminals win and alternating current is chosen; above it the
cheaper line and the absence of charging current win and direct current is
chosen. The break-even collapses to tens of kilometres for submarine or
underground cable, because a buried or submerged AC cable draws a large
charging current that grows with length and consumes the line's capacity
before the resistive loss even matters, which is why every long submarine
link is direct current regardless of the overhead break-even. The structure
is the same crossover the storage market showed: a technology with a high
fixed cost and a low marginal cost wins at large scale, and one with the
reverse wins at small scale, with geography and the cable-versus-overhead
choice moving the crossing.

% Figure: total transmission cost against corridor length for HVAC and HVDC.
% HVAC has cheaper terminals but a steeper per-kilometre slope (line cost
% plus reactive-compensation stations and higher loss); HVDC pays a large
% fixed converter cost but a shallower slope. The lines cross at the
% break-even distance, above which DC is cheaper. Straight-line cost models.
% Values illustrative, current as of 2025.
\begin{figure}[htbp]
\centering
\begin{tikzpicture}
\begin{axis}[
  width=13cm, height=8cm,
  xlabel={corridor length (km)},
  ylabel={total cost (USD, billions)},
  xmin=0, xmax=2000, ymin=0, ymax=3.5,
  xtick={0,400,800,1200,1600,2000},
  ytick={0,0.5,1.0,1.5,2.0,2.5,3.0,3.5},
  axis lines=left,
  tick label style={font=\scriptsize},
  label style={font=\small},
  legend style={font=\scriptsize, at={(0.03,0.97)}, anchor=north west, draw=none},
]
% HVAC: low fixed 0.20 B, steep slope 1.30 B per 1000 km
\addplot[engblue, very thick, domain=0:2000, samples=2] {0.20 + 1.30e-3*x};
\addlegendentry{HVAC (cheap terminals, steep line)}
% HVDC: high fixed 0.80 B (converters), shallow slope 0.65 B per 1000 km
\addplot[engorange, very thick, domain=0:2000, samples=2] {0.80 + 0.65e-3*x};
\addlegendentry{HVDC (converter cost, cheap line)}
% break-even: 0.20 + 1.30e-3 d = 0.80 + 0.65e-3 d => 0.65e-3 d = 0.60 => d = 923 km
\addplot[engred, thick, dashed] coordinates {(923,0)(923,2.0)};
\addplot[engred, only marks, mark=*, mark size=2pt] coordinates {(923,1.40)};
\node[anchor=west, font=\scriptsize, color=engred]
  at (axis cs:940,1.05) {break-even $\approx 920$ km};
\end{axis}
\end{tikzpicture}
\caption[Transmission cost against distance for HVAC and HVDC]{Total cost of
a bulk transmission corridor against its length for alternating and direct
current (current as of 2025). Alternating current pays little for its
terminal equipment but a steep cost per kilometre, because the line itself,
the reactive-compensation stations, and the higher loss all grow with
length. Direct current pays a large fixed cost for its rectifier and inverter
stations but a shallow cost per kilometre and draws no charging current, so
its total-cost line starts high and rises slowly. The two cross at a
break-even distance, below which alternating current wins on the cheaper
terminals and above which direct current wins on the cheaper line, which is
why the longest overhead corridors and all long submarine links are built as
direct current. Quantities illustrative.}
\label{fig:vol04ch14:transmission-cost-distance}
\end{figure}


\paragraph{Battery C-rate and the internal-resistance penalty.} A battery's
round-trip efficiency is not a constant; it falls as the charge and discharge
rate rises, because the internal resistance dissipates power that grows with
the square of the current. Define the \engterm{C-rate} as the current
expressed in multiples of the rate that would fully discharge the cell in one
hour, so a $2\,\text{C}$ discharge empties the cell in half an hour. Take a
cell of usable energy $E = 0.64\,\si{\kilo\watt\hour}$, nominal voltage
$V = 3.2\,\si{\volt}$, and internal resistance $r = 8\,\si{\milli\ohm}$. At a
$1\,\text{C}$ discharge the current is $I = E/(V\cdot 1\,\si{\hour}) =
640/3.2 = 200\,\si{\ampere}$, so the resistive loss is $I^2 r = (200)^2\times
0.008 = 320\,\si{\watt}$ against a delivered power of $V I = 640\,\si{\watt}$,
a loss fraction of $\sim 33\%$ at that crude single-cell reading. Real packs
run the cells far below their $1\,\text{C}$ full-power point at grid duty: at
a $0.25\,\text{C}$ discharge ($50\,\si{\ampere}$) the loss is $(50)^2\times
0.008 = 20\,\si{\watt}$ against $160\,\si{\watt}$ delivered, $\sim 12.5\%$,
and the resistive loss scales as the square of the C-rate while the delivered
power scales linearly, so the loss fraction rises linearly with the C-rate.
This is why a battery rated for many hours of storage is more efficient per
cycle than the same cells sized for a fast burst: a $4\,\si{\hour}$ system
discharges at $0.25\,\text{C}$ and loses little to resistance, while a
$15$-minute frequency-response system discharges at $4\,\text{C}$ and pays a
steep resistive penalty and a large heat-rejection load. The duration a
storage system is sized for therefore sets not only its energy cost, through
the power-energy split of the earlier mastery box, but its round-trip
efficiency, through the C-rate.

\paragraph{EV depot power from the diversity factor.} The fleet-charging
design exercise sized a depot on the energy floor; the number that actually
sets the service is the coincident peak, and the coincident peak follows from
the diversity of the fleet, not the sum of the chargers. Take $100$ vehicles
each drawing up to $11\,\si{\kilo\watt}$, so the connected load is $1100\,\si{
\kilo\watt}$ if every charger runs flat out at once. They do not: each vehicle
needs only $\sim 60\,\si{\kilo\watt\hour}$ over a $13\,\si{\hour}$ window, an
energy-limited average of $4.6\,\si{\kilo\watt}$, so the aggregate energy
floor is $460\,\si{\kilo\watt}$. The coincident peak sits between the floor
and the connected load, and the \engterm{diversity factor}, the ratio of the
sum of individual peaks to the coincident peak, measures how much the
staggering helps. Uncontrolled, vehicles plug in on arrival and charge at
full power until done, so the peak clusters near the $18{:}00$ arrival and the
coincidence is high, a diversity factor of perhaps $1.5$ giving a peak of
$\sim 1100/1.5 \approx 730\,\si{\kilo\watt}$. Smart charging spreads the
draw across the window and holds the aggregate near the $460\,\si{\kilo\watt}$
floor, a diversity factor of $\sim 2.4$, which is the difference between a
$0.75\,\si{\mega\watt}$ interconnection and a $0.5\,\si{\mega\watt}$ one. The
design lesson is that management buys down the coincident peak, and the peak,
not the connected load or the energy, sets the transformer and the
interconnection cost; a depot designed to the connected load pays for a
service it will never use, and one designed to the energy floor without
management will trip on the arrival cluster.

\paragraph{Single-axis tracking and the energy it buys.} A fixed-tilt panel
faces one direction all day; a single-axis tracker rotates about a
north-south axis to follow the sun east to west, and the energy gain is worth
pricing against the tracker's cost and failure exposure. The instantaneous
capture of a flat collector is the cosine of the angle between its normal and
the sun, so a fixed panel that is normal to the sun at noon captures
$\cos\theta$ as the sun moves off by an hour-angle $\theta$, losing the
morning and evening to the cosine. A tracker holds the panel near normal
across the day, so it captures close to the full beam whenever the sun is up.
Modelling a clear day as the sun sweeping $\pm 90^\circ$ of hour angle, the
fixed panel collects $\int_{-\pi/2}^{\pi/2}\cos\theta\,d\theta = 2$ in units
of the noon rate, while the tracker collects the full $\pi$ over the same
sweep (the sun is above the horizon for the half-day the integral covers), a
ratio of $\pi/2 \approx 1.57$ in the idealised limit. Real single-axis
trackers deliver a more modest $\sim 15$ to $25\%$ annual energy gain over
fixed tilt, because diffuse light (which a tracker does not concentrate),
backtracking to avoid row-to-row shading, and cloudy hours all erode the
clear-sky geometric gain. The gain is largest at low latitudes and in clear
climates and shrinks toward the poles and in cloudy ones. The engineering
trade is the $\sim 15$ to $25\%$ more energy from the same panels against the
tracker's added capital, its moving parts and their maintenance, and its
larger land footprint from the row spacing that avoids self-shading, which is
why utility-scale plants in sunny regions track and rooftop and high-latitude
plants usually do not.

\paragraph{The maximum-power point and the fill factor.} A photovoltaic
module does not deliver a fixed power; it delivers whatever the product of its
operating voltage and the current that voltage draws happens to be, and the
job of the inverter's tracker is to sit at the product's maximum. Take a
module whose short-circuit current is $I_{\text{sc}} = 10\,\si{\ampere}$ (the
current when the terminals are shorted and the voltage is zero) and whose
open-circuit voltage is $V_{\text{oc}} = 40\,\si{\volt}$ (the voltage when no
current is drawn). The power is zero at both ends, at $V = 0$ because there is
no voltage and at $V = V_{\text{oc}}$ because there is no current, and it
peaks in between at the \engterm{maximum-power point}, where the current is
just beginning to fall off the flat part of the $I$--$V$ curve (Figure
\ref{fig:vol04ch14:pv-iv-curve}). For this module the peak sits near
$V_{\text{mp}} \approx 32\,\si{\volt}$ and $I_{\text{mp}} \approx 9.0\,\si{
\ampere}$, so the peak power is $P_{\text{mp}} = 32\times 9.0 = 288\,\si{
\watt}$. The quality of the cell is measured by the \engterm{fill factor}, the
ratio of the peak power to the product of the two end-point values,
\[
\text{FF} = \frac{P_{\text{mp}}}{V_{\text{oc}}\,I_{\text{sc}}}
= \frac{288}{40\times 10} = 0.72,
\]
a good silicon module reaching $\sim 0.75$ to $0.82$ and a poor or degraded
one falling below $0.70$. The fill factor is the geometric measure of how
square the $I$--$V$ knee is: a perfect cell would hold its full current right
up to $V_{\text{oc}}$ and give a fill factor of one, and the internal series
resistance and the diode physics round the knee and pull it below. The
tracker matters because the maximum-power point moves: rising irradiance lifts
the current and shifts the point up, and rising temperature lowers the voltage
and shifts it left, so a fixed operating voltage would spill a large fraction
of the available power, which is why every grid-tied inverter runs a
maximum-power-point-tracking loop that continually hunts the peak.

\paragraph{Temperature derating of a module.} The nameplate of a module is
quoted at a cell temperature of $25\,{}^{\circ}\mathrm{C}$, and a module in
service runs hotter, so the delivered power sits below the nameplate whenever
the sun is strong enough to matter. Silicon carries a negative temperature
coefficient of power near $\gamma \approx -0.4\,\%$ per kelvin, so the power
relative to rating is $P/P_{\text{rated}} = 1 + \gamma\,(T_{\text{cell}} -
25)$ (Figure \ref{fig:vol04ch14:pv-temperature-derate}). A cell temperature
rises above the air temperature under sun; a common rule puts the cell at
$\sim 30\,\si{\kelvin}$ above ambient at full irradiance, so a
$30\,{}^{\circ}\mathrm{C}$ afternoon gives a cell temperature near
$60\,{}^{\circ}\mathrm{C}$ and a derating of
\[
\frac{P}{P_{\text{rated}}} = 1 + (-0.004)(60 - 25) = 1 - 0.14 = 0.86,
\]
a $14\,\%$ loss of power exactly when the irradiance is highest. The
derating compounds with the inverter and soiling losses to pull the delivered
energy below the naive product of peak rating and insolation, and it is why a
hot desert with abundant sun does not deliver proportionally more energy per
installed watt than the peak rating suggests, why array mounting is designed
to let air circulate behind the modules, and why the same nameplate array
delivers more annual energy in a cool, bright climate than in a hot one. The
effect runs the other way in the cold: a bright winter day at
$0\,{}^{\circ}\mathrm{C}$ cell temperature delivers $1 + (-0.004)(0 - 25) =
1.10$, ten percent above nameplate, so the instantaneous peak a system must be
wired for is set by cold, bright conditions, not by the hottest hour.

\paragraph{Net-load-duration and the firm capacity that remains.} The
load-duration section argued that variable renewables reshape the curve the
dispatchable fleet must serve; the reshaping is quantified by subtracting the
renewable output hour by hour and re-sorting. Take a grid whose gross
load-duration curve runs from a peak of $1.0$ per unit down to a base of
$0.35$ per unit (Figure \ref{fig:vol04ch14:net-load-duration}, blue), with an
annual energy that is the area under it, here $\sim 0.60$ per unit of
peak-times-year. Add a variable renewable fleet supplying $\sim 40\,\%$ of the
annual energy. Subtracting its hourly output and re-sorting gives the
net-load-duration curve (orange dashed): the middle is hollowed where midday
output meets demand, the right end dips below zero where output exceeds demand
and the surplus is curtailed or stored, and, decisively, the top-left barely
moves, because the annual peak arrives on a winter evening when the renewable
output is low. The dispatchable fleet is built to the net curve, so its
required energy falls by the $\sim 40\,\%$ the renewables supply, but its
required firm capacity, read off the left edge, falls hardly at all: the peak
net load is still $\sim 0.95$ per unit against the gross $1.0$. This is the
capacity-credit argument written as a curve: renewables buy down energy fast
and firm capacity slowly, and the gap between the two is the firm-capacity
procurement the transition must pay for separately. Integrating the negative
tail of the net curve gives the curtailed energy directly, and reading the
left edge gives the firm capacity the grid must still hold; the same
construction, run across many weather years, is how a planner sizes the firm
fleet against the worst renewable drought rather than the average year.

\paragraph{Storage holding time and self-discharge.} A storage system loses
energy not only on the round trip but while it simply holds charge, and the
holding loss decides which technology suits which time scale. Model the state
of charge of an idle store as decaying at a fractional self-discharge rate
$\lambda$ per day, so energy held for $t$ days returns $E(t) = E_0
(1-\eta_{\text{sd}})\,e^{-\lambda t}$ after the one-way conversion efficiency
$\sqrt{\eta_{\text{rt}}}$ on the way out. Lithium-ion self-discharges at
$\sim 0.1\,\%$ per day, so over a day of holding it loses almost nothing, and
its round-trip efficiency of $\sim 0.88$ dominates: a day-cycled battery
returns $\sim 0.88 \times 0.999 \approx 0.88$ of the input. A flywheel or a
supercapacitor self-discharges at several percent per hour, so it is superb
for a burst held for seconds and useless for a store held for a day. A
pumped-hydro reservoir loses to evaporation and seepage at perhaps a fraction
of a percent per day and holds for weeks with little loss, and a chemical
store such as hydrogen holds for months with almost no
self-discharge once sealed, which is why the seasonal role falls to it despite
its poor round trip. The lesson pairs with the power-energy split: a
technology is chosen for a duration not only by its cost-per-hour but by
whether it can hold the charge across that duration without leaking it away,
so the storage market splits along two axes at once, the cost structure and
the holding loss, and both point lithium-ion at the hours, pumped hydro at the
days, and chemical storage at the seasons.

\paragraph{The carbon a marginal electrified load actually saves.} Whether
electrifying a car or a furnace cuts carbon turns on the intensity of the
electricity that serves the new load, and the honest accounting uses the
\engterm{marginal} intensity, the emissions of the generator that ramps to
serve the added kilowatt-hour, not the grid average. Take a heat pump drawing
$5300\,\si{\kilo\watt\hour}/\text{yr}$ (the degree-day example's figure)
against the gas furnace it replaces, which burnt $17{,}800\,\si{\kilo\watt
\hour}$ of fuel emitting $\sim 0.20\,\si{\kilogram}\,\text{CO}_2$ per
kilowatt-hour of fuel, or $\sim 3560\,\si{\kilogram}\,\text{CO}_2/\text{yr}$.
If the marginal grid generator is combined-cycle gas at $\sim 490\,\si{\gram}
\,\text{CO}_2/\si{\kilo\watt\hour}$ (the lifecycle figure of Figure
\ref{fig:vol04ch14:lifecycle-carbon}), the heat pump emits $5300\times 0.490
\approx 2600\,\si{\kilogram}/\text{yr}$, a $\sim 27\,\%$ cut. If the marginal
generator is coal at $\sim 820\,\si{\gram}/\si{\kilo\watt\hour}$, the heat
pump emits $5300\times 0.820 \approx 4350\,\si{\kilogram}/\text{yr}$, worse
than the furnace it replaced. If the marginal generator is a low-carbon
source at $\sim 40\,\si{\gram}/\si{\kilo\watt\hour}$, the emissions collapse
to $\sim 210\,\si{\kilogram}/\text{yr}$, a $\sim 94\,\%$ cut. The factor-of-
three efficiency of the heat pump buys the same three-to-one headroom on the
grid intensity: the heat pump wins whenever the marginal electricity is
cleaner than about three times the fuel's intensity per delivered
kilowatt-hour, which holds on almost every modern grid and improves as the
grid decarbonises. The reading generalises to every electrified load: the
carbon benefit is real but conditional on the grid, and the same electric car
is a large win in a hydro-and-nuclear region and a small one, or none, on a
coal grid, which is why electrification and grid decarbonisation are one
programme, not two.

\paragraph{The largest-single-infeed rule and the reserve it sets.} A grid
carries spinning and standing reserve to survive the loss of its largest
single element, and the rule sizes that reserve directly. The
\engterm{largest single infeed} is the biggest block that can be lost at once:
a $1.6\,\si{\giga\watt}$ nuclear unit, a $2\,\si{\giga\watt}$ HVDC import, a
large busbar. The operating reserve must at least cover it, because the $N-1$
standard demands the grid ride through any single loss without shedding load.
Take a $50\,\si{\giga\watt}$ grid whose largest infeed is a $1.6\,\si{\giga
\watt}$ HVDC link. The primary-and-secondary reserve must be at least
$1.6\,\si{\giga\watt}$, held ready to ramp within the seconds-to-minutes the
frequency dynamics allow, and it is carried as part-loaded synchronous plant,
fast storage, and contracted demand response. The rule has a sharp
consequence for the transition: building ever-larger single units or single
interconnectors raises the reserve the whole grid must carry, so there is a
system cost to bigness that the unit's own economics do not see. A $3\,\si{
\giga\watt}$ interconnector is cheaper per megawatt than a $1.6\,\si{\giga
\watt}$ one, but it forces the receiving grid to carry $3\,\si{\giga\watt}$ of
reserve against its loss, and the reserve is not free. The same logic caps the
size of the largest inverter-based in-feed on a low-inertia grid, because the
frequency nadir after losing it (the swing-equation result) deepens with the
size of the loss and shallows with the inertia, so a grid shedding its
synchronous inertia must also shrink its largest credible loss to keep the
nadir above the load-shedding threshold.

\paragraph{Capacity credit from the effective load-carrying capability.} The
capacity-credit mastery box named the quantity; the effective-load-carrying-
capability construction computes it. A resource's \engterm{effective
load-carrying capability} (ELCC) is the amount of extra firm load the system
can serve, at the same reliability, once the resource is added. Take a
$1\,\si{\giga\watt}$ solar fleet on a summer-peaking grid. Adding it lets the
grid serve more load in the sunny afternoon hours but nothing in the evening
peak after sunset, so the reliability-equivalent firm load it adds is far
below its nameplate. Suppose the loss-of-load hours cluster in the
$18{:}00$--$21{:}00$ window, when the solar output has fallen to $\sim 10\,\%$
of nameplate; then the solar fleet carries $\sim 0.10\times 1000 = 100\,\si{
\mega\watt}$ of firm load at the hours that set reliability, an ELCC of
$\sim 10\,\%$. The same $1\,\si{\giga\watt}$ of wind, if the wind blows
$\sim 30\,\%$ of nameplate on average during the peak hours and is weakly
correlated with them, carries an ELCC of $\sim 25$ to $35\,\%$. A
$1\,\si{\giga\watt}/4\,\si{\hour}$ battery, able to discharge across the whole
evening peak, carries an ELCC near $\sim 90\,\%$ as long as the peak is
shorter than its duration, and its ELCC collapses once the peak stretches
beyond four hours or once enough batteries saturate the evening. The ELCC is
the honest firm-capacity currency: it is measured against the hours that
actually threaten reliability, it falls with penetration as those hours shift,
and it is what a capacity market should pay for, rather than nameplate, which
counts a megawatt of midday solar and a megawatt of dispatchable gas as if
they were the same firm product.

\paragraph{Interconnection smoothing across weakly correlated regions.} The
transmission mastery box asserted that averaging renewables over a wide area
smooths the aggregate; the statistics fix how much. Model $N$ regions each
producing a renewable output with the same mean $\mu$ and standard deviation
$\sigma$, with a pairwise correlation coefficient $\rho$ between any two. The
variance of the average output of the $N$ regions is
\[
\operatorname{Var}\!\left(\bar{P}\right) = \frac{\sigma^2}{N}\bigl[1 +
(N-1)\rho\bigr],
\]
so the relative variability of the fleet, the ratio of the aggregate standard
deviation to its mean, falls from $\sigma/\mu$ for one region toward a floor of
$\sigma\sqrt{\rho}/\mu$ as $N\to\infty$. For independent regions ($\rho = 0$)
the floor is zero and the smoothing follows the familiar $1/\sqrt{N}$: five
independent regions cut the relative variability to $\sim 45\,\%$ of a single
region, and twenty regions to $\sim 22\,\%$. Real regions are not independent,
because the same weather system spans several of them; a typical continental
correlation of $\rho \approx 0.3$ sets a floor at $\sqrt{0.3} \approx 0.55$ of
the single-region variability, so no amount of interconnection removes more
than about $45\,\%$ of the swing, and the first few interconnections capture
most of it. The reading matches the mastery box: the smoothing is real,
cheaper than storage per unit of firmed energy, and subject to diminishing
returns set by the residual correlation, so interconnection and storage are
complements, the wires flattening the correlated bulk and the storage
covering the residual that the correlation floor leaves behind.

\paragraph{Transformer no-load loss over a lifetime.} A distribution
transformer's iron loss runs every hour it is energised, whether or not it
serves any load, and over a transformer's long life that standing loss adds a
running cost worth optimising against the purchase price. Take two
$1\,\si{\mega\watt}$ transformers: a standard unit with a no-load loss of
$2.0\,\si{\kilo\watt}$ costing $\$18{,}000$, and a low-loss amorphous-core
unit with a no-load loss of $0.8\,\si{\kilo\watt}$ costing $\$26{,}000$. The
standing loss runs all $8760\,\si{\hour}$ of the year, so the standard unit
wastes $2.0\times 8760 = 17{,}520\,\si{\kilo\watt\hour}/\text{yr}$ and the
low-loss unit $0.8\times 8760 = 7008\,\si{\kilo\watt\hour}/\text{yr}$, a
saving of $\sim 10{,}500\,\si{\kilo\watt\hour}/\text{yr}$. At an energy price
of $\$0.08\,\si{\per\kilo\watt\hour}$ that saves $\sim \$840/\text{yr}$, and
over a $30$-year transformer life, discounted at $7\,\%$ (the annuity factor
$[1 - 1.07^{-30}]/0.07 \approx 12.4$), the present value of the saving is
$\sim 840\times 12.4 \approx \$10{,}400$, comfortably above the $\$8000$ price
premium. The low-loss unit is the cheaper choice on a lifetime basis even
though it costs more to buy, and the standing nature of the iron loss is what
makes it so: a load-dependent loss can be avoided by idling the asset, but the
iron loss is paid every hour the transformer is energised, which is nearly
every hour of a thirty-year life, so it is discounted lightly and accumulates
into a large lifetime cost. The same total-cost-of-ownership arithmetic drives
minimum-efficiency standards for the millions of distribution transformers on
a grid, where the aggregate standing loss is a measurable fraction of national
generation.

\paragraph{Cooling water and the heat a thermal plant rejects.} A steam plant
turns only a fraction of its fuel heat into electricity, and the rest leaves at
the condenser, where it is carried away by cooling water; the flow that duty
demands is large enough to shape where a plant can be built. Take a
$1\,\si{\giga\watt}$-electric plant at a thermal efficiency of $\eta = 0.38$.
The fuel heat is $P_{\text{th}} = 1000/0.38 \approx 2630\,\si{\mega\watt}$, so
the heat rejected at the condenser is $P_{\text{rej}} = P_{\text{th}} -
P_{\text{elec}} \approx 2630 - 1000 = 1630\,\si{\mega\watt}$, more than the
electricity the plant sells. A once-through cooling design that lets the
cooling water warm by $\Delta T = 10\,\si{\kelvin}$ needs a flow of
\[
\dot{m} = \frac{P_{\text{rej}}}{c_p\,\Delta T}
= \frac{1.63\times 10^9\,\si{\watt}}
{(4180\,\si{\joule\per\kilogram\per\kelvin})(10\,\si{\kelvin})}
\approx 3.9\times 10^4\,\si{\kilogram\per\second},
\]
that is $\sim 39\,\si{\meter\cubed\per\second}$ of water, the flow of a small
river passed through the plant every second. A plant on a river or coast can
draw that once-through and return it warmed; an inland plant without such a
source must reject the heat to the air through a cooling tower, evaporating a
few percent of the flow, $\sim 0.5$ to $1\,\si{\meter\cubed\per\second}$
consumed rather than merely warmed, which is a real water-consumption cost in a
dry region. The rejected heat is set by the efficiency, so the same arithmetic
rewards a higher-efficiency plant twice: a combined-cycle unit at $\eta = 0.60$
rejects only $1000\times(1/0.60 - 1) \approx 670\,\si{\mega\watt}$ for the same
electricity, less than half the once-through flow, which is why the move from
open Rankine plants to combined cycle eased the water and thermal-discharge
constraint as much as it cut the fuel bill.

\paragraph{Ramp rate and the cost of cycling a coal plant.} A large thermal
unit cannot change output instantly, and the ramp limit sets whether it can
follow the net-load swings a high-renewable grid throws at it. Take a
$600\,\si{\mega\watt}$ coal unit with a rated ramp rate of $\sim 3\,\%$ of
rating per minute, so $\sim 18\,\si{\mega\watt}/\si{\minute}$, and a minimum
stable load of $40\,\%$, that is $240\,\si{\mega\watt}$, below which combustion
becomes unstable. To swing from minimum load to full output covers
$600 - 240 = 360\,\si{\mega\watt}$, which at $18\,\si{\mega\watt}/\si{\minute}$
takes $360/18 = 20\,\si{\minute}$. The duck-curve evening ramp of Figure
\ref{fig:vol04ch14:ramp-cycling} can demand the grid raise output far faster
than that across the whole fleet, so a coal unit lags the ramp (the gap the
figure shades) and faster resources must cover the shortfall. The minimum-load
floor has its own cost: through the midday solar trough the unit is held at
$240\,\si{\mega\watt}$ rather than shut down, because a cold start takes
several hours and a thermal-fatigue toll on the boiler and turbine, so it burns
fuel at a poor part-load efficiency (heat rate rising perhaps $10$ to $15\,\%$
below the design point) and emits carbon while earning little. The reading is
that a plant built for steady base-load duty is expensive to use as a
follower: its ramp rate caps how fast it helps, its minimum load forces
inefficient part-load running, and its start cost penalises cycling, which is
why coal retires first from a grid that has come to need flexibility rather
than bulk energy, and why fast-ramping gas, storage, and demand response take
over the following duty.

% Figure: a dispatchable plant following a steep net-load ramp against its
% ramp-rate limit and minimum stable load. The net-load demand (blue) rises
% faster in the evening than the plant can follow at its rated ramp rate, so
% the plant output (orange) lags on the up-ramp and is floored at its minimum
% stable generation overnight rather than shutting down. The shaded gap is the
% unserved ramp that faster storage or another unit must cover. Values
% illustrative.
\begin{figure}[htbp]
\centering
\begin{tikzpicture}
\begin{axis}[
  width=13cm, height=8cm,
  xlabel={hour of day},
  ylabel={power (fraction of plant rating)},
  xmin=0, xmax=24, ymin=0, ymax=1.15,
  xtick={0,4,8,12,16,20,24},
  ytick={0,0.2,0.4,0.6,0.8,1.0},
  axis lines=left,
  tick label style={font=\scriptsize},
  label style={font=\small},
  legend style={font=\scriptsize, at={(0.03,0.97)}, anchor=north west, draw=none},
]
% net-load demand the plant is asked to follow: deep midday trough (solar),
% steep evening ramp, high evening peak
\addplot[engblue, very thick] coordinates {
  (0,0.55)(2,0.48)(4,0.45)(6,0.50)(8,0.60)(10,0.45)(12,0.28)
  (14,0.25)(16,0.40)(17,0.60)(18,0.85)(19,1.00)(20,0.95)
  (21,0.80)(22,0.70)(24,0.58)};
\addlegendentry{net-load demand}
% plant output: floored at minimum stable load 0.30, ramp-limited on the
% steep evening rise so it lags the demand
\addplot[engorange, very thick] coordinates {
  (0,0.55)(2,0.48)(4,0.45)(6,0.50)(8,0.60)(10,0.45)(12,0.30)
  (14,0.30)(16,0.40)(17,0.55)(18,0.72)(19,0.89)(20,0.95)
  (21,0.80)(22,0.70)(24,0.58)};
\addlegendentry{plant output (ramp- and floor-limited)}
% minimum stable load line
\addplot[enggray, thick, dashed] coordinates {(0,0.30)(24,0.30)};
\node[anchor=west, font=\scriptsize, color=enggray]
  at (axis cs:0.5,0.235) {minimum stable load};
% mark the ramp gap in the early evening
\node[anchor=west, font=\scriptsize, color=engred]
  at (axis cs:15.0,0.80) {ramp gap};
\draw[engred, -{Latex[length=2mm]}] (axis cs:16.7,0.78) -- (axis cs:18.0,0.79);
\end{axis}
\end{tikzpicture}
\caption[A dispatchable plant against its ramp and minimum-load limits]{A
dispatchable plant asked to follow the evening net-load ramp (blue) cannot
track it exactly. Two constraints bind. On the steep early-evening rise the
plant output (orange) lags the demand because its rated ramp rate limits how
fast it can raise power, and the gap between demand and output is the unserved
ramp that faster storage or a second unit must cover. Overnight and through
the midday solar trough the output is held at the minimum stable load rather
than shutting down, because cycling the unit off and on stresses it and costs
fuel and time to restart, so it runs part-loaded and inefficient through the
low hours. The steeper the net-load ramp grows with renewable penetration, the
harder both limits bind, which is the operational reason a high-renewable grid
values fast-ramping flexibility over cheap bulk energy. Quantities
illustrative.}
\label{fig:vol04ch14:ramp-cycling}
\end{figure}


\paragraph{Combined heat and power and the primary energy it saves.}
Delivering electricity and heat from one fuel stream beats delivering them
separately, and the saving is a primary-energy calculation the transition
leans on for dense loads. Compare two ways to supply $35\,\si{\kilo\watt}$ of
electricity and $48\,\si{\kilo\watt}$ of heat. The separate path burns fuel in
a power plant at $\eta_e = 0.35$ and in a boiler at $\eta_h = 0.90$: the plant
burns $35/0.35 = 100\,\si{\kilo\watt}$ of fuel and the boiler
$48/0.90 \approx 53\,\si{\kilo\watt}$, a total of $\sim 153\,\si{\kilo\watt}$
of fuel for the two products. A combined-heat-and-power unit burns one fuel
stream and captures the heat the power generation would otherwise reject: at an
electrical efficiency of $0.35$ and a heat-recovery efficiency of $0.50$, a
single stream of $F$ delivers $0.35F$ electricity and $0.50F$ heat, so meeting
the $35\,\si{\kilo\watt}$ electrical demand fixes $F = 100\,\si{\kilo\watt}$
and yields $50\,\si{\kilo\watt}$ of heat, enough to cover the
$48\,\si{\kilo\watt}$ heat demand from the same $100\,\si{\kilo\watt}$ of fuel
(Figure \ref{fig:vol04ch14:chp-energy}). The fuel saving is $153 - 100 =
53\,\si{\kilo\watt}$, a third less fuel and a third less carbon for the same
two products, and the plant's total fuel-utilisation efficiency is
$(35 + 48)/100 = 0.83$ against the $0.54$ of the separate path. The saving is
conditional: the heat must be wanted near the generator and at the time it is
produced, because heat does not travel far without loss, so combined heat and
power suits a district-heating network, an industrial site, or a hospital, and
falls apart for a remote bulk plant whose rejected heat has no local buyer. The
transition reading is that combined heat and power is a bridge technology, a
large primary-energy saving while the heat demand is met by combustion, whose
role shrinks as heat pumps electrify the heat and the electricity itself
decarbonises.

% Figure: primary-energy comparison of separate generation-plus-boiler against
% a combined-heat-and-power (CHP) plant delivering the same electricity and
% heat. The separate path (left bars) burns fuel in a power plant that rejects
% most of its heat and in a separate boiler; the CHP path (right bars) captures
% the power plant's rejected heat and delivers it as useful heat, so it burns
% less fuel for the same two products. Values illustrative.
\begin{figure}[htbp]
\centering
\begin{tikzpicture}
\begin{axis}[
  width=13cm, height=8cm,
  ybar stacked,
  bar width=28pt,
  xlabel={supply path},
  ylabel={fuel and output (units of useful output)},
  xmin=-0.6, xmax=1.6, ymin=0, ymax=2.6,
  xtick={0,1},
  xticklabels={separate: plant $+$ boiler, combined heat and power},
  ytick={0,0.5,1.0,1.5,2.0,2.5},
  axis lines=left,
  tick label style={font=\scriptsize},
  label style={font=\small},
  legend style={font=\scriptsize, at={(0.97,0.97)}, anchor=north east, draw=none},
]
% Separate path: electricity 0.35 useful + rejected 0.65 from a plant burning
% 1.0 fuel for 0.35 electricity; plus a boiler burning 0.53 fuel for 0.48 heat
% at 0.90 efficiency. Show as: useful electricity, useful heat, rejected.
% Bar 0 (separate): useful electricity 0.35, useful heat 0.48, rejected 0.70
% Bar 1 (CHP): useful electricity 0.35, useful heat 0.48, rejected 0.37
\addplot[fill=engblue, draw=engblue] coordinates {(0,0.35) (1,0.35)};
\addlegendentry{useful electricity}
\addplot[fill=engorange, draw=engorange] coordinates {(0,0.48) (1,0.48)};
\addlegendentry{useful heat}
\addplot[fill=enggray!45, draw=enggray] coordinates {(0,0.70) (1,0.37)};
\addlegendentry{rejected / loss}
% annotate total fuel
\node[anchor=south, font=\scriptsize, color=engred] at (axis cs:0,1.55) {fuel $\approx 1.53$};
\node[anchor=south, font=\scriptsize, color=engred] at (axis cs:1,1.22) {fuel $\approx 1.20$};
\end{axis}
\end{tikzpicture}
\caption[Primary energy: separate generation versus combined heat and
power]{Primary-energy accounting for delivering the same electricity and heat
two ways. The separate path (left) burns fuel in a power plant that turns
about a third of it into electricity and rejects the rest, and burns more fuel
in a separate boiler for the heat, so its total fuel is the sum of the two and
its rejected energy is large. The combined-heat-and-power path (right) captures
the power plant's rejected heat and delivers it as the useful heat, so it needs
no separate boiler and burns less total fuel for the same two products, cutting
the rejected energy. The saving is real only where the heat is wanted near the
generation and at the time it is produced, which is why combined heat and power
suits district heating, industrial sites, and dense loads rather than remote
bulk generation. Quantities illustrative.}
\label{fig:vol04ch14:chp-energy}
\end{figure}


\paragraph{Voltage drop and the hosting capacity of a distribution feeder.} A
distribution feeder is not a perfect wire, and the voltage it delivers falls
along its length under load and rises under distributed generation, with the
two limits together bounding what the feeder can carry. Model a radial feeder
carrying real power $P$ and reactive power $Q$ to a load at its end, over a line
of resistance $R$ and reactance $X$. The approximate voltage drop from the
sending end to the receiving end, for a feeder at nominal voltage $V$, is
\[
\Delta V \approx \frac{P R + Q X}{V},
\]
the real power dropping voltage across the resistance and the reactive power
across the reactance. Take an $11\,\si{\kilo\volt}$ feeder of
$R = 3\,\si{\ohm}$ and $X = 4\,\si{\ohm}$ carrying $2\,\si{\mega\watt}$ at
power factor $0.95$, so $Q = 2\tan(\arccos 0.95) \approx 0.66\,\text{Mvar}$.
The drop is $\Delta V \approx (2\times 10^6\times 3 + 0.66\times 10^6\times 4)/
(11\times 10^3) \approx (6.0 + 2.6)\times 10^6/(1.1\times 10^4) \approx
780\,\si{\volt}$, that is $\sim 7\,\%$ of nominal, already past the common
$\pm 5\,\%$ statutory band and the sag the blue curve of Figure
\ref{fig:vol04ch14:feeder-voltage} traces. Now let rooftop solar at the feeder
end export $2\,\si{\mega\watt}$ at midday: the flow reverses, the sign of the
drop flips, and the receiving-end voltage rises by a comparable amount above
nominal (the orange curve), approaching the upper limit. The reverse-flow
overvoltage, not the thermal rating of the wire, is what caps how much
distributed solar a feeder can host, the feeder's \engterm{hosting capacity}.
The levers that raise it are the ones that manage the two voltages: on-load tap
changers and line voltage regulators, inverter volt-var control that absorbs
reactive power to pull the voltage back down, and, at the limit, reinforcing
the line to lower $R$ and $X$. The engineering reading is that the constraint on
distributed generation is often local voltage on a weak feeder rather than the
bulk grid's ability to absorb the energy, which is why hosting-capacity maps,
not national generation balances, govern where rooftop solar can connect
without an upgrade.

% Figure: voltage profile along a radial distribution feeder. Under heavy load
% (blue) the voltage sags with distance as current draws a drop across the line
% resistance and reactance, approaching the lower statutory limit at the feeder
% end. With distributed rooftop PV exporting near the end of the feeder (orange)
% the flow reverses in the sunny midday and the voltage rises above nominal,
% approaching the upper limit, the reverse-flow overvoltage that limits how much
% PV a feeder can host. Values illustrative.
\begin{figure}[htbp]
\centering
\begin{tikzpicture}
\begin{axis}[
  width=13cm, height=8cm,
  xlabel={distance along feeder (fraction of length)},
  ylabel={voltage (per unit)},
  xmin=0, xmax=1, ymin=0.92, ymax=1.08,
  xtick={0,0.2,0.4,0.6,0.8,1.0},
  ytick={0.92,0.95,0.98,1.00,1.02,1.05,1.08},
  axis lines=left,
  tick label style={font=\scriptsize},
  label style={font=\small},
  legend style={font=\scriptsize, at={(0.03,0.03)}, anchor=south west, draw=none},
]
% statutory band +/- 5%
\addplot[enggray, thick, dashed] coordinates {(0,1.05)(1,1.05)};
\addplot[enggray, thick, dashed] coordinates {(0,0.95)(1,0.95)};
\node[anchor=west, font=\scriptsize, color=enggray] at (axis cs:0.02,1.062) {upper limit $+5\%$};
\node[anchor=west, font=\scriptsize, color=enggray] at (axis cs:0.02,0.938) {lower limit $-5\%$};
% nominal 1.0
\addplot[black, thin] coordinates {(0,1.00)(1,1.00)};
% heavy-load sag: voltage falls roughly linearly (cumulative I R drop) toward
% the end of the feeder
\addplot[engblue, very thick] coordinates {
  (0,1.00)(0.2,0.988)(0.4,0.976)(0.6,0.965)(0.8,0.956)(1.0,0.950)};
\addlegendentry{heavy evening load (voltage sags)}
% midday PV export: reverse flow lifts the far-end voltage above nominal
\addplot[engorange, very thick] coordinates {
  (0,1.00)(0.2,1.012)(0.4,1.024)(0.6,1.035)(0.8,1.044)(1.0,1.050)};
\addlegendentry{midday PV export (voltage rises)}
\end{axis}
\end{tikzpicture}
\caption[Voltage profile along a radial distribution feeder]{Voltage along a
radial distribution feeder under two opposite conditions. Under a heavy evening
load (blue) the current drawn along the feeder produces a cumulative drop
across the line resistance and reactance, so the voltage sags with distance and
is lowest at the far end, where it approaches the statutory lower limit. When
distributed rooftop solar exports near the end of the feeder in the sunny
midday (orange) the power flow reverses, and the same line impedance now lifts
the voltage above nominal, highest at the far end, approaching the upper limit.
The two limits together bound how much load and how much distributed generation
a feeder can carry, and the reverse-flow overvoltage is the constraint that
caps a feeder's photovoltaic hosting capacity, managed by voltage regulators,
by inverter volt-var control, or by reinforcing the line. Quantities
illustrative.}
\label{fig:vol04ch14:feeder-voltage}
\end{figure}


\paragraph{Flywheel energy for a fast frequency-response burst.} Not every
storage duty is measured in hours; the fastest grid service is arrested in
seconds, and a flywheel sized for it stores far less energy than a battery but
delivers it instantly. A flywheel stores kinetic energy $E = \tfrac12 J
\omega^2$ in a spinning rotor, and because the energy scales with the square of
the speed, doubling the speed quadruples the store. Take a rotor of moment of
inertia $J = 500\,\si{\kilogram\meter\squared}$ spinning between an upper speed
of $\omega_1 = 1600\,\si{\radian\per\second}$ and a lower usable speed of
$\omega_2 = 800\,\si{\radian\per\second}$ (below which the power electronics can
no longer hold the output). The usable energy is
\[
E = \tfrac12 J(\omega_1^2 - \omega_2^2)
= \tfrac12 (500)(1600^2 - 800^2)
= \tfrac12 (500)(1.92\times 10^6) \approx 4.8\times 10^8\,\si{\joule}
\approx 133\,\si{\kilo\watt\hour},
\]
a modest energy, but the flywheel can deliver it at a high power for a short
burst: discharging the $133\,\si{\kilo\watt\hour}$ over $20\,\si{\second}$
gives $133\times 3600/20 \approx 2.4\times 10^4\,\si{\kilo\watt} = 24\,\si{
\mega\watt}$. The flywheel is therefore a power device, not an energy device: it
suits the seconds-scale duty of arresting the frequency nadir after a trip (the
synthetic-inertia and fast-frequency-response role the inertia mastery box
named) and the ride-through of a momentary sag, and it is useless for the
hours-scale shifting a battery does, because its several-percent-per-hour
self-discharge would drain it (the holding-time example). The engineering
reading is that the storage market splits not only by the cost-per-hour and the
holding loss but by the response speed: a flywheel or a supercapacitor holds
the seconds, a battery the hours, and a reservoir or a fuel the days and
seasons, and a grid losing its synchronous inertia buys the seconds back with
exactly these fast devices.

\paragraph{The abatement cost of a carbon-cutting measure.} The transition is a
sequence of measures that each cut carbon at some cost, and comparing them
needs a common currency, the \engterm{cost of abatement} in dollars per tonne
of carbon dioxide avoided. The measure's abatement cost is its net annual cost
divided by the annual emissions it removes. Take replacing a gas furnace with a
heat pump from the earlier example: the heat pump draws $5300\,\si{\kilo\watt
\hour}/\text{yr}$ against the furnace's $17{,}800\,\si{\kilo\watt\hour}$ of
fuel, and on a grid at $\sim 200\,\si{\gram}\,\text{CO}_2/\si{\kilo\watt\hour}$
it emits $5300\times 0.200 \approx 1.06\,\text{t}/\text{yr}$ against the
furnace's $17{,}800\times 0.200 \approx 3.56\,\text{t}/\text{yr}$, an
abatement of $\sim 2.5\,\text{t}\,\text{CO}_2/\text{yr}$. If the heat pump
costs $\$4000$ more than the furnace it replaces, at the $9.4\,\%$
capital-recovery factor that is $\sim \$376/\text{yr}$, and if it also saves
$\sim \$600/\text{yr}$ on energy (fuel bought at $\$0.06\,\si{\per\kilo\watt
\hour}$ against electricity at $\$0.15$: $17{,}800\times 0.06 - 5300\times 0.15
\approx 1068 - 795 \approx \$270/\text{yr}$ saved, plus avoided furnace
maintenance), the measure can have a negative net cost. At a net saving of
$\sim \$100/\text{yr}$ the abatement cost is $-100/2.5 \approx -\$40\,
\text{/t}\,\text{CO}_2$: the measure pays for itself and cuts carbon, a
\engterm{negative-cost} abatement. A measure with a positive net cost, say a
$\$200/\text{yr}$ premium removing the same $2.5\,\text{t}$, has an abatement
cost of $+\$80\,\text{/t}$. Ranking every available measure by its
abatement cost builds the marginal-abatement-cost curve the next mastery box
develops, and the reading here is the arithmetic behind one point on it: the
cheapest carbon to cut is the carbon whose removal also saves money, and the
abatement cost, not the raw capital, is the number that orders the measures.

% Figure: the current-voltage and power-voltage curves of a photovoltaic
% module, with the maximum-power point. The I-V curve holds nearly constant
% current up to the knee near the open-circuit voltage; the P-V product peaks
% at the maximum-power point, where a tracker holds the operating voltage.
% Values illustrative for a single module: short-circuit current 10 A,
% open-circuit voltage 40 V, so the P-V peak sits near 0.8 Voc and 0.9 Isc,
% a fill factor of about 0.75.
\begin{figure}[htbp]
\centering
\begin{tikzpicture}
\begin{axis}[
  width=12.5cm, height=7.5cm,
  xlabel={module voltage $V$ (V)},
  ylabel={current $I$ (A)},
  xmin=0, xmax=42, ymin=0, ymax=11,
  xtick={0,10,20,30,40},
  ytick={0,2,4,6,8,10},
  axis y line=left,
  axis x line=bottom,
  tick label style={font=\scriptsize},
  label style={font=\small},
  legend style={font=\scriptsize, at={(0.03,0.03)}, anchor=south west, draw=none},
]
% I-V curve: nearly flat at Isc = 10 A, knee near 32 V, falling to Voc = 40 V.
% Model I(V) = Isc * (1 - exp((V - Voc)/Vt)) with Vt chosen for a sharp knee.
\addplot[engblue, very thick, domain=0:40, samples=200]
  {10*(1 - exp((x-40)/3.2))};
\addlegendentry{$I$--$V$ curve}
% maximum-power point marker on the I-V curve near V = 32 V, I = 9.0 A
\addplot[engred, only marks, mark=*, mark size=2.4pt] coordinates {(32,9.0)};
\node[anchor=south east, font=\scriptsize, color=engred] at (axis cs:31.5,9.2) {MPP};
\end{axis}
% second axis for the power curve, scaled to share the frame
\begin{axis}[
  width=12.5cm, height=7.5cm,
  xmin=0, xmax=42, ymin=0, ymax=330,
  xtick=\empty,
  axis y line=right,
  axis x line=none,
  ylabel={power $P=VI$ (W)},
  ytick={0,100,200,300},
  tick label style={font=\scriptsize},
  label style={font=\small},
  legend style={font=\scriptsize, at={(0.03,0.20)}, anchor=south west, draw=none},
]
% P-V curve: V times I(V), peaking at the MPP near 288 W
\addplot[engorange, very thick, dashed, domain=0:40, samples=200]
  {x*10*(1 - exp((x-40)/3.2))};
\addlegendentry{$P$--$V$ curve}
\addplot[engred, only marks, mark=*, mark size=2.4pt] coordinates {(32,288)};
\end{axis}
\end{tikzpicture}
\caption[Current-voltage and power-voltage curves of a PV module]{The
current-voltage and power-voltage curves of a photovoltaic module. The
current holds near the short-circuit value $I_{\text{sc}}$ across most of the
voltage range and collapses through a knee toward the open-circuit voltage
$V_{\text{oc}}$, where the current falls to zero. The power, the product
$P=VI$ (orange, right axis), rises with voltage, peaks at the maximum-power
point (MPP), and falls to zero at both ends. A maximum-power-point tracker
holds the operating voltage at the peak as irradiance and temperature move
it. The ratio of the peak power to the product $V_{\text{oc}}I_{\text{sc}}$
is the fill factor, here about $0.75$. Values illustrative for a single
module with $I_{\text{sc}}=10\,\text{A}$ and $V_{\text{oc}}=40\,\text{V}$.}
\label{fig:vol04ch14:pv-iv-curve}
\end{figure}


% Figure: photovoltaic module output against cell temperature, showing the
% negative temperature coefficient of power. A silicon module loses about
% 0.4 percent of its power per kelvin above the 25 C rating, so a hot roof at
% 60 C delivers roughly 14 percent less than the nameplate suggests. The
% relative-power line falls linearly with cell temperature; a second marker
% shows the standard-test-condition point at 25 C.
\begin{figure}[htbp]
\centering
\begin{tikzpicture}
\begin{axis}[
  width=12.5cm, height=7.2cm,
  xlabel={cell temperature $T_{\text{cell}}$ (${}^{\circ}\mathrm{C}$)},
  ylabel={power relative to rating (\%)},
  xmin=0, xmax=75, ymin=78, ymax=112,
  xtick={0,15,25,40,55,70},
  ytick={80,85,90,95,100,105,110},
  axis lines=left,
  tick label style={font=\scriptsize},
  label style={font=\small},
  legend style={font=\scriptsize, at={(0.03,0.03)}, anchor=south west, draw=none},
]
% relative power P/Prated = 1 - 0.004*(T - 25), in percent
\addplot[engblue, very thick, domain=0:75, samples=76] {100 - 0.4*(x-25)};
\addlegendentry{$-0.4\,\%/\text{K}$ silicon}
% standard-test-condition point at 25 C, 100 percent
\addplot[englime, only marks, mark=*, mark size=2.5pt] coordinates {(25,100)};
\node[anchor=south west, font=\scriptsize, color=englime] at (axis cs:26,100.5) {STC ($25\,{}^{\circ}\mathrm{C}$)};
% hot-roof operating point at 60 C
\addplot[engred, only marks, mark=*, mark size=2.5pt] coordinates {(60,86)};
\node[anchor=north east, font=\scriptsize, color=engred] at (axis cs:59,85.5) {hot roof $60\,{}^{\circ}\mathrm{C}$};
\end{axis}
\end{tikzpicture}
\caption[PV output against cell temperature]{Photovoltaic module output
against cell temperature. Silicon has a negative temperature coefficient of
power near $-0.4\,\%$ per kelvin, so a module hotter than its
$25\,{}^{\circ}\mathrm{C}$ standard-test-condition rating delivers less than
its nameplate. A cell at $60\,{}^{\circ}\mathrm{C}$, a common summer roof
temperature, loses about $14\,\%$ of its rated power. The effect is why real
capacity factors sit below the value the peak rating and the annual
insolation would predict, why module rating is quoted at a fixed cell
temperature, and why array ventilation and cooler climates lift the delivered
energy per installed watt. Values illustrative for a crystalline-silicon
module.}
\label{fig:vol04ch14:pv-temperature-derate}
\end{figure}


% Figure: how variable renewables reshape the load-duration curve. The gross
% load-duration curve (the year's hourly demands re-sorted high to low) is a
% smooth monotone fall. Subtracting variable renewable output hour by hour and
% re-sorting gives the net-load-duration curve, which is hollowed in the
% middle, steepened at the top, and pushed below zero at the right, where the
% renewable surplus exceeds demand and must be curtailed or stored. The
% dispatchable fleet is built to the net curve, not the gross one. Both
% curves in per-unit of the annual peak against the fraction of hours.
\begin{figure}[htbp]
\centering
\begin{tikzpicture}
\begin{axis}[
  width=12.5cm, height=7.5cm,
  xlabel={fraction of hours in the year (sorted)},
  ylabel={load (per-unit of gross peak)},
  xmin=0, xmax=1, ymin=-0.35, ymax=1.05,
  xtick={0,0.2,0.4,0.6,0.8,1.0},
  ytick={-0.2,0,0.2,0.4,0.6,0.8,1.0},
  axis lines=left,
  tick label style={font=\scriptsize},
  label style={font=\small},
  legend style={font=\scriptsize, at={(0.97,0.97)}, anchor=north east, draw=none},
]
% gross load-duration curve: smooth monotone fall from 1.0 to about 0.35
\addplot[engblue, very thick, domain=0:1, samples=120]
  {0.35 + 0.65*(1-x)^1.6};
\addlegendentry{gross load duration}
% net load-duration curve: subtract renewable output and re-sort; hollowed
% middle, steeper top, dips negative at the right. Modelled as a steeper
% monotone that crosses zero near x = 0.85.
\addplot[engorange, very thick, dashed, domain=0:1, samples=120]
  {0.90*(1-x)^2.3 - 0.15};
\addlegendentry{net load duration}
% zero line
\addplot[enggray, thin, domain=0:1, samples=2] {0};
% shade the region below zero as curtailable surplus (marker text)
\node[anchor=south east, font=\scriptsize, color=engorange] at (axis cs:0.99,-0.32) {surplus (curtail/store)};
\node[anchor=north west, font=\scriptsize, color=engblue] at (axis cs:0.02,0.55) {firm-capacity need at the peak};
\end{axis}
\end{tikzpicture}
\caption[Gross and net load-duration curves]{How variable renewables reshape
the load-duration curve. The gross curve (blue) is the year's hourly demands
re-sorted from the highest hour to the lowest, a smooth monotone fall whose
left edge is the annual peak and whose area is the annual energy. Subtracting
the hourly renewable output and re-sorting gives the net-load-duration curve
(orange dashed): its middle is hollowed as midday renewable output meets
demand, its top-left is steeper because the peak still arrives when the
resource is low, and its right end dips below zero, where renewable output
exceeds demand and the surplus must be curtailed or stored. The dispatchable
fleet is built to the net curve, not the gross one, so high renewable
penetration cuts the energy the fleet must supply while leaving, and even
sharpening, the firm capacity it must hold for the peak.}
\label{fig:vol04ch14:net-load-duration}
\end{figure}


% Figure: lifecycle carbon intensity of electricity by generation source,
% in grams of CO2-equivalent per kilowatt-hour, spanning the construction,
% fuel, operation, and decommissioning of each. Coal and gas dominate; the
% low-carbon sources cluster two orders of magnitude below, their small
% residue coming from manufacturing and construction rather than operation.
% Values are representative median figures (current as of 2024); the point is
% the two-order-of-magnitude gap, not the exact bar heights.
\begin{figure}[htbp]
\centering
\begin{tikzpicture}
\begin{axis}[
  width=13cm, height=7.5cm,
  xbar,
  xlabel={lifecycle carbon intensity ($\si{\gram}\,\text{CO}_2\text{eq}/\si{\kilo\watt\hour}$)},
  xmin=0, xmax=980,
  xtick={0,200,400,600,800},
  symbolic y coords={hydro,wind,nuclear,solar PV,geothermal,gas CCGT,coal},
  ytick=data,
  y dir=reverse,
  nodes near coords,
  nodes near coords style={font=\scriptsize},
  every node near coord/.append style={anchor=west},
  bar width=8pt,
  tick label style={font=\scriptsize},
  label style={font=\small},
  enlarge y limits=0.12,
]
\addplot[fill=engblue, draw=engblue] coordinates {
  (24,hydro) (11,wind) (12,nuclear) (43,solar PV) (38,geothermal) (490,gas CCGT) (820,coal)
};
\end{axis}
\end{tikzpicture}
\caption[Lifecycle carbon intensity of electricity by source]{Lifecycle
carbon intensity of electricity by generation source, in grams of
carbon-dioxide-equivalent per kilowatt-hour, counting construction, fuel,
operation, and decommissioning. Coal and combined-cycle gas stand two orders
of magnitude above the low-carbon sources, whose small residue comes from
manufacturing and construction (the concrete and steel of a dam, the silicon
of a panel, the towers of a wind farm) rather than from operation. The gap,
not the exact bar height, is the content: decarbonising the electricity
carrier means moving generation from the top two bars to the cluster at the
bottom. Representative median figures, current as of 2024.}
\label{fig:vol04ch14:lifecycle-carbon}
\end{figure}


\begin{estimation}
\textbf{Question.} How much land would the United States need to cover
with solar panels to supply its entire electricity demand, ignoring
storage and transmission?

\textbf{Working.} US electricity demand is $\sim 4000\,\text{TWh}/
\text{yr}$ (current as of 2024). Average insolation over the
contiguous US is $\sim 200\,\si{\watt\per\meter\squared}$ averaged over
day and night and over the year; a utility solar farm converts at
$\sim 20\,\%$ panel efficiency but packs panels at $\sim 40\,\%$ ground
coverage and loses $\sim 15\,\%$ in conversion and downtime, so the
delivered areal power is roughly $200 \times 0.20 \times 0.40 \times
0.85 \approx 14\,\si{\watt\per\meter\squared}$. The annual energy per
square metre is $14\,\si{\watt} \times 8760\,\si{\hour} \approx 120\,
\si{\kilo\watt\hour}/\si{\meter\squared}$. The area is
\[
A \approx \frac{4000 \times 10^{9}\,\si{\kilo\watt\hour}}
{120\,\si{\kilo\watt\hour}/\si{\meter\squared}}
\approx 3.3 \times 10^{10}\,\si{\meter\squared}
= 3.3 \times 10^{4}\,\si{\kilo\meter\squared},
\]
a square about $180\,\si{\kilo\meter}$ on a side, around $0.4\,\%$ of
the contiguous US land area. The order-of-magnitude verdict, that the
land is available, holds under every reasonable assumption; the binding
constraints are storage (the calculation ignored night and winter),
transmission, and the spread of the resource, not the raw area. The
estimate also shows why ``solar cannot power a country because it would
cover the country'' is wrong by two orders of magnitude.
\end{estimation}

\section{The mathematics of grid operation}\pathtag{standard}

The operational quantities used throughout the chapter rest on four
derivations that the preceding sections stated and used without proof. We
give them here at working depth: the power-flow equations that relate bus
voltages to the power they exchange, the discounted-cash-flow construction
behind the levelized cost, the swing equation and droop control that govern
frequency, and the capacity-outage table that turns unit reliability into a
loss-of-load probability. Each derivation ends in a number the rest of the
chapter has already leaned on.

\subsection{The two-bus power-flow equations and the DC approximation}

Consider the simplest non-trivial network: two buses joined by a single
transmission line. Bus~1 holds voltage phasor $V_1\angle\delta_1$ and bus~2
holds $V_2\angle\delta_2$, and the line between them has series impedance
$Z = R + jX$, with admittance $Y = 1/Z = G + jB$. The current from bus~1
into the line is $I_{12} = Y(V_1\angle\delta_1 - V_2\angle\delta_2)$, and the
complex power injected at bus~1 is $S_{12} = V_1\angle\delta_1 \cdot
I_{12}^{\ast}$. Writing the angle difference as $\delta = \delta_1 -
\delta_2$ and expanding the conjugate product gives the real and reactive
power leaving bus~1:
\begin{align}
P_{12} &= V_1^2 G - V_1 V_2\,(G\cos\delta + B\sin\delta),\\
Q_{12} &= -V_1^2 B - V_1 V_2\,(G\sin\delta - B\cos\delta).
\end{align}
A transmission line is dominated by its reactance, $X \gg R$, so to a good
first approximation $G \to 0$ and $B \to -1/X$ (the susceptance of a pure
reactance is negative). The real-power flow collapses to the form the
chapter has used since the grid section:
\[
P_{12} = \frac{V_1 V_2}{X}\sin\delta,
\]
the \engterm{power-angle equation}. Real power flows from the leading bus to
the lagging bus, driven by the voltage-angle difference $\delta$, and the
transfer reaches a maximum of $P_{\max} = V_1 V_2 / X$ at $\delta = 90^\circ$
(Figure \ref{fig:vol04ch14:power-angle}). Beyond $90^\circ$ the transfer
falls with rising angle, so a load on the falling branch is statically
unstable: a small increase in demanded power lowers the transfer, the angle
runs away, and the machines slip out of synchronism. The usable operating
range is the rising branch, and the gap between the working angle and
$90^\circ$ is the \engterm{static stability margin}. The reactive-power
equation, with the same approximation, becomes $Q_{12} = (V_1^2 - V_1 V_2
\cos\delta)/X$, which for small $\delta$ is $\approx V_1(V_1 - V_2)/X$:
reactive power flows from the higher-voltage bus to the lower-voltage one and
is the lever that controls voltage magnitude, while real power follows the
angle. The decoupling of the two, real power with angle and reactive power
with voltage, is the working simplification that makes large-grid analysis
tractable.

% Figure: the power-angle curve of a two-bus AC link. Real power transfer
% P = (V1 V2 / X) sin(delta) rises with the voltage-angle difference delta
% to a maximum at 90 degrees and falls beyond, the static stability limit.
% A horizontal operating-power line cuts the curve at a stable point on the
% rising branch and an unstable point on the falling branch. Values
% illustrative (per-unit), with V1 = V2 = 1.0 pu and reactance X = 0.2 pu so
% the peak transfer is 1/0.2 = 5 pu.
\begin{figure}[htbp]
\centering
\begin{tikzpicture}
\begin{axis}[
  width=12.5cm, height=7.5cm,
  xlabel={voltage-angle difference $\delta$ (degrees)},
  ylabel={real power transfer $P$ (pu)},
  xmin=0, xmax=180, ymin=0, ymax=5.6,
  xtick={0,30,60,90,120,150,180},
  ytick={0,1,2,3,4,5},
  axis lines=left,
  tick label style={font=\scriptsize},
  label style={font=\small},
  legend style={font=\scriptsize, at={(0.97,0.03)}, anchor=south east, draw=none},
]
% power-angle curve P = 5 sin(delta)
\addplot[engblue, very thick, domain=0:180, samples=181] {5*sin(x)};
\addlegendentry{$P=(V_1V_2/X)\sin\delta$}
% operating power line at P = 3.5 pu
\addplot[engred, thick, dashed, domain=0:180, samples=2] {3.5};
\addlegendentry{operating power $=3.5$ pu}
% stable operating point: 5 sin(delta) = 3.5 -> delta = 44.4 deg
\addplot[engorange, only marks, mark=*, mark size=2.2pt] coordinates {(44.4,3.5)};
\node[anchor=west, font=\scriptsize, color=engorange] at (axis cs:46,3.0) {stable};
% unstable operating point: delta = 135.6 deg
\addplot[enggray, only marks, mark=o, mark size=2.2pt] coordinates {(135.6,3.5)};
\node[anchor=east, font=\scriptsize, color=enggray] at (axis cs:133,3.0) {unstable};
% peak / static stability limit at 90 deg
\addplot[englime, only marks, mark=triangle*, mark size=2.5pt] coordinates {(90,5)};
\node[anchor=south, font=\scriptsize, color=englime] at (axis cs:90,5.05) {limit $P_{\max}=V_1V_2/X$};
\end{axis}
\end{tikzpicture}
\caption[The power-angle curve of a two-bus AC link]{The power-angle curve
of a lossless two-bus AC link: real power transfer follows
$P=(V_1V_2/X)\sin\delta$ in the voltage-angle difference $\delta$ across the
connecting reactance $X$. Transfer rises to a maximum at $\delta=90^\circ$,
the static stability limit $P_{\max}=V_1V_2/X$, and falls beyond it. A given
operating power (red) is met at two angles: a stable point on the rising
branch, where a small increase in load draws a restoring increase in
transfer, and an unstable point on the falling branch, where the same
increase reduces transfer and the machines lose synchronism. Operating
margin is the gap between the working angle and $90^\circ$. Per-unit values
illustrative, with $V_1=V_2=1.0$ pu and $X=0.2$ pu.}
\label{fig:vol04ch14:power-angle}
\end{figure}


The \engterm{DC power-flow approximation} pushes the simplification one step
further for planning-scale studies. Assume every bus voltage holds at its
nominal $1.0$ per unit ($V_1 = V_2 = 1$), neglect the line resistance, and
take the angle differences small enough that $\sin\delta \approx \delta$ in
radians. The power-angle equation linearises to
\[
P_{12} \approx \frac{\delta_1 - \delta_2}{X},
\]
the exact electrical analogue of Ohm's law with angle in the role of voltage,
power in the role of current, and reactance in the role of resistance. The
whole network's flows then follow from a linear system $\mathbf{P} =
\mathbf{B}\,\boldsymbol{\delta}$, with $\mathbf{B}$ the susceptance matrix
built the way a conductance matrix is built for a resistive circuit. The DC
approximation drops voltage magnitude, reactive power, and losses, so it is
wrong in detail, but it is fast, it never fails to converge, and it captures
the quantity a planner most needs: which lines carry how much real power
under a given dispatch, and which will overload first. The full AC power flow,
solved by Newton-Raphson iteration on the nonlinear equations above, is the
tool when voltage and reactive power matter; the DC flow is the tool when the
question is contingency screening across thousands of cases.

\paragraph{Worked example: two-bus transfer and the stability margin.} A
$220\,\si{\kilo\volt}$ line of reactance $X = 40\,\si{\ohm}$ connects a
generator bus held at $V_1 = 1.02$ per unit to a load bus at $V_2 = 0.98$ per
unit, with a measured angle difference of $\delta = 18^\circ$. On a
$100\,\si{\mega\watt}$, $220\,\si{\kilo\volt}$ base the per-unit reactance is
$X_{\text{pu}} = X / Z_{\text{base}} = 40 / (220^2/100) = 40/484 \approx
0.0826$ per unit. The real-power transfer is $P = (V_1 V_2 / X_{\text{pu}})
\sin\delta = (1.02 \times 0.98 / 0.0826)\sin 18^\circ \approx 12.1 \times
0.309 \approx 3.74$ per unit, that is $\sim 374\,\si{\mega\watt}$. The maximum
transfer is $P_{\max} = V_1 V_2 / X_{\text{pu}} \approx 12.1$ per unit, $\sim
1210\,\si{\mega\watt}$, so the line runs at $374/1210 \approx 31\,\%$ of its
static stability limit, a comfortable margin. Halving the voltages to $V_1 =
V_2 = 0.90$ per unit (a depressed-voltage emergency) would cut $P_{\max}$ to
$0.81/0.0826 \approx 9.8$ per unit and force the angle to $\delta =
\arcsin(3.74/9.8) \approx 22.5^\circ$ to carry the same power, eating into the
margin; this is why voltage collapse and angle instability arrive together,
and why reactive support that holds voltage up also buys angle-stability
margin.

\subsection{Levelized cost of energy from discounted cash flow}

The chapter has quoted levelized costs and used a capital-recovery factor
without deriving either. Both follow from one idea: money has a time value,
so a cost or a revenue that arrives in year $t$ is worth less today than the
same sum now, discounted by $(1+i)^{-t}$ at a discount rate $i$. The
\engterm{levelized cost of energy} is the constant per-unit price $\text{LCOE}$
that makes the discounted lifetime revenue equal the discounted lifetime cost.
Writing the year-$t$ cost as a capital outlay $I_t$ plus operations,
maintenance, and fuel $M_t$, and the year-$t$ energy as $E_t$,
\[
\sum_{t=0}^{N} \frac{\text{LCOE}\cdot E_t}{(1+i)^t}
= \sum_{t=0}^{N} \frac{I_t + M_t}{(1+i)^t},
\qquad\text{so}\qquad
\text{LCOE} = \frac{\displaystyle\sum_{t=0}^{N} (I_t + M_t)/(1+i)^t}
{\displaystyle\sum_{t=0}^{N} E_t/(1+i)^t}.
\]
The discounting of the denominator is the part that surprises: a plant's
future energy is discounted at the same rate as its future cost, because a
megawatt-hour delivered in year twenty is worth, in present terms, far less
than one delivered in year one. For the common case of a single upfront
capital cost $I_0$, a constant annual running cost $M$, and a constant annual
energy $E$ over $N$ years, the denominators are geometric series. The sum
$\sum_{t=1}^{N}(1+i)^{-t} = [1 - (1+i)^{-N}]/i$ is the annuity factor, and its
reciprocal is the \engterm{capital-recovery factor}
\[
\text{CRF} = \frac{i}{1 - (1+i)^{-N}},
\]
the constant annual payment that repays one unit of capital over $N$ years at
rate $i$. The levelized cost then reduces to
\[
\text{LCOE} = \frac{I_0 \cdot \text{CRF} + M}{E},
\]
the annualised capital plus the running cost, divided by the annual energy,
which is exactly the bookkeeping the case studies used. The capital-recovery
factor carries the entire dependence on financing: at $i = 0.07$ and $N = 20$,
$\text{CRF} = 0.07/(1 - 1.07^{-20}) \approx 0.0944\,\text{yr}^{-1}$; at $i =
0.03$ the same plant has $\text{CRF} \approx 0.0672$, almost a third lower, so
a capital-heavy low-carbon plant is far cheaper under cheap capital than under
dear capital. The discount rate is therefore not a financial footnote; it is
the single number that most moves the relative cost of capital-heavy
renewables and nuclear against fuel-heavy gas.

\paragraph{Worked example: how the discount rate reorders the merit of
generation.} Compare a solar plant ($I_0 = \$900\,\si{\per\kilo\watt}$,
$M = \$15\,\si{\per\kilo\watt}$ per year, no fuel, $\text{CF} = 0.18$,
$N = 30$) with a combined-cycle gas plant ($I_0 = \$1000\,\si{\per\kilo\watt}$,
$M + \text{fuel} = \$45\,\si{\per\mega\watt\hour}$ delivered, $\text{CF} =
0.50$, $N = 30$). The annual energy per kilowatt is $\text{CF}\times 8760$:
$1577\,\si{\kilo\watt\hour}$ for solar, $4380\,\si{\kilo\watt\hour}$ for gas.
At a $7\,\%$ discount rate, $\text{CRF}_{30} = 0.07/(1 - 1.07^{-30}) \approx
0.0806\,\text{yr}^{-1}$. Solar: $\text{LCOE} = (900\times 0.0806 +
15)/1.577 \approx (72.5 + 15)/1.577 \approx \$55\,\si{\per\mega\watt\hour}$.
Gas: $\text{LCOE} = (1000\times 0.0806)/4.380 + 45 \approx 18.4 + 45 \approx
\$63\,\si{\per\mega\watt\hour}$. Solar wins by a small margin. Now drop the
discount rate to $3\,\%$: $\text{CRF}_{30} \approx 0.051$. Solar falls to
$(900\times 0.051 + 15)/1.577 \approx \$38\,\si{\per\mega\watt\hour}$ while
gas falls only to $(1000\times 0.051)/4.380 + 45 \approx \$57\,\si{\per\mega
\watt\hour}$, because gas is dominated by the undiscounted fuel cost that no
financing rate touches. Cheap capital widens solar's lead from $\$8$ to $\$19
\,\si{\per\mega\watt\hour}$; dear capital ($i = 12\,\%$, $\text{CRF}_{30}
\approx 0.124$) reverses it, lifting solar to $(900\times 0.124 + 15)/1.577
\approx \$80$ against gas at $(1000\times 0.124)/4.380 + 45 \approx \$73$. The
levelized cost is not a property of a technology; it is a property of a
technology at a financing rate, and the same two plants change places as the
cost of capital moves.

\begin{mastery}
\textbf{The learning curve: cost falls with cumulative production, not with
time.} The dramatic fall in solar and battery cost that recurs through this
chapter follows a regularity that predates both, \engterm{Wright's law}: the
unit cost of a manufactured good falls by a constant fraction for every
doubling of cumulative production. Writing $C(x)$ for the cost at cumulative
volume $x$, the law is a power law $C(x) = C_0\,(x/x_0)^{-b}$, and the
\engterm{learning rate} is the fractional cost drop per doubling, $\text{LR} =
1 - 2^{-b}$. Solar photovoltaics have held a learning rate near $\sim 20\,\%$
for four decades: every doubling of the panels ever made has cut the price by
about a fifth, and because the cumulative volume has doubled many times the
compounded fall is the $\sim 90\,\%$ decline since 2010 the chapter quotes.
Lithium-ion cells sit near a $\sim 18\,\%$ learning rate, wind turbines lower
at $\sim 10$ to $15\,\%$, and mature thermal and nuclear plants show little or
negative learning, because their cost is dominated by site-specific
construction and regulation rather than by a mass-manufactured unit that
accumulates volume. The law reframes the transition as a positive feedback:
deploying a technology lowers its cost, which drives more deployment, which
lowers the cost again, so a policy that buys down the early volume of a
high-learning-rate technology can move it along its curve to a cost that no
longer needs support, which is the mechanism behind the feed-in tariffs and
deployment mandates that carried solar and batteries down their curves.
The limit is that the law is empirical, not guaranteed: it holds while the
technology keeps finding manufacturing and materials improvements and can
stall when a raw-material cost floor (the lithium or the silver in a cell) or
a physical efficiency ceiling (the Shockley-Queisser limit for a single-
junction cell) is reached, so the extrapolation is a guide, not a promise, and
the durable lesson is that cost is bought with cumulative volume, which is why
the technologies that scale fastest also cheapen fastest.
\end{mastery}

\subsection{Primary frequency control: the swing equation and droop}

The frequency dynamics the chapter sketched in the grid section follow from
Newton's second law applied to the rotating mass of the synchronous fleet.
Let $J$ be the combined moment of inertia of all generators and motors locked
to the grid, $\omega$ the electrical angular frequency, $T_m$ the mechanical
torque driving the rotors, and $T_e$ the electrical torque the load draws.
The rotor obeys $J\,d\omega/dt = T_m - T_e$. Multiplying by $\omega$ converts
torque to power, $J\omega\,d\omega/dt = P_m - P_e$, and normalising by the
system rating $S$ and the nominal frequency $\omega_0$ defines the
\engterm{inertia constant} $H = \tfrac12 J\omega_0^2 / S$, the stored kinetic
energy in seconds of rated output. In terms of the per-unit frequency
deviation $\Delta f = (f - f_0)/f_0$, the result is the \engterm{swing
equation}:
\[
2H\,\frac{d(\Delta f)}{dt} = \Delta P_m - \Delta P_e - D\,\Delta f,
\]
where $D$ is a small load-damping term (loads draw slightly less power as
frequency falls). The equation says the rate of change of frequency is the
power imbalance divided by twice the inertia. At the instant a block $\Delta
P$ of generation is lost, before any governor has moved, $\Delta P_m = 0$ and
the initial rate of change of frequency is
\[
\left.\frac{d(\Delta f)}{dt}\right|_{0} = -\frac{\Delta P}{2H},
\qquad\text{or in hertz}\qquad
\left.\frac{df}{dt}\right|_{0} = -\frac{f_0\,\Delta P}{2H S},
\]
the expression the inertia mastery box quoted. Halving the inertia doubles the
initial slope and deepens the nadir, the regime a high-inverter grid moves
into.

The governor closes the loop. Each generator's mechanical power is commanded
to rise as frequency falls, in proportion set by the \engterm{droop} $R$:
$\Delta P_m = -\Delta f / R$, where $R$ is the per-unit frequency change that
swings the unit from no load to full load, conventionally $\sim 0.05$ (a
$5\,\%$ droop means full frequency authority over a $5\,\%$, that is $2.5\,
\si{\hertz}$ on a $50\,\si{\hertz}$ grid, frequency swing). Substituting the
governor law into the swing equation and solving for the steady state, where
$d(\Delta f)/dt = 0$, gives the settled frequency deviation after a sustained
imbalance $\Delta P$:
\[
\Delta f_{\text{ss}} = -\frac{\Delta P}{1/R + D} \approx -R\,\Delta P
\quad(\text{since } 1/R \gg D),
\]
so a $5\,\%$ droop holds the steady-state frequency error to $R$ times the
per-unit power deficit. With $n$ units of droops $R_k$ sharing the load, the
common settled frequency is read by all of them, and each picks up power
$\Delta P_k = -\Delta f_{\text{ss}}/R_k$, so the units share the imbalance in
inverse proportion to their droops (Figure
\ref{fig:vol04ch14:droop-sharing}). Droop control is therefore a
communication-free load-sharing scheme: every machine measures the one global
signal, frequency, and contributes its share automatically. Secondary control
(automatic generation control) then trims the residual steady-state error back
to zero over minutes by shifting the units' no-load setpoints, the slow
integral action on top of the fast proportional droop.

% Figure: governor droop characteristics of two synchronous units sharing a
% load increase. Each unit's frequency-power line has a negative slope set by
% its droop R; at a common settled frequency the units pick up power in
% inverse proportion to their droop. A load increase drops the common
% frequency from f0 to f1, and the two units move along their lines to new
% outputs. Frequency in Hz on the vertical axis (inverted sense: rising power
% lowers frequency), power in per-unit of each unit's rating on the
% horizontal. Values illustrative: unit A droop 4 percent, unit B droop 6
% percent, nominal 50 Hz.
\begin{figure}[htbp]
\centering
\begin{tikzpicture}
\begin{axis}[
  width=12.5cm, height=7.5cm,
  xlabel={unit output (pu of unit rating)},
  ylabel={grid frequency (Hz)},
  xmin=0, xmax=1.0, ymin=49.4, ymax=50.1,
  xtick={0,0.2,0.4,0.6,0.8,1.0},
  ytick={49.4,49.6,49.8,50.0},
  axis lines=left,
  tick label style={font=\scriptsize},
  label style={font=\small},
  legend style={font=\scriptsize, at={(0.03,0.03)}, anchor=south west, draw=none},
]
% Unit A: droop 4 percent. f = 50 - 0.04*50*P = 50 - 2P
\addplot[engblue, very thick, domain=0:1.0, samples=2] {50 - 2*x};
\addlegendentry{unit A, droop $4\,\%$}
% Unit B: droop 6 percent. f = 50 - 0.06*50*P = 50 - 3P
\addplot[engamber, very thick, dashed, domain=0:1.0, samples=2] {50 - 3*x};
\addlegendentry{unit B, droop $6\,\%$}
% common frequency before load step: f0 = 49.9 -> A at 0.05, B at 0.0333
\draw[enggray, thin, dash dot] (axis cs:0,49.9) -- (axis cs:1.0,49.9);
\node[anchor=west, font=\scriptsize, color=enggray] at (axis cs:0.62,49.93) {$f_0=49.90$};
% common frequency after load step: f1 = 49.7 -> A at 0.15, B at 0.10
\draw[engred, thin, dash dot] (axis cs:0,49.7) -- (axis cs:1.0,49.7);
\node[anchor=west, font=\scriptsize, color=engred] at (axis cs:0.62,49.73) {$f_1=49.70$};
% operating points after the step
\addplot[engblue, only marks, mark=*, mark size=2pt] coordinates {(0.15,49.7)};
\addplot[engamber, only marks, mark=square*, mark size=2pt] coordinates {(0.10,49.7)};
\end{axis}
\end{tikzpicture}
\caption[Governor droop characteristics and load sharing]{Governor droop
characteristics of two synchronous units. Each unit's settled output rises
as the common grid frequency falls, along a line whose slope is set by the
droop setting $R$ (the per-unit frequency change that moves the unit from no
load to full load). A load increase pulls the common frequency from $f_0$
down to $f_1$, and the two units slide along their lines to pick up the
extra power in inverse proportion to their droop: the stiffer unit A (lower
droop) takes the larger share. The common settled frequency, not the
individual outputs, is what every unit on the grid reads, so droop control
shares load among many machines with no communication between them. Values
illustrative; nominal $50\,\si{\hertz}$.}
\label{fig:vol04ch14:droop-sharing}
\end{figure}


\paragraph{Worked example: nadir and droop sharing after a trip.} A
$30\,\si{\giga\watt}$ grid with an inertia constant $H = 5\,\si{\second}$
loses a $900\,\si{\mega\watt}$ unit. The per-unit deficit is $\Delta P =
900/30000 = 0.03$. The initial rate of frequency fall is $df/dt = -f_0 \Delta
P/(2H) = -50 \times 0.03/(2\times 5) = -0.15\,\si{\hertz}/\si{\second}$, so in
the first second, before governors respond, the frequency drops $\sim 0.15\,
\si{\hertz}$. The surviving fleet carries an effective droop of $R = 0.05$, so
the steady-state deviation is $\Delta f_{\text{ss}} \approx -R\,\Delta P =
-0.05\times 0.03 = -0.0015$ per unit, that is $-0.075\,\si{\hertz}$: the grid
settles near $49.93\,\si{\hertz}$, inside the emergency band, and secondary
control restores nominal. If the inertia were halved to $H = 2.5\,\si{\second}$
(a high-renewable hour), the initial slope would double to $-0.30\,\si{\hertz}/
\si{\second}$ and the nadir would reach deeper before the governors arrest it,
which is the case for synthetic inertia and faster primary response. The
steady-state deviation does not change with inertia, since it depends only on
droop, but the transient nadir does, and the nadir is what trips automatic
load shedding. Two units of droops $R_A = 0.04$ and $R_B = 0.06$ carrying the
$0.03$ deficit at a common settled frequency share it through their droops:
unit A picks up $\Delta P_A = -\Delta f_{\text{ss}}/R_A$ and unit B $\Delta
P_B = -\Delta f_{\text{ss}}/R_B$, in the ratio $R_B : R_A = 0.06 : 0.04 = 3 :
2$, so the stiffer unit A carries $60\,\%$ of the shortfall and unit B
$40\,\%$. The droop setting, not the unit size, fixes the share.

\subsection{Reliability: the capacity-outage table and loss-of-load
probability}

The reserve margin the case studies held at a planning target rests on a
probabilistic construction. Each generating unit is either available or
forced out, and its \engterm{forced-outage rate} (FOR) is the long-run
probability it is unavailable when called, typically a few percent for a
well-maintained thermal unit. The system's reliability is the probability
that enough units are available to meet the load. For $n$ identical units each
of capacity $C$ and forced-outage rate $q$, the number forced out is binomial:
the probability that exactly $k$ units are out is $\binom{n}{k} q^k
(1-q)^{n-k}$, and the \engterm{capacity-outage probability table} (COPT) lists
the cumulative probability that the capacity on outage equals or exceeds each
level (Figure \ref{fig:vol04ch14:copt-lolp}). The
\engterm{loss-of-load probability} (LOLP) is read off the table at the deficit
between installed capacity and load: if the installed capacity is $n C$ and the
peak load is $L$, the load is lost whenever more than $nC - L$ of capacity is
forced out, so
\[
\text{LOLP} = \Pr\!\left(\text{capacity on outage} > nC - L\right)
= \sum_{k > (nC-L)/C} \binom{n}{k} q^k (1-q)^{n-k}.
\]
Summed over the hours of a year against the load-duration curve, the same
construction gives the \engterm{loss-of-load expectation} (LOLE) in hours per
year, the planning currency the case studies held at $\sim 0.1\,\si{\hour}/
\text{yr}$. The reserve margin is chosen so the LOLP sits below the target,
and the table shows why the last increment of reliability is dear: each added
unit of reserve shifts the cumulative curve down by roughly a factor of $q$
per step, so reaching a stringent target costs a whole unit of capacity that
runs only in the rare hours the rest of the fleet is short.

% Figure: the capacity-outage probability table as a cumulative curve. The
% horizontal axis is capacity on outage (MW); the vertical axis is the
% probability that the outage equals or exceeds that level, on a log scale.
% The available-generation deficit relative to a peak load sets the
% loss-of-load probability: the curve read at (installed minus load) gives
% the probability the surviving capacity cannot meet the load. Adding a unit
% shifts the curve down (more reserve). Values illustrative for a small
% five-unit system with per-unit forced-outage rate of about 5 percent.
\begin{figure}[htbp]
\centering
\begin{tikzpicture}
\begin{axis}[
  width=12.5cm, height=7.5cm,
  xlabel={capacity on outage $X$ (MW)},
  ylabel={$\Pr(\text{outage}\ge X)$},
  xmin=0, xmax=500, ymin=1e-5, ymax=1.2,
  ymode=log,
  log basis y=10,
  xtick={0,100,200,300,400,500},
  ytick={1,0.1,0.01,0.001,0.0001,0.00001},
  yticklabels={$1$,$10^{-1}$,$10^{-2}$,$10^{-3}$,$10^{-4}$,$10^{-5}$},
  grid=both,
  grid style={enggray!25},
  tick label style={font=\scriptsize},
  label style={font=\small},
  legend style={font=\scriptsize, at={(0.97,0.97)}, anchor=north east, draw=none},
]
% five 100-MW units, FOR = 0.05: cumulative outage staircase (const plot).
% Pr(>=0)=1; Pr(>=100)=1-(0.95)^5=0.2262; Pr(>=200)=1-binom up to 1 out
% =0.0226; Pr(>=300)=0.001158; Pr(>=400)=3.05e-5; Pr(>=500)=3.125e-7.
\addplot[engblue, very thick, const plot] coordinates {
  (0,1)(100,1)(100,0.2262)(200,0.2262)(200,0.02259)(300,0.02259)(300,0.001158)(400,0.001158)(400,3.0e-5)(500,3.0e-5)
};
\addlegendentry{$5\times100$ MW, FOR $=0.05$}
% sixth unit added (reserve): curve shifts so the same deficit is rarer.
\addplot[engamber, very thick, dashed, const plot] coordinates {
  (0,1)(100,1)(100,0.2649)(200,0.2649)(200,0.03277)(300,0.03277)(300,0.00223)(400,0.00223)(400,8.64e-5)(500,8.64e-5)
};
\addlegendentry{$6\times100$ MW (added reserve)}
% load line: deficit at which load is lost (installed minus peak = 200 MW)
\draw[engred, thick, dash dot] (axis cs:200,1e-5) -- (axis cs:200,1.2);
\node[engred, font=\scriptsize, anchor=south, rotate=90] at (axis cs:200,2e-4) {deficit $=$ installed $-$ load};
\end{axis}
\end{tikzpicture}
\caption[Capacity-outage probability table and loss-of-load probability]{The
capacity-outage probability table drawn as a cumulative curve: the
probability that the capacity on outage equals or exceeds a level $X$, for a
small system of identical units each with a forced-outage rate of about
$5\,\%$ (log vertical scale). Read at the deficit between installed capacity
and peak load (red line), the curve gives the loss-of-load probability, the
chance the surviving units cannot meet the load. Adding a unit of reserve
(amber) shifts the curve down, so the same deficit becomes rarer, and the
reserve margin is chosen to push the loss-of-load probability below the
planning target. The steep staircase is why the last increment of reliability
is bought with a whole unit of capacity that runs only in the rare hours it
is needed. Values illustrative.}
\label{fig:vol04ch14:copt-lolp}
\end{figure}


\paragraph{Worked example: loss-of-load probability of a small system.} A
small grid has five $100\,\si{\mega\watt}$ units, each with a forced-outage
rate $q = 0.05$, serving a peak load of $L = 300\,\si{\mega\watt}$. Installed
capacity is $500\,\si{\mega\watt}$, so the load is lost when more than $500 -
300 = 200\,\si{\mega\watt}$, that is more than two units, are forced out. The
loss-of-load probability is the probability that three, four, or five units
are out:
\[
\text{LOLP} = \sum_{k=3}^{5} \binom{5}{k} (0.05)^k (0.95)^{5-k}
= 0.001158 + 3.0\times 10^{-5} + 3.1\times 10^{-7}
\approx 1.19\times 10^{-3}.
\]
Over $8760\,\si{\hour}$ at the peak that is $\sim 10.4\,\si{\hour}/\text{yr}$,
far above a $0.1\,\si{\hour}/\text{yr}$ target. Adding a sixth $100\,\si{\mega
\watt}$ unit raises installed capacity to $600\,\si{\mega\watt}$, so the load
is now lost only when more than three units are out:
\[
\text{LOLP} = \sum_{k=4}^{6} \binom{6}{k} (0.05)^k (0.95)^{6-k}
\approx 8.6\times 10^{-5} + 1.8\times 10^{-6} + 1.6\times 10^{-8}
\approx 8.8\times 10^{-5},
\]
a drop of more than a factor of ten from one added unit, the staircase of
Figure \ref{fig:vol04ch14:copt-lolp}. The reserve margin has risen from
$500/300 - 1 = 67\,\%$ to $100\,\%$, large for this toy system because the
units are coarse relative to the load; a real grid with many small units
reaches the same reliability at a far smaller percentage margin, which is one
reason a fleet of many modest units is more reliable per installed megawatt
than a few large ones. The construction also exposes the failure mode the
Texas 2021 event \cite{acc:ferc-nerc-texas-2021} embodied: the binomial
assumes the units fail independently, and a common-mode stress (extreme cold
de-rating many gas units at once) makes the outages correlated, which fattens
the tail of the distribution far beyond the independent-failure table and
voids the reserve margin computed from it.

\begin{mastery}
\textbf{Reactive power is the voltage currency, and it does not travel.} The
power-angle derivation split the two-bus flow into a real part carried by the
angle and a reactive part carried by the voltage difference, and the split is
the working basis of voltage control. Reactive power $Q$ is the part of the
apparent power $S = P + jQ$ that oscillates between the source and the
inductive and capacitive elements of the network without doing net work; it
energises the magnetic fields of transformers and motors and the electric
fields of lines and cables. A grid must supply it, because most load is
partly inductive (motors, transformers), and the reactive demand sets the
voltage: too little reactive supply and the voltage sags, too much and it
rises. Unlike real power, reactive power does not travel well, because moving
it along a reactive line itself consumes reactive power ($Q_{\text{loss}} =
I^2 X$, and $X$ dominates the line), so voltage must be supported locally,
near where the reactive load sits, rather than imported from a distant
generator the way real power is. The local sources are synchronous generators
and condensers (which can absorb or supply reactive power by adjusting their
field excitation), shunt capacitor banks (supply, switched in as load rises),
shunt reactors (absorb, switched in on lightly loaded lines whose charging
overvolts them), and modern power-electronic compensators (static var
compensators and STATCOMs, which adjust continuously and fast). The energy
transition complicates this: retiring synchronous generators removes their
reactive-support and voltage-regulating capability along with their inertia,
so a high-inverter grid must manufacture reactive support too, through
grid-forming inverters and dedicated compensators. Voltage collapse, the
runaway in which sagging voltage raises current, which raises reactive loss,
which sags voltage further, is the reactive-power analogue of the angle
instability the power-angle curve showed, and the two arrive together in a
stressed grid, which is why reactive reserve is planned alongside real-power
reserve.
\end{mastery}

\subsection{An applied vignette: a rooftop-solar net-metering household}

The chapter's household project models a year of dispatch; a shorter
calculation settles the question a homeowner actually asks, whether the
rooftop array pays for itself, and shows how the answer turns on the billing
arrangement rather than on the physics. Take a household at a sunny
mid-latitude site that installs a $6\,\si{\kilo\watt\textsubscript{p}}$ array
at an installed cost of $\$2.50\,\si{\per\watt}$ (current as of 2024), so
$\$15\,000$ before any subsidy, and consumes $9000\,\si{\kilo\watt\hour}/
\text{yr}$. At a $20\,\%$ capacity factor the array generates $6 \times 0.20
\times 8760 \approx 10{,}500\,\si{\kilo\watt\hour}/\text{yr}$, slightly more
than the household uses.

The payback depends on what each kilowatt-hour is worth, and that is set by the
tariff, not by the panel. Under \engterm{full net metering}, the meter runs
backward when the array exports and the household is billed only on its net
annual draw, so every generated kilowatt-hour offsets one bought kilowatt-hour
at the full retail rate of, say, $\$0.20\,\si{\per\kilo\watt\hour}$. The annual
saving is then $10{,}500 \times 0.20 \approx \$2100/\text{yr}$ (capped at the
$\$1800$ the household's own consumption is worth, with the small surplus
credited or lost depending on the rule), and the simple payback is $\sim
15{,}000/1800 \approx 8.3$ years, comfortably inside the array's $\sim
25$-year life. Under \engterm{net billing}, the more common arrangement as
distributed solar has grown (current as of 2024), exports are paid not at the
retail rate but at a lower export rate near the wholesale value, say $\$0.06\,
\si{\per\kilo\watt\hour}$, while only the self-consumed share offsets the full
retail rate. If the household self-consumes $40\,\%$ of its generation directly
(the no-battery midday overlap) and exports the other $60\,\%$, the saving is
$0.40\times 10{,}500\times 0.20 + 0.60\times 10{,}500\times 0.06 \approx 840 +
378 \approx \$1218/\text{yr}$, and the payback stretches to $\sim 15{,}000/1218
\approx 12.3$ years. The same hardware, the same sunshine, and a payback that
moves by four years on the billing rule alone.

The lever the household controls under net billing is self-consumption, which
is why a battery changes the economics only under that tariff. The
self-consumption saturation calculation earlier in the chapter gives the shape:
a battery that lifts self-consumption from $40\,\%$ to $\sim 70\,\%$ moves
$0.30\times 10{,}500 \approx 3150\,\si{\kilo\watt\hour}/\text{yr}$ from the
$\$0.06$ export rate to the $\$0.20$ retail offset, worth $3150\times(0.20 -
0.06) \approx \$440/\text{yr}$. Against a $\$6000$ battery at the
$9.4\,\%$ capital-recovery factor, that is $\$564/\text{yr}$ of annualised cost
to capture $\$440/\text{yr}$ of value, so the battery does not pay for itself
on arbitrage alone at these prices and must earn the rest from backup value or
a steeper tariff spread. Under full net metering the battery saves almost
nothing, because the grid already serves as a free, lossless, infinite store
at retail value, which is precisely why utilities have moved away from full net
metering as solar penetration has risen: the arrangement that makes the
homeowner's array pay fastest also shifts the most grid cost onto
non-solar customers, and the export rate is where that conflict is settled.
The numbers carry their year and move with hardware cost, retail rate, and the
export rule, but the structure is durable: rooftop solar is a sound investment
on physics and falling hardware cost, and the variable that decides the
payback is regulatory, the price the grid pays for an exported kilowatt-hour.

\subsection{A second applied vignette: a data-centre behind-the-meter power system}

A large computing facility is one of the fastest-growing loads on modern
grids, and its power system is a compact illustration of the chapter's
quantities running together at the scale of a single site. Take a facility
with an information-technology load of $40\,\si{\mega\watt}$, the power drawn
by the servers themselves. The site draws more than that from the grid,
because cooling, power conversion, and distribution all add overhead, and the
ratio of total facility power to IT power is the \engterm{power usage
effectiveness} (PUE). A well-run modern facility reaches a PUE of $\sim 1.2$,
so the grid draw is $40\times 1.2 = 48\,\si{\mega\watt}$, of which
$8\,\si{\mega\watt}$ is the cooling and conversion overhead and $40\,\si{
\mega\watt}$ is the useful compute. An older facility at a PUE of $1.6$ would
draw $64\,\si{\mega\watt}$ for the same compute, so the $\sim 16\,\si{\mega
\watt}$ difference is the return on the cooling and power-train engineering,
worth $\sim \$14\text{M}/\text{yr}$ at $\$100\,\si{\per\mega\watt\hour}$ and
a near-constant load, which is why cooling design is where data-centre energy
engineering concentrates.

The load shape is the feature that distinguishes a data centre from the
households and grids earlier in the chapter: it is nearly flat. The servers
run around the clock at a high and slowly varying utilisation, so the
facility's capacity factor is $\sim 0.85$ to $0.95$, far above any renewable
and above most dispatchable plant. A flat, high-capacity-factor load is the
easiest kind to serve with firm low-carbon generation and the hardest to
serve with unpaired solar, because the load does not follow the sun; the
midday solar peak matches only a third of the day, and the night and the
cloudy stretches must come from the grid, from storage, or from firm
generation. The annual energy is $48\,\si{\mega\watt}\times 0.90\times 8760
\approx 378\,\si{\giga\watt\hour}/\text{yr}$, comparable to the whole remote
community of the microgrid case but drawn by a single building, and the
matching problem is why the large computing operators contract for firm
low-carbon power (nuclear, geothermal, or round-the-clock renewable-plus-
storage) rather than for the cheaper but intermittent energy an unpaired
solar contract would supply.

The reliability requirement inverts the household's tolerance for grid
dependence. A data centre is engineered for continuity measured in the last
fractions of a percent: a ``four nines'' target is $99.99\%$ availability,
$\sim 53\,\si{\minute}$ of downtime a year, and the higher tiers push toward
``five nines'', $\sim 5\,\si{\minute}$. The site cannot ride through a grid
outage on the grid alone, so it carries its own bridging and backup: an
uninterruptible-power-supply battery that carries the full $48\,\si{\mega
\watt}$ for the seconds-to-minutes it takes the standby generators to start
and take load, and diesel or gas generators sized to carry the site through a
prolonged grid outage. The battery here is sized for power, not energy: it
must deliver the full load for perhaps $5\,\si{\minute}$, so its energy is
only $48\,\si{\mega\watt}\times (5/60)\,\si{\hour} = 4\,\si{\mega\watt\hour}$,
a $16\,\text{C}$ discharge that pays the steep resistive penalty the C-rate
worked example described, which is why the bridging battery is a different
design from the energy-optimised grid battery even when the chemistry is the
same. The redundancy is specified as $N{+}1$ or $2N$ on every element (the
same contingency logic the grid applies to generators and lines applies here
to the transformers, the UPS strings, and the generators), so a single
failure of any component never drops the load.

The site's growing role as a controllable grid resource is the integration
twist. A facility with a bridging battery, standby generators, and a workload
that can be shifted or throttled is a flexible load the grid can call on: it
can shed to its own generators during a grid emergency (the direct-load-
control mechanism at industrial scale), shift deferrable compute (batch jobs,
model training) to hours of surplus renewable output, and offer its UPS
battery for fast frequency response in the seconds after a trip. The same
attributes that make a data centre a hard load to serve, its size and its
appetite for continuity, make it, once instrumented, one of the more valuable
flexible resources on a high-renewable grid, and the operators that co-locate
compute with cheap renewable generation and time the deferrable workload to
the resource are running the sector-coupling argument of the transition at the
scale of a single interconnection point.

\section{An end-to-end case study: sizing a solar-plus-storage microgrid}\pathtag{core}

The chapter's quantities are easiest to trust when they run end to end on
one system. We size a solar-plus-storage microgrid for a small remote
community from its load profile through its generation and storage to its
levelized cost, and compare the result against the diesel generation it
would replace. The numbers are representative of a community of about
$250$ households at a sunny mid-latitude site (capacity factor $\sim
20\,\%$) and carry their year; the method transfers to any site once the
local load and irradiance data replace the illustrative figures.

\subsection{The load}

The community draws an average $42\,\si{\kilo\watt}$ with a morning rise
and a pronounced evening peak near $84\,\si{\kilo\watt}$, the shape of the
blue trace in Figure \ref{fig:vol04ch14:microgrid-week}. The annual energy
is
\[
E_{\text{load}} = 42\,\si{\kilo\watt}\times 8760\,\si{\hour}
\approx 368\,\si{\mega\watt\hour}/\text{yr}.
\]
The evening peak is the design driver: it arrives after sunset, so it must
be served from storage or from a backup generator, not directly from the
array. The ratio of peak to average, here about $2{:}1$, sets how much
discharge power the battery must deliver.

\subsection{Sizing the PV array}

The array must supply the annual energy with margin for conversion losses
and for the share that cycles through the battery (at a round-trip cost)
rather than serving load directly. Take a target where the array's annual
generation overbuilds the load by $\sim 1.3\times$, the value the
reliability calculation above identified as the knee of the
loss-of-load curve. The required AC energy is $1.3\times 368 \approx
478\,\si{\mega\watt\hour}/\text{yr}$. At a $20\,\%$ capacity factor the AC
nameplate is
\[
P_{\text{AC}} = \frac{E}{\text{CF}\times 8760}
= \frac{478\times 10^3\,\si{\kilo\watt\hour}}{0.20\times 8760\,\si{\hour}}
\approx 273\,\si{\kilo\watt}.
\]
Allowing a DC-to-AC ratio of $\sim 1.25$ (the array is oversized relative
to the inverter to fatten the shoulders of the daily curve, accepting a
little midday clipping) gives a DC nameplate of $\sim 340\,\si{\kilo\watt
\textsubscript{p}}$. We round the build to $250\,\si{\kilo\watt}$ AC /
$340\,\si{\kilo\watt\textsubscript{p}}$ DC, the array behind the amber
trace in Figure \ref{fig:vol04ch14:microgrid-week}.

\subsection{Sizing the battery for autonomy and reliability}

The battery serves the evening peak every day and carries the community
through cloudy stretches. Two requirements size it. The daily cycling
requirement is the evening energy that arrives after sunset: integrating
the load above the PV trace from sunset to sunrise gives $\sim
350\,\si{\kilo\watt\hour}$ on a typical day. The autonomy requirement,
from the reliability calculation, is $\sim 2.8\,$ days of average load at
the $1.3\times$ overbuild to reach a $99.95\,\%$ reliability target;
sizing storage for the full autonomy would demand $2.8\times 42\times 24
\approx 2.8\,\si{\mega\watt\hour}$, which is expensive. The least-cost
design does not buy full autonomy from the battery. We size the battery to
$600\,\si{\kilo\watt\hour}$ usable, enough for the daily cycle with margin
and for one cloudy day, and let a small diesel genset cover the rare
multi-day deficit. The battery power rating is set by the evening peak it
must serve from storage, $\sim 84\,\si{\kilo\watt}$ minus whatever PV
remains at dusk; a $150\,\si{\kilo\watt}$ inverter covers the peak with
headroom. The battery is therefore $600\,\si{\kilo\watt\hour} / 150\,\si{
\kilo\watt}$, a $4\,\si{\hour}$ system, the band where lithium-ion is the
cheaper storage by the crossover calculation above.

\subsection{Dispatch over a representative week}

Figure \ref{fig:vol04ch14:microgrid-week} runs the sized system over a
representative week with the greedy dispatch of
\texttt{code/household\_dispatch.py} scaled to the community: at each hour,
serve load from PV first, charge the battery from any surplus until full,
discharge the battery into any deficit until empty, and call the diesel
only when the battery reaches its floor. On the three clear days the
array refills the battery to its ceiling by early afternoon and the
evening peak is served entirely from storage, with the surplus beyond the
battery curtailed. The cloudy stretch on days four and five is the test:
the reduced array cannot refill the battery, the state of charge falls
toward its floor, and on the second cloudy evening the diesel runs for a
few hours (the red window). Summed over a representative year of such
weeks, the diesel supplies $\sim 4\,\%$ of the annual energy and the
loss-of-load expectation lands near the $99.95\,\%$ target; the PV-plus-
battery supplies the other $\sim 96\,\%$, and curtailment runs at $\sim
12\,\%$ of generation, consistent with the $1.3\times$ overbuild at
$4\,\si{\hour}$ storage in Figure
\ref{fig:vol04ch14:overbuild-curtailment}.

% Figure: representative-week dispatch for the solar-plus-storage microgrid
% case study. Upper panel: hourly load and PV generation (kW). Lower panel:
% battery state of charge (kWh) tracking charge from midday surplus and
% discharge into evening, with one cloudy stretch drawing the diesel backup.
% Values are illustrative of the sized 250 kWp PV / 600 kWh battery system.
\begin{figure}[htbp]
\centering
\begin{tikzpicture}
\begin{axis}[
  width=13cm, height=5.2cm,
  name=gen,
  xlabel={hour of representative week},
  ylabel={power (kW)},
  xmin=0, xmax=168, ymin=0, ymax=260,
  xtick={0,24,48,72,96,120,144,168},
  ytick={0,100,200},
  axis lines=left,
  tick label style={font=\scriptsize},
  label style={font=\small},
  legend style={font=\scriptsize, at={(0.5,1.02)}, anchor=south, legend columns=2, draw=none},
]
% PV generation: clear-day humps, with days 4-5 cloudy (low humps)
\addplot[engamber, very thick] coordinates {
  (0,0)(6,0)(9,90)(12,230)(15,150)(18,10)(24,0)
  (30,0)(33,95)(36,235)(39,160)(42,12)(48,0)
  (54,0)(57,90)(60,225)(63,155)(66,10)(72,0)
  (78,0)(81,30)(84,70)(87,45)(90,8)(96,0)
  (102,0)(105,25)(108,55)(111,35)(114,6)(120,0)
  (126,0)(129,95)(132,230)(135,158)(138,11)(144,0)
  (150,0)(153,92)(156,228)(159,156)(162,10)(168,0)
};
\addlegendentry{PV generation}
% load: morning and evening peaks, ~42 kW average
\addplot[engblue, thick] coordinates {
  (0,28)(6,30)(8,55)(12,40)(18,75)(20,80)(22,50)(24,30)
  (30,55)(36,40)(42,78)(44,82)(48,30)
  (54,55)(60,40)(66,76)(68,80)(72,30)
  (78,55)(84,42)(90,78)(92,82)(96,30)
  (102,55)(108,42)(114,80)(116,84)(120,30)
  (126,55)(132,40)(138,76)(140,80)(144,30)
  (150,55)(156,40)(162,78)(164,82)(168,30)
};
\addlegendentry{household load}
\end{axis}
\begin{axis}[
  width=13cm, height=5.2cm,
  at={(gen.below south)}, anchor=north, yshift=-0.7cm,
  xlabel={hour of representative week},
  ylabel={state of charge (kWh)},
  xmin=0, xmax=168, ymin=0, ymax=650,
  xtick={0,24,48,72,96,120,144,168},
  ytick={0,300,600},
  axis lines=left,
  tick label style={font=\scriptsize},
  label style={font=\small},
]
% SoC: fills each clear day, partly depleted overnight; runs low days 4-5
\addplot[englime!70!black, very thick] coordinates {
  (0,300)(6,200)(12,520)(16,600)(20,360)(24,250)
  (30,180)(36,510)(40,600)(44,330)(48,220)
  (54,160)(60,500)(64,590)(68,310)(72,200)
  (78,150)(84,250)(88,260)(92,90)(96,40)
  (102,30)(108,120)(112,130)(116,30)(120,20)
  (126,30)(132,480)(136,600)(140,330)(144,230)
  (150,170)(156,500)(160,600)(164,340)(168,250)
};
% diesel-backup window marker during the cloudy low-SoC stretch
\draw[engred, dashed] (axis cs:88,0) -- (axis cs:88,650);
\draw[engred, dashed] (axis cs:120,0) -- (axis cs:120,650);
\node[engred, font=\scriptsize, anchor=south] at (axis cs:104,610) {diesel backup};
\end{axis}
\end{tikzpicture}
\caption[Microgrid dispatch over a representative week]{Representative-week
dispatch of the sized solar-plus-storage microgrid (case study). The upper
panel plots hourly PV generation against household load; the lower panel
plots the battery state of charge as the running integral of net power.
Clear days (1 to 3, 6 to 7) refill the battery to its usable ceiling and
ride the evening peak from storage. The cloudy stretch on days 4 and 5
drains the battery toward empty and forces the diesel backup to run for a
few evening hours (red window), the single contributor to the year's
loss-of-load expectation. Values are illustrative of the sized system.}
\label{fig:vol04ch14:microgrid-week}
\end{figure}


\subsection{Levelized cost and the diesel comparison}

The cost combines the annualised capital of each component with its
running cost over the delivered energy. Using representative installed
costs (current as of 2024): PV at $\$1000\,\si{\per\kilo\watt}$ AC gives
$\$250\text{k}$; battery at $\$300\,\si{\per\kilo\watt\hour}$ usable
(installed, including the inverter and balance-of-system) gives
$\$180\text{k}$; the $150\,\si{\kilo\watt}$ diesel genset and the
microgrid controls and interconnection add $\sim \$120\text{k}$. The
capital is $\sim \$550\text{k}$. The capital-recovery factor at a
$7\,\%$ discount rate over a $20$-year life is $\text{CRF} = i/(1 -
(1+i)^{-N}) = 0.07/(1 - 1.07^{-20}) \approx 0.094\,\text{yr}^{-1}$, so the
annualised capital is $\sim \$52\text{k}/\text{yr}$. Running costs add the
diesel fuel for the $4\,\%$ of energy it serves ($\sim 14.7\,\si{\mega
\watt\hour}/\text{yr}$ at $\sim 0.28\,\si{\liter\per\kilo\watt\hour}$ and
$\$1.20\,\si{\per\liter}$ is $\sim \$5\text{k}/\text{yr}$), plus
operations and maintenance at $\sim \$15\text{k}/\text{yr}$, a total
running cost of $\sim \$20\text{k}/\text{yr}$. Over the delivered
$368\,\si{\mega\watt\hour}/\text{yr}$ the levelized cost is
\[
\text{LCOE} = \frac{52\text{k} + 20\text{k}}{368\,\si{\mega\watt\hour}}
\approx \$0.20\,\si{\per\kilo\watt\hour}.
\]
A diesel-only microgrid serving the same load would carry a smaller
capital ($\sim \$200\text{k}$ for gensets, controls, and fuel storage)
but burn fuel for all $368\,\si{\mega\watt\hour}/\text{yr}$: at the same
fuel intensity that is $\sim 103{,}000\,\si{\liter}/\text{yr}$ and $\sim
\$124\text{k}/\text{yr}$ in fuel alone, plus heavier maintenance from the
high running hours, for an LCOE of $\sim \$0.40$ to $\$0.50\,\si{\per\kilo
\watt\hour}$. The solar-plus-storage-plus-diesel hybrid roughly halves the
delivered cost of energy and cuts the fuel burn (and its emissions and its
supply-chain exposure) by $\sim 96\,\%$.

\subsection{Verdict}

The hybrid wins on cost and on fuel security, and the design choices that
make it win are the chapter's recurring lessons in miniature: overbuild
the array to the knee of the reliability curve rather than to full
autonomy; size the battery for the daily cycle and one cloudy day, not for
the rare multi-day event; and keep a small dispatchable backup for the
long tail rather than paying for the storage to reach $100\,\%$ on
renewables alone. The verdict is sensitive to three numbers that carry
their year: the battery cost (falling, which favours more storage and less
diesel over time), the diesel fuel price (rising or volatile, which
favours the hybrid further), and the site's capacity factor (a cloudier
site needs more overbuild and more backup). The design is a screen, not a
final engineering specification; the full design adds the protection,
the grounding, the inverter grid-forming controls of the off-grid section,
and the financing structure, but the sizing logic and the verdict hold.

\section{A second case study: a regional grid decarbonisation pathway}\pathtag{standard}

The microgrid case ran the chapter's quantities through one small system at
one moment. A regional grid runs them through a whole network across a
quarter-century, and the integration questions change character: the
problem is no longer to size a battery but to steer a generation mix from a
fossil baseline to a low-carbon target while holding reliability, managing
curtailment, and paying for the firm capacity the renewables do not supply.
We take an illustrative mid-size temperate grid, a region with a 2025 peak
demand of $\sim 34\,\si{\giga\watt}$ and annual generation of $\sim
188\,\text{TWh}/\text{yr}$, and trace a representative pathway to 2050. The
numbers are illustrative of a real region's scale and carry their year; the
method, not the figures, is the transferable content.

\subsection{The 2025 baseline and the demand trajectory}

The baseline mix (current as of 2025) is coal-heavy: $\sim 37\,\%$ coal,
$\sim 29\,\%$ gas, $\sim 13\,\%$ nuclear, $\sim 10\,\%$ hydro, $\sim 6\,\%$
wind, $\sim 4\,\%$ solar, the left edge of Figure
\ref{fig:vol04ch14:decarb-pathway}. The grid carbon intensity is $\sim
380\,\si{\kilogram}\,\text{CO}_2/\si{\mega\watt\hour}$, the red point at
2025 in Figure \ref{fig:vol04ch14:emissions-cost-trajectory}, set mostly by
the coal block. Demand does not hold still over the pathway: electrifying
road transport and building heat shifts load off direct fuel and onto the
grid, and the standard planning assumption is a $\sim 1.5\,\%/\text{yr}$
rise in electricity demand even as total final energy falls, because
electric drives and heat pumps reject far less than the combustion they
replace. Over twenty-five years that compounds to a $\sim 45\,\%$ rise, so
the 2050 grid serves $\sim 270\,\text{TWh}/\text{yr}$ against the 2025
$\sim 188$, the rising total height of the stack in Figure
\ref{fig:vol04ch14:decarb-pathway}. The pathway must therefore decarbonise
a moving target: it has to displace the fossil generation and serve the
new electrified load at the same time.

\subsection{Generation mix: capacity versus energy}

The pathway retires coal first, because coal is both the highest-carbon and
the highest-marginal-cost bulk generator, so retiring it cuts emissions and
wholesale cost together. Gas is retired more slowly and not to zero: a
shrinking gas fleet is retained as firm dispatchable capacity for the
low-renewable tail, running at a falling capacity factor as the pathway
proceeds. Nuclear and hydro are life-extended and held roughly flat,
cheap firm low-carbon energy whose capital is already sunk. Wind and solar
carry nearly all of the growth, rising from a combined $\sim 10\,\%$
of generation in 2025 to $\sim 76\,\%$ in 2050. The distinction between
capacity and energy is the recurring trap. By 2050 the wind-and-solar
nameplate is large: matching $\sim 206\,\text{TWh}/\text{yr}$ of
wind-and-solar energy at blended capacity factors of $\sim 35\,\%$ (wind)
and $\sim 18\,\%$ (solar) needs on the order of $80\,\si{\giga\watt}$ of
combined renewable nameplate against a $\sim 40\,\si{\giga\watt}$ peak
demand, so the grid carries roughly twice its peak demand in renewable
nameplate alone. The overbuild is not waste; it is the price of covering
the low-resource periods, and it is exactly the overbuild-versus-curtailment
trade of Figure \ref{fig:vol04ch14:overbuild-curtailment} applied at
regional scale.

\subsection{Storage, curtailment, and the firm-capacity gap}

Two integration problems grow as the renewable share rises. The first is
curtailment: with the renewable nameplate at roughly twice peak demand and
a wind-and-solar energy share above $\sim 70\,\%$, the surplus at times of
high output exceeds what the grid can use or store, and the curtailed
fraction climbs toward $\sim 15$ to $20\,\%$ of renewable generation by
2050 even with substantial storage, consistent with the high-overbuild end
of Figure \ref{fig:vol04ch14:overbuild-curtailment}. The pathway absorbs
part of that surplus into storage and part into sector coupling
(electrolytic hydrogen, flexible EV charging, thermal storage), and accepts
the rest as the spilled cost of covering the lulls. The storage build rises
from near zero in 2025 to several gigawatts of $4$-to-$8$-hour lithium-ion
for the daily cycle by 2050, plus a smaller pumped-hydro or long-duration
component for the multi-day swings, the duration split the crossover
calculation predicts. The second problem is the firm-capacity gap. The
$\sim 40\,\si{\giga\watt}$ 2050 peak must be met with a reserve margin of
$\sim 15\,\%$, so $\sim 46\,\si{\giga\watt}$ of firm capacity, and the
capacity credit of the renewable fleet at that peak (a winter evening,
sun down, wind uncertain) is low. The firm capacity is supplied by the
retained nuclear and hydro ($\sim 6\,\si{\giga\watt}$), the shrunken gas
fleet ($\sim 12\,\si{\giga\watt}$), storage with enough duration to cover
the evening peak ($\sim 10$ to $15\,\si{\giga\watt}$ of firm contribution
after derating for duration), demand response ($\sim 3\,\si{\giga\watt}$),
and interconnection to neighbouring regions whose peaks do not coincide.
The pathway does not reach a pure-renewable grid; it reaches a
high-renewable grid with a firm low-carbon and dispatchable backbone, which
is the least-cost reading of the reliability constraint.

\subsection{Reliability and the reserve margin through the transition}

The reserve margin is the quantity that must not slip during the
transition, and it is easiest to lose in the middle years, when coal has
retired faster than the firm replacement has been built. The planning
discipline is to retire fossil capacity only as fast as firm low-carbon
capacity, storage, demand response, and interconnection replace its
capacity credit, not merely its energy. A coal plant retired for its energy
(displaced by cheap renewables) still carried firm capacity that the
renewables did not replace, and retiring it without a firm substitute
thins the reserve margin precisely in the low-renewable hours. The loss-of-
load expectation, the reliability currency of the earlier worked example,
is held at the standard $\sim 0.1\,\si{\hour}/\text{yr}$ planning target
across every year of the pathway by sequencing the retirements behind the
firm replacements, not by the average energy balance, which looks healthy
long before the reserve margin does. The Texas 2021 cold-weather event
\cite{acc:ferc-nerc-texas-2021} is the standing warning: a correlated loss
of firm capacity (gas plants de-rating together in extreme cold, with no
interconnection to import from) violated the independence the reserve
margin assumed, and the average energy balance was no protection. A
decarbonisation pathway that watches only annual energy and carbon, and not
the firm-capacity reserve in the worst correlated hour, is the same mistake
written across a quarter-century.

\subsection{Emissions and cost trajectory}

Figure \ref{fig:vol04ch14:emissions-cost-trajectory} traces the two
outputs the pathway is steered by. Carbon intensity falls from $\sim 380$
to $\sim 18\,\si{\kilogram}\,\text{CO}_2/\si{\mega\watt\hour}$, a $\sim
95\,\%$ cut, steeply through the coal retirement and then more slowly as
the last fossil generation is the firm gas tail that runs rarely but is
hard to remove. The cumulative emissions, the area under the curve, are
what the climate constraint actually bounds, and front-loading the coal
retirement cuts that area more than an equal-paced glide would, which is the
case for retiring the dirtiest capacity first. The system cost dips through
the middle of the pathway as low-marginal-cost wind and solar enter and
push the merit-order stack to the right (the merit-order effect of the
worked example), then rises modestly toward 2050 as the last emissions are
removed by paying for firm capacity, storage, transmission, and curtailed
overbuild rather than for cheap bulk energy. The cost curve is the
shape the capacity-credit mastery box predicts: the first emissions are
cheap to cut because cheap renewables displace expensive coal, and the last
emissions are dear because they are removed by firming and overbuild whose
value is capacity and flexibility, not energy. The endpoint cost of $\sim
\$74\,\si{\per\mega\watt\hour}$ sits modestly above the mid-pathway trough,
and the comparison the region actually faces is not against the trough but
against the counterfactual of holding the coal fleet, whose cost rises with
carbon pricing and fuel and whose external cost in emissions the pathway
removes.

\subsection{Verdict}

The pathway reaches a $\sim 95\,\%$ carbon cut by 2050 at a system cost
modestly above the cheapest mid-transition point and well below a
carbon-priced fossil counterfactual, and the design choices that make it
work are the chapter's lessons at regional scale: retire the dirtiest and
most expensive capacity first; hold firm low-carbon capacity (nuclear,
hydro) that is already paid for; keep a shrinking dispatchable tail rather
than chasing the last few percent of emissions with storage alone;
overbuild renewables to the knee of the curtailment curve and couple the
surplus to flexible loads; and sequence every fossil retirement behind a
firm replacement so the reserve margin never thins. The binding constraints
are not the cost of generation, which falls, but the build rate of
transmission and the procurement of firm capacity and flexibility, the
three products the capacity-credit argument insists be valued separately.
The pathway is a planning screen, not a dispatch model; a real plan runs
the hour-by-hour dispatch of the simulation exercise across many weather
years to size the storage and firm capacity against the worst correlated
events, and it carries the financing, supply-chain, and workforce
constraints the chapter names but does not quantify. The numbers carry
their year and will look different by 2050; the structure of the problem,
steering capacity and energy and firm reserve together while curtailment
and cost trade against each other, is durable.

% Figure: regional generation-mix trajectory over a decarbonisation pathway.
% Stacked area of annual energy by source from a 2025 baseline to a 2050
% target, illustrative of a mid-size temperate grid. Stack plots via
% cumulative coordinates. Values illustrative, current as of 2025.
\begin{figure}[htbp]
\centering
\begin{tikzpicture}
\begin{axis}[
  width=13cm, height=8cm,
  xlabel={year},
  ylabel={annual generation (TWh/yr)},
  xmin=2025, xmax=2050, ymin=0, ymax=260,
  xtick={2025,2030,2035,2040,2045,2050},
  ytick={0,50,100,150,200,250},
  axis lines=left,
  stack plots=y,
  area style,
  legend pos=north west,
  legend cell align=left,
  tick label style={font=\scriptsize},
  label style={font=\small},
  legend style={font=\scriptsize, fill opacity=0.85, draw=none},
]
% coal, retired by mid-pathway
\addplot[draw=engblack, fill=enggray!70] coordinates {
  (2025,70)(2030,40)(2035,15)(2040,2)(2045,0)(2050,0)} \closedcycle;
\addlegendentry{coal}
% gas, declining but firm tail
\addplot[draw=engred, fill=engred!45] coordinates {
  (2025,55)(2030,55)(2035,45)(2040,32)(2045,20)(2050,12)} \closedcycle;
\addlegendentry{gas}
% nuclear, life-extended flat
\addplot[draw=engblue, fill=engblue!55] coordinates {
  (2025,25)(2030,25)(2035,25)(2040,24)(2045,22)(2050,22)} \closedcycle;
\addlegendentry{nuclear}
% hydro, flat
\addplot[draw=enggray, fill=enggray!35] coordinates {
  (2025,18)(2030,18)(2035,19)(2040,19)(2045,20)(2050,20)} \closedcycle;
\addlegendentry{hydro}
% wind, rising
\addplot[draw=englime!60!black, fill=englime!50] coordinates {
  (2025,12)(2030,30)(2035,55)(2040,78)(2045,95)(2050,108)} \closedcycle;
\addlegendentry{wind}
% solar, rising fastest
\addplot[draw=engamber, fill=engamber!50] coordinates {
  (2025,8)(2030,28)(2035,52)(2040,72)(2045,88)(2050,98)} \closedcycle;
\addlegendentry{solar}
\end{axis}
\end{tikzpicture}
\caption[Regional decarbonisation generation-mix pathway]{An illustrative
generation-mix trajectory for a mid-size temperate grid from a 2025 baseline
to a 2050 low-carbon target (current as of 2025). Coal retires first, gas
shrinks to a firm dispatchable tail, life-extended nuclear and hydro hold
roughly flat, and wind and solar carry the growth. Total generation rises
because electrification of transport and heat lifts demand even as it raises
end-use efficiency. The stacked height is gross annual energy, not firm
capacity; the firm-capacity and reliability accounting of the pathway is a
separate question worked in the case study. Quantities illustrative.}
\label{fig:vol04ch14:decarb-pathway}
\end{figure}


% Figure: emissions and system-cost trajectory over the decarbonisation
% pathway. Two y-axes: grid carbon intensity (left) falling toward zero,
% and average system cost of electricity (right) dipping then rising
% slightly as the firm-capacity and storage build is paid for. Illustrative,
% current as of 2025.
\begin{figure}[htbp]
\centering
\begin{tikzpicture}
\begin{axis}[
  width=13cm, height=8cm,
  axis y line*=left,
  xlabel={year},
  ylabel={grid carbon intensity (kg CO$_2$/MWh)},
  xmin=2025, xmax=2050, ymin=0, ymax=420,
  xtick={2025,2030,2035,2040,2045,2050},
  ytick={0,100,200,300,400},
  tick label style={font=\scriptsize},
  label style={font=\small},
  legend style={font=\scriptsize, at={(0.97,0.97)}, anchor=north east, draw=none},
]
\addplot[engred, very thick, mark=*, mark size=1.3pt] coordinates {
  (2025,380)(2030,290)(2035,185)(2040,95)(2045,42)(2050,18)
};
\addlegendentry{carbon intensity}
\end{axis}
\begin{axis}[
  width=13cm, height=8cm,
  axis y line*=right,
  axis x line=none,
  ylabel={system cost (USD/MWh)},
  xmin=2025, xmax=2050, ymin=0, ymax=140,
  ytick={0,40,80,120},
  tick label style={font=\scriptsize},
  label style={font=\small},
  legend style={font=\scriptsize, at={(0.97,0.80)}, anchor=north east, draw=none},
]
\addplot[engblue, very thick, dashed, mark=square*, mark size=1.3pt] coordinates {
  (2025,85)(2030,72)(2035,64)(2040,62)(2045,68)(2050,74)
};
\addlegendentry{system cost}
\end{axis}
\end{tikzpicture}
\caption[Emissions and cost trajectory of the pathway]{Grid carbon
intensity (red, left axis) and average system cost of electricity (blue
dashed, right axis) over the illustrative decarbonisation pathway (current
as of 2025). Carbon intensity falls steeply as low-marginal-cost wind and
solar displace coal and gas. System cost dips through the middle of the
pathway as cheap renewables enter, then rises modestly toward 2050 as the
last fraction of emissions is removed by paying for firm capacity, storage,
and transmission rather than for cheap bulk energy. The cost curve is the
shape that the capacity-credit and curtailment arguments predict: the first
emissions are cheap to cut, the last are dear. Quantities illustrative.}
\label{fig:vol04ch14:emissions-cost-trajectory}
\end{figure}


\section{A third case study: a green-hydrogen power-to-X plant}\pathtag{standard}

The two preceding cases moved electricity in time and steered a grid across
decades. The third takes the surplus that overbuild and curtailment leave
on the table and converts it into a storable chemical, the power-to-X path
that the capacity-credit mastery box named as the home for cheap curtailed
energy. We size a green-hydrogen plant end to end, from the renewable
surplus that feeds it through the electrolyser, compression, and storage to
the cost per kilogram and the round trip back to electricity, and we read
off where the economics live. The numbers are representative of a
mid-2020s plant co-located with an overbuilt wind-and-solar farm and carry
their year; the method transfers once the local electricity price and
electrolyser cost replace the illustrative figures.

\subsection{The energy source and the electrolyser size}

The plant runs on the surplus of an overbuilt renewable farm: the energy
that would otherwise be curtailed, available at a low price but only when
the farm is generating above local demand. Take a farm whose curtailed
surplus totals $\sim 175\,\si{\giga\watt\hour}/\text{yr}$, delivered in the
intermittent pattern of high-output hours rather than steadily. The
electrolyser cannot run at a high capacity factor on this supply: the
surplus is present for perhaps $40\%$ of the hours of the year, so an
electrolyser sized to absorb the surplus peaks runs at a capacity factor of
$\sim 0.30$ to $0.40$. We size the electrolyser to a $40\,\si{\mega\watt}$
input rating, which at a $0.30$ capacity factor draws
\[
E_{\text{in}} = 40\,\si{\mega\watt}\times 0.30\times 8760\,\si{\hour}
\approx 105\,\si{\giga\watt\hour}/\text{yr},
\]
absorbing the bulk of the surplus while leaving the rare highest-output
hours curtailed (oversizing the electrolyser to catch them would lower its
capacity factor and raise the cost per kilogram, the same overbuild trade
the grid faces). The choice of $40\,\si{\mega\watt}$ over the surplus is
the power-to-X version of the array-to-inverter ratio in the microgrid
case: size the converter below the resource peak and accept spilling the
tail.

\subsection{Electrolysis efficiency and the overpotential}

A water-electrolysis cell splits water by passing charge through it, and
the energy per kilogram of hydrogen is set by the cell voltage. The
thermodynamic minimum at standard conditions is the reversible voltage
$E_{\text{rev}} \approx 1.23\,\si{\volt}$, or $1.48\,\si{\volt}$ at the
thermoneutral point that also supplies the reaction's heat. Real cells run
above the thermoneutral voltage by an \engterm{overpotential} that drives
the reaction faster than equilibrium: activation overpotential at each
electrode, ohmic loss in the electrolyte and membrane, and concentration
loss at high current. A modern PEM or alkaline cell operates near
$V_{\text{cell}}\approx 1.9\,\si{\volt}$ at its rated current density, so
its voltage efficiency referenced to the thermoneutral value is
\[
\eta_V = \frac{1.48\,\si{\volt}}{1.9\,\si{\volt}} \approx 0.78,
\]
and after the auxiliary loads (the balance of plant: pumps, gas drying,
power electronics) the system efficiency lands near $\eta_{\text{sys}}
\approx 0.65$. The specific electricity consumption is the inverse of that
efficiency against hydrogen's heating value. Using the higher heating value
$\sim 39.4\,\si{\kilo\watt\hour\per\kilogram}$,
\[
e_{\text{H}_2} = \frac{39.4\,\si{\kilo\watt\hour\per\kilogram}}{0.65}
\approx 50\,\si{\kilo\watt\hour\per\kilogram} \;\text{(HHV basis)},
\]
the $\sim 50\,\si{\kilo\watt\hour}/\si{\kilogram}$ that recurs across the
power-to-X literature. At that intensity the $105\,\si{\giga\watt\hour}/
\text{yr}$ of input produces
\[
m_{\text{H}_2} = \frac{105\times 10^6\,\si{\kilo\watt\hour}}
{50\,\si{\kilo\watt\hour\per\kilogram}} \approx 2.1\times 10^6\,\si{\kilogram}
/\text{yr} \approx 2100\,\text{t}/\text{yr}
\]
of hydrogen. The efficiency cascade of Figure
\ref{fig:vol04ch14:hydrogen-chain} starts from this electrolyser stage and
follows the energy to the far end.

% Figure: green-hydrogen power-to-power efficiency cascade. Each bar is the
% energy remaining after a stage, normalised to 100 units of input
% electricity, so the bar height is the cumulative efficiency. The last bar
% is the electricity recovered through the fuel cell; the gaps are the
% stage losses. Values illustrative, current as of 2024.
\begin{figure}[htbp]
\centering
\begin{tikzpicture}
\begin{axis}[
  width=13cm, height=7.5cm,
  ylabel={energy remaining (\% of input electricity)},
  symbolic x coords={input, electrolysis, compression, storage, fuelcell},
  xtick={input, electrolysis, compression, storage, fuelcell},
  xticklabels={input\\electricity, after\\electrolysis, after\\compression, after\\storage, after\\fuel cell},
  xticklabel style={align=center, font=\scriptsize},
  ymin=0, ymax=105,
  ytick={0,20,40,60,80,100},
  ybar,
  bar width=20pt,
  nodes near coords,
  nodes near coords style={font=\scriptsize},
  tick label style={font=\scriptsize},
  label style={font=\small},
  enlarge x limits=0.12,
]
\addplot[draw=engblue, fill=engblue!55] coordinates {
  (input,100)
  (electrolysis,65)
  (compression,59)
  (storage,56)
  (fuelcell,31)
};
\end{axis}
\end{tikzpicture}
\caption[Green-hydrogen power-to-power efficiency cascade]{The
power-to-power efficiency cascade for green hydrogen, normalised to $100$
units of input electricity (current as of 2024). Electrolysis keeps about
$65$ units, compression to $70\,\si{\mega\pascal}$ costs another few,
storage and handling a little more, and the fuel cell recovers about $55\%$
of what reaches it, leaving roughly $31$ units as electricity at the far
end. The losses compound multiplicatively, so the round trip returns about
a third of the input. The two large drops, at the electrolyser and at the
fuel cell, are where the engineering effort to raise the round trip
concentrates. The cascade is why hydrogen is reserved for long-duration
and hard-to-electrify uses where the round-trip penalty is paid rarely.}
\label{fig:vol04ch14:hydrogen-chain}
\end{figure}


\subsection{Compression and storage}

Hydrogen leaves a PEM electrolyser at modest pressure and must be
compressed for storage and transport. Compressing to a storage pressure of
$70\,\si{\mega\pascal}$ costs work; the ideal isothermal work to compress a
gas by a pressure ratio $r$ is $w = (RT/M)\ln r$ per unit mass, and for
hydrogen from $\sim 3\,\si{\mega\pascal}$ at the electrolyser to $70\,\si{
\mega\pascal}$ the ratio is $\sim 23$, so the ideal work is
\[
w \approx \frac{(8.314\,\si{\joule\per\mole\per\kelvin})(298\,\si{\kelvin})}
{2.016\times 10^{-3}\,\si{\kilogram\per\mole}}\ln 23
\approx 3.85\,\si{\mega\joule\per\kilogram} \approx 1.07\,\si{\kilo\watt
\hour\per\kilogram},
\]
which real multi-stage compressors with intercooling deliver at $\sim 60\%$
efficiency, so $\sim 1.8\,\si{\kilo\watt\hour\per\kilogram}$, about $5\%$ of
the hydrogen's energy. Storage in steel or composite pressure vessels adds
a small parasitic loss to handling and slow leakage; the cascade keeps
$\sim 56$ of the original $100$ units of electricity by the time the
hydrogen is stored and ready to use. The storage volume is set by the gas
density at $70\,\si{\mega\pascal}$, $\sim 40\,\si{\kilogram\per\meter
\cubed}$, so a day of the plant's output ($\sim 5.8\,\text{t}$) needs
$\sim 145\,\si{\meter\cubed}$ of pressurised volume, a bank of vessels
rather than a reservoir; this is the high power cost and low energy cost
that the power-energy mastery box assigned to chemical storage.

\subsection{End use: the fork in the path}

The stored hydrogen can go back to electricity through a fuel cell or be
used directly, and the two paths have different economics because the fuel
cell discards more than half the energy again. Through a fuel cell at $\sim
55\%$ efficiency the power-to-power round trip keeps $\sim 0.65\times 0.95
\times 0.95\times 0.55 \approx 0.31$ of the input (Figure
\ref{fig:vol04ch14:hydrogen-chain}), the $\sim 31\%$ that makes hydrogen a
poor general-purpose battery and a reasonable seasonal one, used rarely.
The direct-use paths skip the fuel-cell loss: hydrogen fed to ammonia
synthesis, to direct-reduced-iron steelmaking, or to a refinery replaces
fossil feedstock at the $\sim 56\%$ delivered efficiency of the stored gas,
not the $\sim 31\%$ of the round trip, and these hard-to-electrify uses are
where green hydrogen competes first. The plant's verdict therefore depends
on what it feeds, and the cost per kilogram is the quantity that decides.

\subsection{The cost per kilogram}

The levelized cost of hydrogen splits into an electricity term and a
capital term, and the split is the whole story. The electricity term is the
specific consumption times the electricity price: at a low surplus price of
$\$20\,\si{\per\mega\watt\hour}$,
\[
c_{\text{elec}} = 50\,\si{\kilo\watt\hour\per\kilogram}\times
\frac{\$20}{1000\,\si{\kilo\watt\hour}} = \$1.0\,\si{\per\kilogram}.
\]
The capital term amortises the electrolyser over its output. At an
installed cost of $\sim \$1200\,\si{\per\kilo\watt}$ for the $40\,\si{\mega
\watt}$ unit, the capital is $\sim \$48\text{M}$; at the $9.4\%$
capital-recovery factor of the earlier examples that is $\sim \$4.5\text{M}/
\text{yr}$, and over the $2100\,\text{t}/\text{yr}$ of output it adds
\[
c_{\text{cap}} = \frac{\$4.5\text{M}/\text{yr}}{2.1\times 10^6\,\si{
\kilogram}/\text{yr}} \approx \$2.1\,\si{\per\kilogram}.
\]
Compression and balance-of-plant running cost add $\sim \$0.4\,\si{\per
\kilogram}$, for a levelized cost of $\sim \$3.5\,\si{\per\kilogram}$ at the
$0.30$ electrolyser capacity factor. Two readings follow, both visible in
Figure \ref{fig:vol04ch14:hydrogen-cost}. First, the capital term falls
steeply with utilisation: running the same electrolyser at a $0.60$
capacity factor (a steadier supply, or a grid connection that fills the
gaps) halves $c_{\text{cap}}$ to $\sim \$1.1\,\si{\per\kilogram}$ and the
total to $\sim \$2.5\,\si{\per\kilogram}$, which is the tension at the heart
of power-to-X: the cheap curtailed electricity that justifies the plant is
exactly the intermittent supply that starves the electrolyser of running
hours and inflates its capital per kilogram. Second, the electricity term
dominates at any realistic price: every $\$10\,\si{\per\mega\watt\hour}$ on
the electricity price adds $\$0.5\,\si{\per\kilogram}$, so green hydrogen is
cheap only where electricity is cheap, and the slope of Figure
\ref{fig:vol04ch14:hydrogen-cost} is the same for both utilisation cases
because it is set by the $50\,\si{\kilo\watt\hour}/\si{\kilogram}$ specific
consumption alone.

% Figure: levelized cost of green hydrogen against electricity price, for
% two electrolyser capacity factors. Cost per kg = electricity term
% (price * specific energy / efficiency) plus an amortised capex term that
% falls with utilisation. Shows electricity price as the dominant lever.
% Values illustrative, current as of 2024.
\begin{figure}[htbp]
\centering
\begin{tikzpicture}
\begin{axis}[
  width=12.5cm, height=7.5cm,
  xlabel={electricity price (USD/MWh)},
  ylabel={levelized cost of hydrogen (USD/kg)},
  xmin=0, xmax=80, ymin=0, ymax=8,
  xtick={0,20,40,60,80},
  ytick={0,2,4,6,8},
  grid=both,
  grid style={enggray!25},
  tick label style={font=\scriptsize},
  label style={font=\small},
  legend style={font=\scriptsize, at={(0.03,0.97)}, anchor=north west, draw=none},
]
% low utilisation (CF 0.30): high capex term, slope from 50 kWh/kg / 0.65
\addplot[engorange, very thick, mark=*, mark size=1.3pt] coordinates {
  (0,2.4)(20,3.94)(40,5.48)(60,7.02)(80,8.0)
};
\addlegendentry{electrolyser CF $0.30$}
% high utilisation (CF 0.60): lower capex term, same electricity slope
\addplot[engblue, very thick, dashed, mark=square*, mark size=1.3pt] coordinates {
  (0,1.2)(20,2.74)(40,4.28)(60,5.82)(80,7.36)
};
\addlegendentry{electrolyser CF $0.60$}
% reference band for fossil (grey) hydrogen, ~1.5-2 USD/kg
\draw[enggray, thick] (axis cs:0,1.75) -- (axis cs:80,1.75);
\node[anchor=south east, font=\scriptsize, color=enggray] at (axis cs:80,1.75) {fossil ``grey'' H$_2$};
\end{axis}
\end{tikzpicture}
\caption[Levelized cost of green hydrogen]{Levelized cost of green
hydrogen against the electricity price, for a well-utilised electrolyser
(capacity factor $0.60$) and a poorly utilised one ($0.30$), current as of
2024. The slope is the same for both because it is set by the electricity
energy per kilogram divided by the electrolyser efficiency, about $50\,\si{
\kilo\watt\hour}/\si{\kilogram}$ at $\sim 65\%$ efficiency; the vertical
offset is the amortised capital, which is larger per kilogram when the
electrolyser runs fewer hours. Electricity price is the dominant lever:
halving it cuts the cost far more than any plausible capital reduction. The
grey line is the rough cost of fossil-derived hydrogen, the incumbent green
hydrogen must beat. Values illustrative.}
\label{fig:vol04ch14:hydrogen-cost}
\end{figure}


\subsection{Verdict}

The plant turns $\sim 105\,\si{\giga\watt\hour}/\text{yr}$ of otherwise
curtailed renewable surplus into $\sim 2100\,\text{t}/\text{yr}$ of
hydrogen at $\sim \$3.5\,\si{\per\kilogram}$, above the $\sim \$1.5$ to
$\$2\,\si{\per\kilogram}$ of the fossil-derived hydrogen it must displace
(current as of 2024) but on a cost curve that falls with cheaper
electrolysers and cheaper surplus electricity. The design choices echo the
chapter's recurring lessons: size the converter below the resource peak and
spill the tail; chase utilisation, because the capital per kilogram is set
by running hours; and send the hydrogen to a direct use that skips the
fuel-cell loss rather than to a round trip that throws away two thirds of
the energy. The verdict is sensitive to three numbers that carry their
year: the electrolyser capital cost (falling, which favours more plants),
the surplus electricity price (which depends on how overbuilt the grid
becomes), and the carbon or policy premium on the fossil incumbent (which
sets the price green hydrogen must beat). The plant is a screen, not a final
design; the full design adds the water treatment, the gas purification to
the purity the end use demands, the safety engineering for a wide-flammable-
range gas, and the offtake contract that gives the output a buyer. The
sizing logic and the verdict hold: power-to-X is the sink that turns the
curtailment problem of an overbuilt grid into a supply of low-carbon
molecules for the sectors electricity cannot reach directly.

\begin{mastery}
\textbf{The electricity bill is the hydrogen cost.} A green-hydrogen plant
inverts the cost structure of the storage technologies earlier in the
chapter. Pumped hydro and lithium-ion are capital-dominated: once built,
their marginal cost per cycle is small. An electrolyser is
energy-dominated, because it consumes $\sim 50\,\si{\kilo\watt\hour}$ of
electricity per kilogram of product and that electricity is bought at
market price every hour the plant runs. The consequence is a sharp
sensitivity that the levelized-cost arithmetic makes plain: the cost per
kilogram is $e_{\text{H}_2}\cdot p_{\text{elec}} + c_{\text{cap}}/(
\text{CF}\cdot 8760)$, a line in the electricity price whose slope is the
specific consumption and whose intercept is the amortised capital divided
by the running hours. Two levers move it, and they pull against each other.
Cheaper electricity lowers the slope-driven term and is the dominant lever,
which is why green hydrogen is sited where renewable electricity is cheapest
and most abundant. Higher utilisation lowers the capital term, but the
cheapest electricity is the curtailed surplus, which is intermittent and
caps the utilisation, so the plant cannot have both the cheapest power and
the highest running hours from the same source. The resolution in practice
is a blend: run on cheap surplus when it is there and on cheap firm
low-carbon power (nuclear, hydro, or a power-purchase contract) to fill the
gaps, trading a little higher average electricity price for the lower
capital-per-kilogram that high utilisation buys. The same arithmetic
governs every power-to-X conversion, synthetic methane, methanol, ammonia,
and sustainable aviation fuel, because each adds further conversion losses
and capital on top of the hydrogen, multiplying the electricity per unit of
final product and sharpening the dependence on a cheap, abundant,
low-carbon electron.
\end{mastery}

\section{Two applied vignettes: an island interconnector and a flexible steel mill}\pathtag{standard}

The chapter's quantities appear again in two systems it has not yet touched: a
small island grid deciding between a subsea cable and more diesel, and an
electric-arc steel mill that turns a heavy industrial load into a grid
resource. Each is a fresh setting for the same reasoning, and each carries its
year.

\subsection{An island grid and the subsea interconnector}

A populated island runs its own small grid, and the recurring question is
whether to keep generating locally, usually on imported diesel, or to lay a
subsea cable to the mainland grid. Take an island with a peak demand of
$60\,\si{\mega\watt}$, an average load of $\sim 35\,\si{\mega\watt}$, and an
annual energy of $\sim 300\,\si{\giga\watt\hour}/\text{yr}$, served today by
diesel gensets burning fuel shipped in at a delivered price that makes the
island's electricity cost $\sim \$0.30\,\si{\per\kilo\watt\hour}$, several
times the mainland wholesale price. Three options compete. The first is more
of the same: diesel is cheap to install and dispatchable, but its fuel cost
and its exposure to a supply chain that a storm can cut are the reasons to
leave it. The second is a local renewable build, wind and solar plus a
battery, which lowers the fuel burn but faces the small-system stability
problem of the off-grid section: the island's inertia is tiny, so a
high-inverter build needs grid-forming inverters and fast storage to hold the
frequency, and the winter renewable drought still calls for a diesel backup
the island cannot retire. The third is a subsea HVDC or HVAC interconnector to
the mainland.

The interconnector is a transmission problem of the kind the HVDC-versus-HVAC
example worked. A subsea link is almost always direct current once it passes a
few tens of kilometres, because an AC subsea cable draws a large charging
current that grows with length and consumes the cable's capacity before the
resistive loss matters, exactly the mechanism that collapses the AC-versus-DC
break-even for cable to $\sim 50\,\si{\kilo\meter}$. For a $50\,\si{\kilo
\meter}$ crossing carrying the island's $60\,\si{\mega\watt}$ peak, a modest
HVDC link with converter stations at each end suffices, and its cost is
dominated by the two converter terminals plus the cable, of order a few
hundred million dollars. The cable buys three things at once: it replaces the
expensive diesel energy with cheap mainland wholesale energy, it supplies the
firm capacity and inertia the island lacked (the mainland grid becomes the
island's reserve and its frequency reference), and it lets the island export
its own renewable surplus rather than curtailing it. The reading is the
transmission mastery box at island scale: a megawatt-kilometre of cable that
moves energy that already exists on the mainland is cheaper per unit of firmed
energy than the storage the island would otherwise need, and the binding
constraint is the capital and the permitting of the crossing, not the physics.
The cable also introduces a single largest-infeed the island must plan for:
losing the link at peak import is the loss of the island's largest source, so
the island retains a shrunken diesel or battery reserve sized to that
contingency, the largest-single-infeed rule applied to a grid whose largest
infeed is now a wire. The verdict for most inhabited islands within cable
reach is the interconnector plus a local renewable build plus a small reserve,
with the diesel retained for the cable-outage contingency rather than for
routine energy, which cuts the delivered cost of energy and the fuel burn
together, the same hybrid logic the microgrid case reached by a different
route.

\subsection{An electric-arc-furnace steel mill as a flexible load}

Primary steelmaking from iron ore uses a coke-fired blast furnace, one of the
largest single carbon sources in industry; recycling scrap steel in an
\engterm{electric arc furnace} instead melts the scrap with an electric arc,
replacing the coke with electricity and turning a chemical process into a
large electrical load. Take a mill with a $120\,\si{\mega\watt}$ furnace that
melts a $150\,\text{t}$ heat of scrap in a $\sim 50\,\si{\minute}$ tap-to-tap
cycle. The electricity per tonne is set by the energy to heat and melt steel,
whose specific heat and latent heat of fusion demand a thermodynamic minimum
near $\sim 350\,\si{\kilo\watt\hour}/\text{t}$, and a real furnace uses
$\sim 400$ to $500\,\si{\kilo\watt\hour}/\text{t}$ after the arc, chemical, and
radiation losses, so a $150\,\text{t}$ heat draws $\sim 60$ to $75\,\si{\mega
\watt\hour}$ over the melt. The furnace draws that energy in a short, intense
burst at a high and fluctuating power, and the fluctuation is itself an
engineering problem: the arc is electrically noisy, drawing an unbalanced,
rapidly varying current that causes \engterm{voltage flicker} on the supplying
grid, which is why arc-furnace mills sit on strong transmission connections and
carry reactive compensation (static var compensators) to hold the voltage
against the swinging load.

The mill's value to a high-renewable grid is that the melt can be scheduled.
A heat is a discrete, deferrable block of energy: the mill can choose to run
its furnace in the hours of cheap, abundant renewable output and idle it in
the expensive scarcity hours, shifting a large load off the peak and onto the
surplus the overbuild-and-curtailment argument leaves on the table. A mill
running $\sim 20$ heats a day has real freedom to cluster them in the cheap
hours, so it behaves as a flexible load that flattens the net-load curve from
the demand side, the demand-response mechanism at industrial scale and at a
size, tens of megawatts shiftable per mill, that matters to the grid. The
merit-order and demand-response examples give the value directly: moving the
melt from a scarcity hour at $\$140\,\si{\per\mega\watt\hour}$ to a surplus
hour near $\$20$ saves $\sim \$120\,\si{\per\mega\watt\hour}$ on the
$\sim 70\,\si{\mega\watt\hour}$ of a heat, $\sim \$8000$ a heat, and it does
so while relieving the grid of a peak it would otherwise have built capacity to
serve. The mill also decarbonises twice over: it replaces coke with
electricity, and it lets that electricity be the cheap renewable surplus, so an
arc-furnace mill sited near an overbuilt wind-and-solar farm and scheduled to
its output is the sector-coupling argument realised in heavy industry. The
constraints are that scrap availability limits how much steel can be made this
way (primary ore-based steel still needs the hydrogen-and-direct-reduced-iron
route of the transition itemisation), and that the flicker and reactive
demand require a strong grid connection, so the mill is a flexible resource
only where the transmission is stout enough to carry its swinging load.

\subsection{An aggregated electric-vehicle fleet as a distributed battery}

The chapter has treated storage as dedicated hardware, a battery or a
reservoir built to hold energy for the grid. A large fleet of electric
vehicles is a store the grid did not have to build, because the batteries
already exist to move the cars, and the vignette works the arithmetic of
turning parked cars into a grid resource. Take a region with one million
electric vehicles, each carrying a $\sim 60\,\si{\kilo\watt\hour}$ battery, so
the aggregate energy on wheels is $60\,\si{\giga\watt\hour}$, comparable to a
large fleet of grid batteries. The energy is not all available: a car is
driving perhaps $5\,\%$ of the time, so $\sim 95\,\%$ are parked at any
moment, but of the parked cars only the fraction plugged in can respond, and
of that fraction each owner permits only a slice of the battery to be used for
the grid to protect range and cycle life. A working assumption is that at the
evening peak $\sim 30\,\%$ of the fleet is plugged in and each offers
$\sim 10\,\si{\kilo\watt\hour}$ and $\sim 7\,\si{\kilo\watt}$ of discharge
power, so the dispatchable aggregate is $0.30\times 10^6\times 10 = 3\,\si{
\giga\watt\hour}$ of energy and $0.30\times 10^6\times 7 = 2.1\,\si{\giga
\watt}$ of power, a virtual power plant of gigawatt scale assembled from cars.

The value follows the same three currencies the chapter has separated.
Managed unidirectional charging, the simplest mode, shifts the fleet's charging
load off the evening peak and onto the midday and overnight surplus, flattening
the net-load curve from the demand side at no cost to the battery beyond the
smart-charging controller, and it is the mode that scales first because it asks
nothing of the owner but a delay. Bidirectional \engterm{vehicle-to-grid}
discharge, feeding energy back during the evening peak, adds firm capacity and
fast frequency response, the high-value services the capacity-credit and
ancillary-service arguments price, but it cycles the battery and so trades cell
degradation for grid revenue, which is why the early value is in the fast,
shallow, high-price services (frequency response, a short evening peak) rather
than in deep daily energy arbitrage. The effective-load-carrying-capability
reading applies directly: a fleet able to discharge across a three-hour evening
peak carries an ELCC near its available power as long as the peak is shorter
than the energy each car offers, and its capacity credit collapses once the
peak stretches beyond the shallow slice the owners permit, exactly the
duration-limited firmness of a short battery.

The engineering constraints are aggregation and the distribution grid, not the
cars. A single car is noise; a million cars are a resource only through an
aggregator that forecasts availability, bids the fleet into the market, and
dispatches it against the owners' driving needs, the virtual-power-plant
business model at fleet scale. The distribution grid is the binding physical
limit: a neighbourhood of cars charging or discharging together drives the
feeder voltage the feeder-voltage example bounded, so the same hosting-capacity
constraint that caps rooftop solar caps uncoordinated fleet charging, and the
management that smooths the coincident peak (the diversity-factor example) is
what keeps the fleet inside the feeder's limits. The reading is that the
transition's largest new load, the electrified car fleet, is also potentially
its largest distributed store and its most widely held flexible resource, and
whether it becomes a burden on the grid or a support turns on the controls and
the tariffs, not on the batteries, which are already there.

\section{Failure: grid blackouts, runaway batteries, generation outages}\pathtag{core}

Three failure modes of large energy systems recur, each with
characteristic engineering content.

\engterm{Grid blackouts} are the worst-case failure mode of a
synchronous AC grid. They begin with an event (a generator trip, a
transmission-line fault, a sudden load surge) that drives the
frequency or a transmission corridor outside its operational limits.
Protective relays trip generators or lines, the imbalance grows, and
the cascade propagates faster than human or automated response can
arrest. The grid separates into electrically isolated islands; some
islands survive on their own generation, others lose generation-load
balance and collapse to zero voltage.

Historic events trace the recurring patterns. The August 2003
Northeast blackout affected $\sim 50$ million customers in the
northeastern United States and Ontario after a transmission-line
fault in Ohio went unaddressed by an inadequate utility-control-room
alarm system, propagating over the next $90\,\si{\minute}$ across the
interconnection. The February 2021 Texas blackout left $\sim 4.5$
million customers without power for up to four days after natural-
gas-fired generators failed in extreme cold weather, while wind and
solar generation were also reduced; the ERCOT grid is not
synchronously connected to the rest of the United States and could
not import emergency capacity. The 2003 event is described here by
mechanism; the 2021 Texas event is recorded in the named-case registry,
and this chapter narrates it by mechanism rather than invoking the
named case, deferring the formal citation to a chapter that frames the
case directly.

Figure \ref{fig:vol04ch14:cascade-network} traces the transmission-overload
cascade that underlies the historic events. A generator feeds a load over
several parallel corridors, each carrying its share. When the first line
trips, its flow does not vanish; it redistributes onto the survivors, which
now carry more than before and run hotter. If a survivor was already near its
thermal rating, the added flow pushes it over, it trips, and its flow
redistributes again onto an even smaller set of lines. Each redistribution is
larger and faster than the last, so the cascade accelerates past the point
where operators or even automatic schemes can intervene, which is why the
$N-1$ standard requires the survivors to hold the full flow after any single
loss, with no line left near its limit. The mechanism is the same whether the
first loss is a fault, an overload on a hot day that sags the conductor into a
tree, or a relay misoperation; the danger is the redistribution, not the first
loss.

% Figure: the mechanism of a cascading transmission overload. A generator
% feeds a load over three parallel corridors. When the first line trips, its
% flow redistributes onto the survivors; if they are already near their
% thermal limit the next overloads and trips, and the redistribution repeats
% until the corridor collapses. The three panels are the same network at
% three moments: intact, one line lost, two lines lost.
\begin{figure}[htbp]
\centering
\begin{tikzpicture}[
  gen/.style={circle, draw=engblue, very thick, minimum size=10mm, font=\scriptsize},
  load/.style={rectangle, draw=engred, very thick, minimum size=8mm, font=\scriptsize},
  intact/.style={engblue, thick},
  tripped/.style={enggray, thick, dash pattern=on 1pt off 2pt},
  hot/.style={engorange, very thick},
  lbl/.style={font=\scriptsize\itshape, color=enggray},
]
% ----- panel 1: intact -----
\begin{scope}[xshift=0cm]
  \node[gen] (g1) at (0,0) {G};
  \node[load] (l1) at (3.4,0) {L};
  \draw[intact] (g1) to[out=35,in=145] node[above,font=\scriptsize]{$1.0$} (l1);
  \draw[intact] (g1) -- node[above,font=\scriptsize]{$1.0$} (l1);
  \draw[intact] (g1) to[out=-35,in=-145] node[below,font=\scriptsize]{$1.0$} (l1);
  \node[lbl] at (1.7,-1.6) {intact: 3 lines};
\end{scope}
% ----- panel 2: one line lost -----
\begin{scope}[xshift=5.6cm]
  \node[gen] (g2) at (0,0) {G};
  \node[load] (l2) at (3.4,0) {L};
  \draw[tripped] (g2) to[out=35,in=145] node[above,font=\scriptsize]{trip} (l2);
  \draw[hot] (g2) -- node[above,font=\scriptsize]{$1.5$} (l2);
  \draw[hot] (g2) to[out=-35,in=-145] node[below,font=\scriptsize]{$1.5$} (l2);
  \node[lbl] at (1.7,-1.6) {overload on survivors};
\end{scope}
% ----- panel 3: two lines lost -----
\begin{scope}[xshift=11.2cm]
  \node[gen] (g3) at (0,0) {G};
  \node[load] (l3) at (3.4,0) {L};
  \draw[tripped] (g3) to[out=35,in=145] (l3);
  \draw[tripped] (g3) -- (l3);
  \draw[tripped] (g3) to[out=-35,in=-145] node[below,font=\scriptsize]{trip} (l3);
  \node[lbl] at (1.7,-1.6) {corridor collapse};
\end{scope}
\end{tikzpicture}
\caption[The cascading-overload mechanism]{The mechanism of a cascading
transmission overload, the same corridor at three moments. Left: a generator
$\mathrm{G}$ feeds a load $\mathrm{L}$ over three parallel lines, each
carrying its share (per-unit flow shown). Centre: the top line trips, and its
flow redistributes onto the two survivors, which now carry $1.5$ times their
former load and run hot (orange), past their thermal rating. Right: an
overloaded survivor trips in turn, the flow redistributes again onto the last
line, which cannot carry the whole load, and the corridor collapses. Each
redistribution is faster than the last, which is why a cascade outruns
operator intervention once it starts and why the $N-1$ contingency standard
requires the survivors to hold the full flow after any single loss.}
\label{fig:vol04ch14:cascade-network}
\end{figure}


The engineering content of blackout prevention is twofold: build the
grid with operational margin (the $N-1$ contingency standard, in
which the grid must continue to operate after the loss of any single
element), and design the protective-relay system to disconnect
faulted elements rapidly without triggering cascading trips of
unfaulted elements. The IEEE C37 series of protective-relay standards
and the NERC operating standards (in North America, with analogous
ENTSO-E rules in Europe) specify the working framework.

\engterm{Three Mile Island} (28 March 1979) is the canonical reactor
incident in this chapter; the registry entry is at
\cite{acc:kemeny-tmi-1979}. A reactor scram during normal operation
combined with operator confusion about valve states led to partial
core uncovering and significant fuel damage. Off-site radiation
release was minor; the accident's consequences were primarily
economic (loss of the plant) and political (a decades-long pause in
new US nuclear plant orders). The engineering lessons centred on
control-room human factors: simpler indications, clearer alarms, and
training in the diagnosis of multiple simultaneous abnormalities are
the operational content of post-TMI nuclear-plant design.

\engterm{Thermal runaway in lithium-ion batteries} is the failure
mode at the energy-storage scale. A cell that experiences an
internal short circuit (from a manufacturing defect, mechanical
abuse, or extreme overcharge) heats; above $\sim 130\,\si{\celsius}$
the SEI layer in lithium-ion cells decomposes; above $\sim
200\,\si{\celsius}$ the cathode releases oxygen, which reacts with the
electrolyte; the temperature rises rapidly, the cell vents, and
adjacent cells can be propagated to runaway. The engineering response
is multi-layered: cell-level current and thermal sensing with
automatic disconnection; module-level fire-suppression; system-level
fire-rated installation with thermal isolation between modules; and
operational limits on charging rate and depth of discharge that
extend cycle life and reduce the likelihood of internal short
circuits. The McMicken battery storage facility fire (Arizona, 2019)
and the Moss Landing battery facility incidents (California, 2021
and 2024) are the working case studies; neither has a registry entry
in this book yet.

\engterm{Generation outages} are a routine failure mode whose
operational impact varies with the grid's reserve margin. A
$1\,\si{\giga\watt}$ unit trip in a $100\,\si{\giga\watt}$ grid is a
$1\,\%$ load-supply imbalance; the frequency dips by $\sim 0.05\,\si{
\hertz}$ and primary response restores balance within seconds.
A $1\,\si{\giga\watt}$ trip in a $10\,\si{\giga\watt}$ grid is a
$10\,\%$ imbalance; protective load-shedding may be triggered. The
$N-1$ standard ensures that any single trip is tolerable; multiple
correlated failures (e.g., a cold-weather event that simultaneously
de-rates many gas plants) violates the assumed independence and
produces the rarer but more severe failure modes.

\begin{principle}[The grid is a control system at every time scale]
The synchronous AC grid is a real-time control system operating on
time scales from microseconds (the protection-relay response) to
years (long-term planning). Failure to maintain control at any one
time scale propagates to the others: a generation-load imbalance not
caught at the second-scale by governor response triggers second-
scale corrections, which not caught at the minute scale trigger
minute-scale dispatch, which not caught at the hour scale trigger
hour-scale market response, which not caught at the day-to-year
scale fails to provision adequate capacity. The discipline of
power-systems engineering is the discipline of controlled cascades
across these time scales.
\end{principle}

\section{Project}\pathtag{core}

\begin{project}[Model a household with rooftop solar and batteries]
Build a spreadsheet (or programming) model of a household's annual
electricity consumption, rooftop PV generation, and battery storage,
and estimate the grid dependence.

\textbf{Data inputs.}
\begin{itemize}[leftmargin=*]
  \item Hourly household load: a representative residential load
    profile from the reader's utility (US: EIA RBSA; Europe: hourly
    profiles from the local DSO; many utilities publish anonymised
    hourly data).
  \item Hourly solar irradiance: NREL PVWatts or the EU Joint
    Research Centre PVGIS, free and well-calibrated.
  \item Battery specification: round-trip efficiency $90\,\%$, usable
    capacity $10\,\si{\kilo\watt\hour}$, peak charge/discharge power
    $5\,\si{\kilo\watt}$.
\end{itemize}

\textbf{Procedure.}
\begin{enumerate}
  \item Compute the hourly PV generation from the irradiance data
    and a rated DC capacity of $5\,\si{\kilo\watt\textsubscript{p}}$,
    with $0.86$ DC-to-AC conversion efficiency and azimuth and tilt
    optimised for the reader's latitude.
  \item Subtract the hourly load from the hourly generation to
    compute the hourly surplus (positive) or deficit (negative).
  \item Apply battery dispatch: charge from surplus until full,
    discharge from deficit until empty. The greedy dispatch in
    \texttt{code/household\_dispatch.py}, run on the sample profile in
    \texttt{data/household\_load\_pv.csv}, implements exactly this loop
    and reports the annual fractions.
  \item Sum the residual surplus (sold back to grid or curtailed)
    and the residual deficit (imported from grid).
  \item Repeat for a year of data; report the annual grid-import
    fraction, the annual self-consumption fraction, and the
    seasonal pattern.
\end{enumerate}

\textbf{Report.} A one-page report with a sample-week generation-
load-battery plot, the annual statistics, and a sensitivity analysis
to: PV capacity (vary from $3$ to $10\,\si{\kilo\watt
\textsubscript{p}}$), battery capacity (vary from $0$ to $30\,\si{
\kilo\watt\hour}$), and load profile (compare a single-occupant flat
to a four-person household). Expected: an annual self-consumption
fraction of $\sim 70$ to $80\,\%$ for the base case in mid-latitude
locations; the marginal benefit of additional battery capacity
saturates beyond $\sim 15\,\si{\kilo\watt\hour}$. The household
remains grid-dependent in winter unless the PV is significantly
oversized.
\end{project}

\section{Exercises}

\subsection*{Calculation}\pathtag{core}

\begin{exercise}
A $400\,\si{\kilo\volt}$ three-phase transmission line carries
$1\,\si{\giga\watt}$ of real power at power factor $0.95$. Compute
the line current and the $I^{2}R$ loss in a $200\,\si{\kilo\meter}$
line of resistance $0.03\,\si{\ohm\per\kilo\meter}$ per phase.
\paragraph{Solution sketch.}
From $P = \sqrt{3}\,V_L I_L \cos\phi$, the line current is $I_L =
P/(\sqrt{3}\,V_L\cos\phi) = 10^9/(\sqrt{3}\cdot 4\times10^5\cdot 0.95)
\approx 1519\,\si{\ampere}$. The per-phase resistance is $R = 0.03
\times 200 = 6\,\si{\ohm}$. The three-phase loss is $P_{\text{loss}} =
3 I_L^2 R = 3(1519)^2(6) \approx 41.5\,\si{\mega\watt}$, about $4.2\,\%$
of the transmitted power. The point is the $1/V^2$ scaling of loss at
fixed power: doubling the voltage quarters the current and the loss,
which is why bulk power moves at hundreds of kilovolts.
\end{exercise}

\begin{exercise}
A wind turbine has a rotor diameter of $120\,\si{\meter}$ and operates
at $C_{p} = 0.45$ in wind of $10\,\si{\meter\per\second}$ at sea-
level air density. Compute the mechanical power available at the
rotor shaft.
\paragraph{Solution sketch.}
The swept area is $A = \pi(60)^2 \approx 1.13\times10^4\,\si{\meter
\squared}$. With $\rho = 1.225\,\si{\kilogram\per\meter\cubed}$ and
$v = 10\,\si{\meter\per\second}$, $P = \tfrac{1}{2}C_p\rho A v^3 =
0.5(0.45)(1.225)(1.13\times10^4)(10^3) \approx 3.1\,\si{\mega\watt}$.
Note the cube law: dropping the wind to $8\,\si{\meter\per\second}$
would cut the power to $(8/10)^3 = 0.51$ of this, about $1.6\,\si{\mega
\watt}$.
\end{exercise}

\begin{exercise}
A pumped-storage plant has an upper reservoir at $500\,\si{\meter}$
elevation above the lower reservoir, $80\,\%$ round-trip efficiency,
and stores $2\,\si{\giga\watt\hour}$ of electrical energy. Compute
the volume of water cycled at full discharge.
\paragraph{Solution sketch.}
The recoverable electrical energy is $E = \eta_{\text{gen}}\,\rho g V H$.
Splitting the $80\,\%$ round trip evenly gives a generation-leg
efficiency $\eta_{\text{gen}} \approx \sqrt{0.80} \approx 0.89$ (or take
the stated value as the discharge efficiency directly; either reading is
acceptable if stated). Using $\eta_{\text{gen}} = 0.89$, $E = 2\,\si{
\giga\watt\hour} = 7.2\times10^{12}\,\si{\joule}$, so $V = E/(\eta\rho g
H) = 7.2\times10^{12}/(0.89\cdot1000\cdot9.81\cdot500) \approx 1.65
\times10^{6}\,\si{\meter\cubed}$, roughly $1.6\,\si{\meter}$ of
depth over a square kilometre, or a reservoir a few hundred metres on a
side and several metres deep. The head, not the volume, is what makes
pumped storage compact.
\end{exercise}

\begin{exercise}
A residential rooftop PV system has $6\,\si{\kilo\watt
\textsubscript{p}}$ rated DC, $85\,\%$ DC-to-AC efficiency, and a
$22\,\%$ capacity factor. Compute the annual electricity production
in kilowatt-hours.
\paragraph{Solution sketch.}
Read the capacity factor as applying to the AC nameplate after the
inverter, so $P_{\text{AC,rated}} = 0.85\times 6 = 5.1\,\si{\kilo\watt}$
and $E = P_{\text{AC,rated}}\,\text{CF}\,(8760) = 5.1\times0.22\times
8760 \approx 9830\,\si{\kilo\watt\hour}/\text{yr}$. If instead the
$22\,\%$ already refers to the DC rating, $E = 6\times0.22\times8760
\approx 11600\,\si{\kilo\watt\hour}$ before the AC conversion. Either
convention gives $\sim 10\,\si{\mega\watt\hour}/\text{yr}$; state which
capacity the capacity factor is referenced to, since that choice is the
common source of disagreement between quotes.
\end{exercise}

\begin{exercise}
A lithium-iron-phosphate battery has a nominal capacity of $200\,\si{
\ampere\hour}$ at $3.2\,\si{\volt}$ per cell and a $90\,\%$ round-
trip efficiency. Compute the energy stored and the charging energy
required from the grid.
\paragraph{Solution sketch.}
Per cell, $E = QV = 200\,\si{\ampere\hour}\times3.2\,\si{\volt} = 640\,
\si{\watt\hour} = 0.64\,\si{\kilo\watt\hour}$. The energy delivered on
discharge is what the cell holds; the charging energy drawn from the
grid to return that much to load over a full cycle is $E_{\text{in}} =
E_{\text{out}}/\eta_{\text{rt}} = 0.64/0.90 \approx 0.71\,\si{\kilo\watt
\hour}$. The $\sim 11\,\%$ difference is the round-trip loss, dissipated
as heat in the cell internal resistance and the power electronics.
\end{exercise}

\begin{exercise}
A grid is supplying $80\,\si{\giga\watt}$ of demand from a generator
pool with a primary frequency response of $20\,\si{\giga\watt\per
\hertz}$. A $1\,\si{\giga\watt}$ generator trips. Compute the
steady-state frequency dip.
\paragraph{Solution sketch.}
The primary response stiffness relates the steady-state power deficit
to the frequency error: $\Delta P = K\,\Delta f$, so $\Delta f =
\Delta P / K = 1\,\si{\giga\watt}/20\,\si{\giga\watt\per\hertz} = 0.05\,
\si{\hertz}$. The grid settles at $49.95\,\si{\hertz}$ (or $59.95$ on a
$60\,\si{\hertz}$ grid), inside the $\pm 0.05\,\si{\hertz}$ steady-state
band; secondary control then trims the small residual back to nominal.
This is the steady-state value, not the transient nadir, which is deeper
and depends on inertia (the swing-equation model of Figure
\ref{fig:vol04ch14:frequency-response}).
\end{exercise}

\begin{exercise}
A combined-cycle gas turbine plant achieves $60\,\%$ thermal-to-
electric efficiency. Compute the natural-gas consumption in $\si{
\meter\cubed\per\second}$ for a $500\,\si{\mega\watt}$ output, given
a lower heating value of $36\,\si{\mega\joule\per\meter\cubed}$.
\paragraph{Solution sketch.}
The fuel thermal input is $\dot{Q} = P_{\text{out}}/\eta = 500/0.60
\approx 833\,\si{\mega\watt}$. Dividing by the volumetric heating value,
$\dot{V} = \dot{Q}/\text{LHV} = 833\,\si{\mega\watt}/36\,\si{\mega\joule
\per\meter\cubed} \approx 23\,\si{\meter\cubed\per\second}$ of gas, or
about $2\times10^{6}\,\si{\meter\cubed}$ per day. The same plant as a
simple-cycle unit at $\eta \approx 0.38$ would burn $\sim 60\,\%$ more
gas for the same output, which is the case for the bottoming Rankine
cycle.
\end{exercise}

\begin{exercise}
A hydrogen-fuel-cell vehicle has a tank holding $5\,\si{\kilogram}$
of compressed hydrogen at $70\,\si{\mega\pascal}$, and a fuel cell
operating at $55\,\%$ chemical-to-electric efficiency. Compute the
total energy stored in the tank and the equivalent driving range at
$200\,\si{\watt\hour\per\kilo\meter}$ road-load.
\paragraph{Solution sketch.}
Hydrogen's lower heating value is $\sim 120\,\si{\mega\joule\per\kilogram}
\approx 33.3\,\si{\kilo\watt\hour\per\kilogram}$, so the chemical energy
is $5\times33.3 \approx 167\,\si{\kilo\watt\hour}$. The electrical
energy at the wheels (through the $55\,\%$ fuel cell) is $0.55\times167
\approx 92\,\si{\kilo\watt\hour}$. At $200\,\si{\watt\hour\per\kilo
\meter}$, the range is $92000/200 \approx 460\,\si{\kilo\meter}$. The
$5\,\si{\kilogram}$ of hydrogen carries roughly the usable energy of a
$60\,\si{\kilo\watt\hour}$ battery pack but at a fraction of the mass,
which is the case for hydrogen in heavy long-range transport.
\end{exercise}

\subsection*{Derivation}\pathtag{standard}

\begin{exercise}
Derive the Betz limit $C_{p,\max} = 16/27$ for the power extracted
from a horizontal-axis wind turbine, by treating the rotor as a
disc actuator that decelerates the upstream-to-downstream flow.
\paragraph{Solution sketch.}
Let the upstream speed be $v_1$ and the far-downstream speed $v_2$, and
define the axial induction factor $a$ so the disc speed is $v_1(1-a)$.
Mass conservation gives the disc velocity as the average $\tfrac12(v_1
+v_2)$, so $v_2 = v_1(1-2a)$. The power extracted equals the kinetic-
energy flux lost by the air, $P = \tfrac12\dot{m}(v_1^2 - v_2^2)$ with
$\dot{m} = \rho A v_1(1-a)$. Substituting $v_2 = v_1(1-2a)$ and
factoring out $\tfrac12\rho A v_1^3$ gives $P = \tfrac12\rho A
v_1^3\,[\,4a(1-a)^2\,]$, so $C_p = 4a(1-a)^2$. Differentiating,
$dC_p/da = 4(1-a)(1-3a) = 0$ at $a = 1/3$, where $C_p = 4(\tfrac13)
(\tfrac23)^2 = 16/27 \approx 0.593$. Figure
\ref{fig:vol04ch14:betz-curve} is this curve.
\end{exercise}

\begin{exercise}
Derive the optimum tilt angle for a fixed-tilt solar panel as a
function of latitude, given simple geometric arguments about
average noon solar elevation over the year.
\paragraph{Solution sketch.}
At solar noon the sun's elevation above the horizon is $90^\circ -
\phi + \delta$, where $\phi$ is the latitude and $\delta$ is the solar
declination, swinging $\pm 23.4^\circ$ over the year and averaging zero.
A panel collects most when its normal points at the sun, that is when
its tilt from horizontal equals the sun's zenith angle $\phi - \delta$.
Averaging over the year ($\langle\delta\rangle = 0$) gives the
first-order rule $\beta_{\text{opt}} \approx \phi$: tilt the panel at
the latitude angle, facing the equator. The refinement is that
energy-weighting toward the brighter summer months and accounting for
diffuse light pulls the optimum a few degrees below $\phi$ (empirically
$\beta \approx 0.76\phi + 3.1^\circ$ is a common fit); the geometric
$\beta = \phi$ is the defensible starting point.
\end{exercise}

\begin{exercise}
Derive the relationship between battery state of charge and
open-circuit voltage for a lithium-iron-phosphate cell, using a
simplified one-zone thermodynamic model. State the assumptions.
\paragraph{Solution sketch.}
The open-circuit voltage is the Gibbs-energy change per unit charge of
the cell reaction, $V_{\text{oc}} = -\Delta G/(nF)$. Writing $\Delta G
= \Delta G^\circ + RT\ln Q$ with $Q$ the reaction quotient gives a
Nernst form $V_{\text{oc}} = V^\circ - (RT/nF)\ln Q$, and modelling the
electrode as a single solid-solution zone whose lithium activity tracks
the state of charge $s$ as $Q \sim s/(1-s)$ yields $V_{\text{oc}}(s) =
V^\circ - (RT/F)\ln[s/(1-s)]$, a sigmoid that is steep near the ends and
flat in the middle. The flat middle is exactly what LFP shows: a long
voltage plateau near $3.2\,\si{\volt}$ that makes state-of-charge
estimation from voltage alone hard, forcing coulomb-counting instead.
Assumptions: single-phase solid solution (LFP is really two-phase, which
flattens the plateau further), uniform temperature, no kinetic
overpotential, equilibrium measurement.
\end{exercise}

\begin{exercise}
Derive the $\sqrt{3}$ factor in three-phase power: $P = \sqrt{3}V_{L}
I_{L}\cos\phi$, by summing the per-phase contributions and
applying the line-to-phase voltage and current relationships for a
wye-connected load.
\paragraph{Solution sketch.}
Total power is three times the per-phase power, $P = 3 V_{\text{ph}}
I_{\text{ph}}\cos\phi$, where the per-phase quantities are the
phase-to-neutral voltage and the phase current. In a wye connection the
line current equals the phase current, $I_L = I_{\text{ph}}$, while the
line-to-line voltage is the phasor difference of two phase voltages
$120^\circ$ apart, whose magnitude is $V_L = \sqrt{3}\,V_{\text{ph}}$
(the chord of two unit vectors at $120^\circ$ is $\sqrt{3}$). Substitute
$V_{\text{ph}} = V_L/\sqrt{3}$ and $I_{\text{ph}} = I_L$ into $P = 3
V_{\text{ph}}I_{\text{ph}}\cos\phi$ to get $P = 3(V_L/\sqrt3)I_L\cos\phi
= \sqrt{3}\,V_L I_L\cos\phi$. The $\sqrt3$ is the geometry of three
phasors at $120^\circ$, not a property of power.
\end{exercise}

\subsection*{Estimation}\pathtag{core}

\begin{exercise}
Estimate the rooftop solar capacity required to meet the global
electricity demand if no other generation were used. State
assumptions about average insolation, panel efficiency, and the
load factor.
\paragraph{Solution sketch.}
Global electricity demand is $\sim 30000\,\text{TWh}/\text{yr}$
(current as of 2024). With a solar capacity factor of $\sim 15\,\%$
(a global rooftop average), the required nameplate is $P =
E/(\text{CF}\times8760) = 30000\times10^9/(0.15\times8760) \approx
2.3\times10^{13}\,\si{\watt} \approx 23\,\si{\tera\watt}$ of panels,
roughly twenty times the $\sim 1.5\,\si{\tera\watt}$ installed by 2024.
At $\sim 200\,\si{\watt\per\meter\squared}$ peak, that is $\sim 1.1
\times10^{11}\,\si{\meter\squared} \approx 1.1\times10^{5}\,\si{\kilo
\meter\squared}$ of panel area. The estimate ignores storage; supplying
demand at night and in winter from daytime summer sun is the real
constraint, and it multiplies the storage and overbuild requirement
several-fold.
\end{exercise}

\begin{exercise}
Estimate the lithium-ion battery capacity that would be needed to
buffer the United Kingdom's electricity grid through a
$10$-day still-and-cloudy weather event, given national peak demand
of $\sim 45\,\si{\giga\watt}$. State the capital cost at $\$140\,\si{
\per\kilo\watt\hour}$ and compare to the UK's annual electricity
expenditure.
\paragraph{Solution sketch.}
Take the average load over the event as $\sim 70\,\%$ of peak, $\sim
31\,\si{\giga\watt}$. The energy is $31\,\si{\giga\watt}\times24\,\si{
\hour}\times10 \approx 7.6\,\text{TWh} = 7.6\times10^{9}\,\si{\kilo\watt
\hour}$. At $\$140\,\si{\per\kilo\watt\hour}$ the capital cost is
$\sim \$1.06$ trillion (current as of 2024), several times the UK's
annual electricity expenditure of order $\$50$ to $\$100$ billion. The
verdict is that lithium-ion is the wrong technology for multi-day
seasonal buffering; the role belongs to firm low-carbon generation,
interconnection, or long-duration storage with a low energy cost. The
\texttt{code/storage\_sizing.py} routine returns these numbers.
\end{exercise}

\begin{exercise}
Estimate the land area required for a wind farm to produce the same
annual energy as a $1\,\si{\giga\watt}$ nuclear plant, given turbine
spacing of $5$ rotor diameters and a capacity factor of $35\,\%$.
\paragraph{Solution sketch.}
The nuclear plant at a $92\,\%$ capacity factor delivers $\sim 8.1\,
\text{TWh}/\text{yr}$. Matching this with $35\,\%$-capacity-factor
wind needs $8.1/(0.35\times8760)\times10^6 \approx 2.6\,\si{\giga\watt}$
nameplate. Take $4\,\si{\mega\watt}$ turbines of $150\,\si{\meter}$
rotor, so $\sim 660$ turbines. At $5$-diameter spacing each occupies
$\sim (5\times150)^2 = 5.6\times10^{5}\,\si{\meter\squared}$, giving a
gross farm area $\sim 660\times5.6\times10^{5} \approx 370\,\si{\kilo
\meter\squared}$, a square $\sim 19\,\si{\kilo\meter}$ on a side. The
land between turbines remains usable for farming or grazing, so the
exclusive footprint is far smaller than the gross area; the nuclear
plant's own site is $\sim 1\,\si{\kilo\meter\squared}$. The contrast in
land intensity is one of the standard trade-offs.
\end{exercise}

\begin{exercise}
Estimate the round-trip efficiency of a green-hydrogen storage chain
from electricity through electrolysis, compression, transport, and
fuel-cell conversion. State the dominant loss at each stage.
\paragraph{Solution sketch.}
Multiply the stage efficiencies: electrolysis $\sim 0.65$ (the
overpotential needed to split water faster than equilibrium),
compression to $70\,\si{\mega\pascal}$ $\sim 0.90$ (the compressor work,
$\sim 10\,\%$ of the hydrogen's energy), transport and storage $\sim
0.95$ (boil-off and leakage), fuel cell $\sim 0.55$ (the activation and
ohmic losses, plus the thermodynamic ceiling). The product is $0.65
\times0.90\times0.95\times0.55 \approx 0.31$, so roughly two thirds of
the input electricity is lost. The electrolyser and the fuel cell are
the two dominant losses. The low figure is why hydrogen storage is
reserved for long-duration and hard-to-electrify uses where the
round-trip penalty is paid rarely; \texttt{code/storage\_sizing.py}
computes the chain product.
\end{exercise}

\begin{exercise}
Estimate the timing requirement of grid-frequency protection relays
to ride through a $50\,\si{\hertz}\pm 0.5\,\si{\hertz}$ excursion in
a synchronous-AC grid. State the consequence of a poorly chosen
trip-time for grid resilience.
\paragraph{Solution sketch.}
A relay must distinguish a recoverable transient excursion from a
genuine collapse. The nadir of a normal trip is reached in $\sim 5$ to
$10\,\si{\second}$ and recovers over tens of seconds (Figure
\ref{fig:vol04ch14:frequency-response}), so an under-frequency relay set
to trip at $49.5\,\si{\hertz}$ should carry a deliberate time delay of
order seconds (or be set well below the expected nadir) to ride through.
Set the trip too fast or too shallow and the relay sheds generation on a
survivable dip, deepening the very imbalance it was meant to protect
against and turning a transient into a cascade; set it too slow and
equipment is exposed to damaging off-frequency operation. The working
compromise is a frequency-time characteristic: progressively faster
tripping the further the frequency strays, the under-frequency
load-shedding schedule that protects the grid as a whole.
\end{exercise}

\subsection*{Simulation}\pathtag{mastery}

\begin{exercise}
Implement an hour-by-hour simulator of a 100$\,\%$-renewable
electricity system for a temperate-climate country, with a mix of
wind, solar, and pumped-storage hydropower. Vary the wind:solar:
storage ratio and identify the minimum-cost mix subject to a
$99.95\,\%$ reliability constraint.
\paragraph{Solution sketch.}
Reuse the dispatch loop of \texttt{code/household\_dispatch.py} at
national scale: at each hour compute generation $G_t = w\cdot
\text{wind}_t + s\cdot\text{solar}_t$ from normalised wind and solar
traces (the shape of Figure \ref{fig:vol04ch14:renewable-variability}),
subtract demand, and charge or discharge a storage reservoir bounded by
its energy and power ratings. Count the hours of unserved energy; the
reliability is one minus the unserved fraction. Sweep the
wind:solar:storage triple on a grid, keep the points meeting $99.95\,\%$,
and minimise total annualised cost (capital times capacity-recovery
factor plus running cost) over that feasible set. The recurring finding
is that a wind-heavy mix needs less storage in temperate climates
(wind's winter strength matches winter demand), that the last fraction
of a percent of reliability is expensive (a long tail of rare
simultaneous low-wind, low-sun events), and that a small dispatchable
backup is far cheaper than the storage to reach $100\,\%$ on renewables
alone.
\end{exercise}

\begin{exercise}
Implement a second-by-second simulator of frequency response in a
synchronous AC grid to a sudden generator trip. Vary the
inertia constant $H$ (kinetic energy in rotating machinery per
megawatt of generation capacity) and identify the regime in which
under-frequency load shedding is necessary.
\paragraph{Solution sketch.}
Integrate the swing equation with a governor lag, $2H\,d(\Delta f)/dt =
P_g - \Delta P - D\,\Delta f$ with $P_g$ relaxing toward $-\Delta f/R$
on a time constant $T_g$, exactly the model in
\texttt{code/frequency\_response.py}. Apply a step generation loss
$\Delta P$ and record the nadir. Sweeping $H$ shows the nadir deepening
as inertia falls: with $H = 8\,\si{\second}$ the dip is shallow, and as
$H$ drops toward $2\,\si{\second}$ (a high-inverter-penetration grid)
the nadir crosses the under-frequency load-shedding threshold for a
given $\Delta P$. The regime in which load shedding is necessary is the
low-inertia, large-trip corner; the engineering responses are synthetic
inertia from inverters, faster primary response, and limits on the size
of the largest single in-feed.
\end{exercise}

\subsection*{Design}\pathtag{standard}

\begin{exercise}
Design a microgrid for a remote village of $200$ households with
average daily demand of $5\,\si{\kilo\watt\hour}$ per household.
Specify the solar capacity, battery capacity, backup-generator size,
and the working operational reserve. Identify the LCOE.
\paragraph{Solution sketch.}
Daily energy is $200\times5 = 1000\,\si{\kilo\watt\hour}/\text{day}$.
At a $\sim 18\,\%$ solar capacity factor the PV nameplate is $1000/
(0.18\times24) \approx 230\,\si{\kilo\watt\textsubscript{p}}$; oversize
to $\sim 300\,\si{\kilo\watt\textsubscript{p}}$ to cover cloudy days and
charge storage. Size the battery for one to two days of autonomy,
$\sim 1$ to $2\,\si{\mega\watt\hour}$ usable. Specify a diesel backup at
the peak load ($\sim 1000\,\si{\kilo\watt\hour}/\text{day}$ implies an
evening peak of perhaps $150\,\si{\kilo\watt}$, so a $\sim 150\,\si{\kilo
\watt}$ genset) for the rare multi-day deficit, with a reserve margin of
$\sim 20\,\%$. The LCOE combines the annualised PV, battery, and genset
capital with the diesel fuel and O\&M over the delivered energy; for a
PV-battery-diesel hybrid it typically lands at $\$0.20$ to $\$0.40\,\si{
\per\kilo\watt\hour}$ (current as of 2024), well below diesel-only
generation. The \texttt{code/lcoe.py} structure carries the calculation.
\end{exercise}

\begin{exercise}
Design a residential-scale heat pump for a $150\,\si{\meter\squared}$
house at a $-10\,\si{\celsius}$ design outdoor temperature. Specify
the rated capacity, the COP at design conditions, the auxiliary
electric resistance backup, and the working refrigerant.
\paragraph{Solution sketch.}
Estimate the design heat loss from a specific load of $\sim 50\,\si{
\watt\per\meter\squared}$ for a moderately insulated house at a
$\sim 30\,\si{\kelvin}$ indoor-outdoor difference, giving $\sim 7.5\,\si{
\kilo\watt}$. Size the heat pump at $\sim 8\,\si{\kilo\watt}$ thermal
output at design conditions. The Carnot ceiling between $T_h = 293\,\si{
\kelvin}$ and $T_c = 263\,\si{\kelvin}$ is $293/30 \approx 9.8$; real
air-source units reach $\sim 40$ to $50\,\%$ of Carnot, so a COP of
$\sim 2.5$ to $3$ at $-10\,\si{\celsius}$, drawing $\sim 3\,\si{\kilo
\watt}$ electrical. Provide $\sim 5\,\si{\kilo\watt}$ of resistance
backup for the coldest hours and defrost cycles. Modern units use
low-global-warming-potential refrigerants such as R-32 or R-290
(propane); the choice trades efficiency, flammability, and regulatory
phase-down. The lesson is that the COP makes the heat pump deliver more
thermal energy than it consumes electrically, which is what makes
electrified heat efficient.
\end{exercise}

\begin{exercise}
Design an EV-charging infrastructure for a $100$-vehicle commercial
fleet that returns to depot at $18{:}00$ and departs at $07{:}00$
the following morning. Specify the per-vehicle charging power, the
total facility capacity, the demand-side management strategy, and the
working interconnection requirements.
\paragraph{Solution sketch.}
Suppose each vehicle needs $\sim 60\,\si{\kilo\watt\hour}$ overnight. The
window is $13\,\si{\hour}$, so the energy-limited per-vehicle power is
$60/13 \approx 4.6\,\si{\kilo\watt}$, and $100$ vehicles charging flat
out at once would demand $\sim 460\,\si{\kilo\watt}$. The naive design
sizes every charger at, say, $11\,\si{\kilo\watt}$ and the service at
$1.1\,\si{\mega\watt}$, an expensive interconnection. The
demand-side-management design instead uses smart charging: stagger and
modulate the chargers so the aggregate stays near the $\sim 460\,\si{
\kilo\watt}$ energy floor, flattening the depot's load and shrinking the
interconnection to $\sim 0.5\,\si{\mega\watt}$. Add managed charging to
shift load to off-peak tariff hours and, where available,
vehicle-to-grid discharge for grid services. Interconnection
requirements include the service transformer rating, the metering and
protection per the local standard (IEEE 1547 for any export), and
coordination with the utility for the new load. The design lesson is
that coincident peak, not total energy, sets the cost, and management
controls the coincident peak.
\end{exercise}

\begin{exercise}
Design a pumped-storage facility of $1\,\si{\giga\watt}/8\,\si{\hour}$
duration. Specify the reservoir volume, the elevation difference, the
turbine type and size, and the round-trip efficiency at full
operation.
\paragraph{Solution sketch.}
The energy capacity is $1\,\si{\giga\watt}\times8\,\si{\hour} = 8\,\si{
\giga\watt\hour} = 2.88\times10^{13}\,\si{\joule}$. Choose a head: at
$H = 400\,\si{\meter}$ and a generation efficiency $\sim 0.90$, the
usable volume is $V = E/(\eta\rho g H) = 2.88\times10^{13}/(0.90\cdot
1000\cdot9.81\cdot400) \approx 8.2\times10^{6}\,\si{\meter\cubed}$,
roughly a reservoir $\sim 600\,\si{\meter}$ square and $\sim 23\,\si{
\meter}$ deep, replicated as the lower reservoir. The flow at full
$1\,\si{\giga\watt}$ output is $Q = P/(\eta\rho g H) \approx 283\,\si{
\meter\cubed\per\second}$. At this head a Francis pump-turbine is the
standard choice (Pelton above $\sim 600\,\si{\meter}$, Kaplan below
$\sim 50\,\si{\meter}$); size for $\sim 283\,\si{\meter\cubed\per\second}$
across one or more units. Round-trip efficiency is the product of the
pumping and generating legs, $\sim 0.90\times0.88 \approx 0.79$. The
head sets the compactness; the volume scales as $1/H$.
\end{exercise}

\subsection*{Diagnosis and reverse engineering}\pathtag{standard}

\begin{exercise}
A residential solar-PV system shows annual energy production $30\,\%$
below the predicted yield. Identify three candidate causes (shading,
inverter degradation, panel soiling) and the supplementary
measurement that distinguishes them.
\paragraph{Solution sketch.}
Each cause leaves a distinct signature. Shading produces
time-of-day-dependent losses (a notch at the hours when a tree or
chimney casts a shadow) and, on string inverters, a large drop when one
shaded panel pulls down its string; the distinguishing measurement is
per-string or per-panel output versus time of day, or a shading survey.
Inverter degradation or clipping shows as a flat-topped midday curve or
elevated inverter temperature and fault logs; read the inverter
efficiency and DC-versus-AC power. Soiling produces a slow uniform
decline that resets after rain or cleaning; the distinguishing test is a
clean-panel-versus-dirty-panel comparison or the step recovery after a
wash. Compare measured plane-of-array irradiance with output to
separate a resource shortfall from a system fault.
\end{exercise}

\begin{exercise}
A grid-scale battery shows accelerated capacity fade $20\,\%$ over
two years of operation when the warranty specification is $20\,\%$
over ten years. Identify candidate causes (excessive cycling, depth
of discharge beyond specification, thermal-management failure) and
the working diagnostic measurements for each.
\paragraph{Solution sketch.}
A five-fold-fast fade points at the operating envelope. Excessive
cycling: count the equivalent full cycles in the energy-throughput logs
and compare to the warranted cycle count; fade scales with throughput.
Depth of discharge beyond specification: read the state-of-charge
histogram, since deep cycling to near $0\,\%$ or charging to near
$100\,\%$ accelerates degradation far more than mid-range cycling.
Thermal-management failure: read the cell-temperature logs, since
sustained operation above $\sim 35\,\si{\celsius}$ roughly doubles the
calendar-fade rate per $\sim 10\,\si{\kelvin}$ (an Arrhenius
dependence). The single most useful measurement is a controlled
capacity test (full slow charge-discharge) to confirm the fade, plus
the temperature and state-of-charge histories to assign the cause.
\end{exercise}

\begin{exercise}
A wind turbine produces $20\,\%$ less power than its rated curve at
the same wind speed. Identify three candidate causes (rotor pitch
miscalibration, blade contamination, gearbox slip) and the
maintenance procedure for each.
\paragraph{Solution sketch.}
Pitch miscalibration moves the blades off their optimum angle of
attack, lowering $C_p$ across the whole curve; check the pitch-encoder
calibration against a reference and re-zero the actuators. Blade
contamination (insect or ice accretion, leading-edge erosion)
roughens the aerofoil and degrades lift, worst at high wind; inspect
and clean or repair the leading edge. Gearbox or generator slip and
control faults show as a mismatch between rotor speed and electrical
output or as rising drivetrain temperature and vibration; read the
SCADA speed-versus-power data and the condition-monitoring vibration
spectrum. The discriminating step is to plot the measured power curve
against the manufacturer's reference and look at where in the
wind-speed range the deficit sits: a uniform shift suggests pitch, a
high-wind shift suggests aerodynamic degradation, and a speed-power
mismatch suggests the drivetrain.
\end{exercise}

\subsection*{Failure analysis}\pathtag{core}

\begin{exercise}
Summarise the failure mechanism of Three Mile Island Unit 2
(28 March 1979) in three sentences, using the registry entry's
working language at \cite{acc:kemeny-tmi-1979}. Identify the human-
factors lesson that has shaped subsequent nuclear-plant control-room
design.
\paragraph{Solution sketch.}
A model answer follows the registry's reconstruction. A loss of
feedwater raised primary pressure, the pilot-operated relief valve
opened as designed and then stuck open, and primary coolant escaped
while the control-room indicator showed only the valve's command signal
(closed) rather than its actual position (open). Reading the rising
pressuriser level as the system ``going solid'' with water, operators
reduced emergency-core-cooling flow, allowing the core to be partially
uncovered and damaged over several hours. The human-factors lesson is
that control rooms must display actual plant state rather than commanded
state, with clearer alarms and training in diagnosing multiple
simultaneous abnormalities; this reshaped subsequent control-room design
and operator training.
\end{exercise}

\begin{exercise}
A grid-tied PV inverter trips during a brief grid frequency
excursion to $50.3\,\si{\hertz}$. Identify whether the trip is
consistent with the IEEE 1547 frequency-tolerance specification
and the operational implication for high-PV-penetration grids.
\paragraph{Solution sketch.}
A trip at $50.3\,\si{\hertz}$ (a $0.3\,\si{\hertz}$ over-frequency, on a
$50\,\si{\hertz}$ grid) is well inside the normal operating band, and a
modern inverter built to the current IEEE 1547 ride-through requirements
should not trip there; it should remain connected and, under the newer
standard, provide frequency-watt response by curtailing output. A trip
at that excursion is consistent with an older fixed-threshold setting,
not with the present ride-through specification. The operational
implication is the case for ride-through and grid-supporting functions:
on a high-PV grid, a narrow trip setting means a small frequency
disturbance can disconnect a large block of generation at once,
converting a minor excursion into a large imbalance, the inverter
analogue of the cascading-trip mechanism. The fix is mandated
ride-through and frequency-responsive curtailment.
\end{exercise}

\begin{exercise}
A battery-storage facility experiences a thermal runaway that
propagates from one $5\,\si{\mega\watt\hour}$ module to two adjacent
modules. Identify the engineering controls that should have
prevented propagation (thermal isolation, gas venting, sprinkler
suppression, module-level shutdown) and the IEC 62933 standard that
defines them.
\paragraph{Solution sketch.}
Propagation means the defence-in-depth failed at the module-to-module
boundary. The controls that should have contained the event: thermal
isolation (fire-rated barriers and spacing between modules so a runaway
cell cannot heat its neighbours past the SEI-decomposition threshold);
gas venting and deflagration management (lithium-ion runaway releases
flammable gas, so the enclosure must vent or inert it to prevent
explosion); suppression (water or clean-agent systems sized for the
energy released); and module-level electrical shutdown and current
limiting on the first sign of over-temperature. The relevant
safety-design framework is the IEC 62933 series for electrical
energy-storage systems (with IEC 62933-5-2 on safety), alongside the
installation-fire standards. The lesson is that single-cell runaway is
treated as a certainty to be contained, not an event to be prevented;
the design goal is to stop the cascade at the module boundary.
\end{exercise}

\begin{exercise}
A regional grid experiences a cascading blackout during a
heat wave when one transmission line trips on overload and adjacent
lines trip in succession. Identify the failure mechanism and the
$N-1$-and-beyond planning standard that should have prevented it.
\paragraph{Solution sketch.}
The mechanism is load redistribution under overload. When the first
line trips, its power redistributes onto parallel paths; if those are
already near their thermal limits (worsened by the heat wave, which both
raises demand and lowers conductor ampacity through sag and reduced
cooling), the next line overloads and trips, redistributing again, and
the cascade accelerates past the operators' ability to intervene. The
$N-1$ standard requires the grid to survive the loss of any single
element with all others within limits; the heat wave violated the
assumptions behind it by stressing many elements simultaneously and
correlating their margins. The prevention is operating with real
contingency margin (not just nominal $N-1$ on a hot day), dynamic line
ratings that account for weather, and protection schemes that shed load
in a controlled way before uncontrolled cascading begins. This is the
mechanism behind the 2003 Northeast event described in the failure
section.
\end{exercise}

\subsection*{Judgment}\pathtag{standard}

\begin{exercise}
A consulting client must decide whether to invest in $1\,\si{\giga
\watt}$ of new gas-fired generation or $1\,\si{\giga\watt}$ of wind
plus $400\,\si{\mega\watt}/4\,\si{\hour}$ of battery storage.
Identify the trade-off in terms of LCOE, dispatchability, carbon
emissions, capital cost, and operational flexibility.
\paragraph{Solution sketch.}
The gas plant is dispatchable: it delivers $1\,\si{\giga\watt}$ on
demand whenever fuel is available, with a moderate capital cost and a
levelized cost set largely by gas price and carbon cost, and it emits
$\sim 350$ to $450\,\si{\kilogram}\,\text{CO}_2$ per MWh. The
wind-plus-battery option has a lower energy LCOE and near-zero
operating emissions, but the $1\,\si{\giga\watt}$ of wind delivers only
$\sim 35\,\%$ of nameplate on average and the $400\,\si{\mega\watt}/4\,
\si{\hour}$ battery firms only a few hours, so it is not equivalent
firm capacity. The decision turns on what the client is buying: energy
(wind-plus-storage wins on cost and carbon) or firm dispatchable
capacity for the rare multi-day low-wind event (gas, or a much larger
storage build, is needed). The defensible answer names the system
context, the carbon price, and the role the asset must play, rather
than declaring one option universally better.
\end{exercise}

\begin{exercise}
A regulator must decide whether to mandate a $20\,\%$ ancillary-
service reserve margin for wind and solar generation, paid for via a
tariff on the renewable generators, or to charge a market premium
based on actual grid services delivered. Identify the trade-off and
the experience of the European and Texas markets in implementing
each approach.
\paragraph{Solution sketch.}
A fixed mandated reserve is simple and predictable but blunt: it
charges every renewable generator the same, regardless of how much
variability it actually imposes, and it can over- or under-provision as
the fleet changes. A market for the actual services (frequency
response, ramping, firm capacity) prices the need where it arises and
rewards the resources that best supply it (batteries, demand response,
flexible plant), but it is more complex to design and can be gamed or
produce volatile prices. The European experience leans toward
service-based ancillary markets coordinated across interconnected
systems; the Texas (ERCOT) energy-only market historically relied on
scarcity pricing rather than a mandated capacity reserve, and the 2021
event drove changes toward stronger reserve and weatherisation
requirements. The judgement is that price-based mechanisms allocate
efficiently when markets are well designed and liquid, while mandates
provide a floor of reliability when they are not; most systems blend
the two.
\end{exercise}

\begin{exercise}
A municipal utility must decide whether to invest in district
heating (a network of insulated pipes carrying hot water from a
central plant to consumers) or in distributed heat pumps for
residential heating. Identify the engineering, economic, and
political trade-offs.
\paragraph{Solution sketch.}
District heating concentrates the heat source, so it can use waste heat
from industry, data centres, or combined heat-and-power, and it suits
dense urban areas where the pipe length per customer is short; the
capital is large, lumpy, and front-loaded (digging up streets), and it
locks in a single supplier. Distributed heat pumps are modular and
deployable house by house, exploit the COP to deliver more heat than
they consume electrically, and need no shared network, but they load
the electricity grid in winter and depend on each building's
insulation and retrofit. The economic trade-off is shared
high-utilisation infrastructure (district heat in dense areas) versus
incremental low-coordination investment (heat pumps in dispersed
areas). The political dimension is that district heating requires
collective commitment and right-of-way, while heat pumps devolve the
choice to individual owners. The defensible answer maps the choice to
density, available waste heat, grid capacity, and the local policy
appetite for coordinated infrastructure.
\end{exercise}

\begin{exercise}
A national government must decide whether to retain its existing
fleet of large nuclear plants beyond their original $40$-year
operating licence or to replace them with renewables-plus-storage.
Identify the trade-off in terms of LCOE, decommissioning cost,
schedule, and grid reliability.
\paragraph{Solution sketch.}
Life-extending an existing nuclear plant has a low marginal LCOE: the
capital is sunk, so the running cost (fuel, O\&M, the refurbishment to
extend the licence) buys reliable, high-capacity-factor, low-carbon
firm power, and it defers the large decommissioning liability. New
renewables-plus-storage has a lower energy LCOE but does not deliver
equivalent firm capacity without a large (expensive) storage build, and
it takes time to permit and construct. On schedule, retention keeps
existing low-carbon capacity online now, while replacement risks a gap
filled by fossil generation if the build lags. On reliability, the
nuclear fleet provides inertia and firm output that a
renewables-plus-storage system must reproduce through other means
(synthetic inertia, firm backup, interconnection). The judgement
usually favours life extension of a safe, performing plant as the
cheaper near-term low-carbon option, with new renewables added in
parallel rather than as a like-for-like swap; the answer should weigh
the plant's safety margin, the decommissioning timing, and the credible
build rate of the alternative.
\end{exercise}
